\documentclass[11pt,reqno]{amsart}

% ================================================================
% Packages
% ================================================================
\usepackage{amsmath,amssymb,amsthm}
\usepackage{mathrsfs}
\usepackage{enumerate}
\usepackage[shortlabels]{enumitem}
\usepackage{booktabs}
\usepackage{array}
\usepackage{microtype}
\usepackage[dvipsnames]{xcolor}
\usepackage[colorlinks=true,linkcolor=blue!60!black,
 citecolor=green!40!black,urlcolor=blue!60!black]{hyperref}

% ================================================================
% Page geometry
% ================================================================
\usepackage[margin=1.15in]{geometry}

% ================================================================
% Theorem environments
% ================================================================
\newtheorem{theorem}{Theorem}[section]
\newtheorem{proposition}[theorem]{Proposition}
\newtheorem{lemma}[theorem]{Lemma}
\newtheorem{corollary}[theorem]{Corollary}
\newtheorem{conjecture}[theorem]{Conjecture}
\theoremstyle{definition}
\newtheorem{definition}[theorem]{Definition}
\newtheorem{construction}[theorem]{Construction}
\newtheorem{example}[theorem]{Example}
\theoremstyle{remark}
\newtheorem{remark}[theorem]{Remark}

% ================================================================
% Macros
% ================================================================
\newcommand{\cA}{\mathcal{A}}
\newcommand{\cH}{\mathcal{H}}
\newcommand{\cM}{\mathcal{M}}
\newcommand{\cW}{\mathcal{W}}
\newcommand{\barB}{\bar{B}}
\newcommand{\Defcyc}{\mathrm{Def}_{\mathrm{cyc}}}
\newcommand{\Ran}{\mathrm{Ran}}
\newcommand{\Mbar}{\overline{\mathcal{M}}}
\newcommand{\MC}{\mathrm{MC}}
\newcommand{\Sh}{\mathrm{Sh}}
\newcommand{\FM}{\mathrm{FM}}
\newcommand{\fg}{\mathfrak{g}}
\newcommand{\fh}{\mathfrak{h}}
\newcommand{\gmod}{\mathfrak{g}^{\mathrm{mod}}_\cA}
\newcommand{\OPE}{\mathrm{OPE}}
\newcommand{\Vir}{\mathrm{Vir}}
\newcommand{\ChirHoch}{\mathrm{ChirHoch}}

\DeclareMathOperator{\gr}{gr}
\DeclareMathOperator{\rk}{rk}
\DeclareMathOperator{\depth}{depth}
\DeclareMathOperator{\Res}{Res}
\DeclareMathOperator{\Aut}{Aut}

\numberwithin{equation}{section}

% ================================================================
\begin{document}

\title[The Shadow Obstruction Tower]
{The Shadow Obstruction Tower}

\author{Raeez Lorgat}
\address{Perimeter Institute for Theoretical Physics,
 Waterloo, ON N2L 2Y5, Canada}
\email{rlorgat@perimeterinstitute.ca}

\date{}

\begin{abstract}
The chiral bar construction of a vertex algebra~$\cA$ on an
algebraic curve produces a Maurer--Cartan element
$\Theta_\cA$ in the modular convolution algebra, encoding
the full genus expansion. The degree-$r$ projections
of~$\Theta_\cA$ form the \emph{shadow obstruction tower},
an infinite sequence of invariants $S_2, S_3, S_4, \ldots$
determined by the operator product expansion.

We prove five results.
First, the weighted generating function
$H(t) = \sum_{r \ge 2} r\, S_r\, t^r$ satisfies
$H(t)^2 = t^4\, Q_L(t)$, where the \emph{shadow metric}
$Q_L(t) = 4\kappa^2 + 12\kappa\alpha\, t
+ (9\alpha^2 + 16\kappa S_4)\, t^2$
is a quadratic polynomial: the tower is algebraic
of degree~$2$, and every $S_r$ for $r \ge 5$ is
determined by three seeds $(\kappa, \alpha, S_4)$.
Second, the critical discriminant $\Delta = 8\kappa S_4$
partitions all chirally Koszul algebras into four depth
classes: Gaussian ($d_{\mathrm{alg}} = 0$), Lie
($d_{\mathrm{alg}} = 1$), contact ($d_{\mathrm{alg}} = 2$),
and mixed ($d_{\mathrm{alg}} = \infty$).
Third, the algebraic depth satisfies
$d_{\mathrm{alg}} \in \{0, 1, 2, \infty\}$; no
value~$3$ or any finite value $\ge 3$ is realized.
Fourth, a chirally Koszul vertex algebra is Swiss-cheese
formal if and only if it belongs to class~$\mathbf{G}$.
Fifth, the shadow obstruction tower is the
$\Sigma_n$-coinvariant image of the ordered $R$-matrix
tower in the $E_1$ convolution algebra.

We compute explicit shadow towers for all standard families
(Heisenberg, affine Kac--Moody, $\beta\gamma$, Virasoro,
$\cW_N$), extract the quartic contact invariant
$Q^{\mathrm{contact}}_{\Vir} = 10/[c(5c+22)]$,
determine the shadow connection and its Koszul
monodromy~$-1$, and compute the genus-$2$ cross-channel
correction $\delta F_2(\cW_3) = (c + 204)/(16c)$
by graph-by-graph stable graph summation over the seven
genus-$2$ stable graphs.
\end{abstract}

\subjclass[2020]{17B69, 14H10, 18M70, 81T40}

\keywords{Chiral algebras, bar construction, Koszul duality,
moduli of curves, vertex algebras, shadow tower,
W-algebras, depth classification}

\maketitle

\tableofcontents

% ================================================================
% 1. INTRODUCTION
% ================================================================

\section{Introduction}\label{sec:intro}

\subsection{The problem}\label{subsec:intro-problem}

Classical Koszul duality is a phenomenon of graded algebras
over a field. A quadratic algebra~$A$ determines a dual
coalgebra~$A^!$; the bar construction $B(A) = T^c(s^{-1}\bar{A})$
mediates between them; and the Koszul property (bar
cohomology concentrated in bar degree~$1$) is equivalent to
the existence of a minimal resolution. The theory is the
subject of Priddy~\cite{Priddy70},
Beilinson--Ginzburg--Soergel~\cite{BGS96}, and
Loday--Vallette~\cite{LV12}.

On an algebraic curve~$X$, the generators become
sections of a $\mathcal{D}_X$-module, the quadratic
relations become operator product expansions with
meromorphic singularities, and the bar differential
becomes an integral transform whose kernel is the
logarithmic $1$-form
$\eta_{ij} = d\log(z_i - z_j)$ on the Fulton--MacPherson
compactification $\FM_n(X)$. The resulting object, the
\emph{chiral bar complex}
$\barB^{\mathrm{ch}}(\cA) = T^c(s^{-1}\bar{\cA})$,
is a factorization coalgebra on the Ran space
$\Ran(X)$~\cite{BD04, FG12, CG17}.

At genus~$0$, the classical theory embeds via
restriction to the formal disk. The embedding is not
an equivalence: the configuration-space geometry of
$\mathrm{Conf}_n(\mathbb{A}^1)$, the Arnol'd algebra of
logarithmic forms, the Fulton--MacPherson
compactifications, and the $E_1/E_\infty$ bar
distinction are already present on~$\mathbb{A}^1$
but absent over a bare point; and the passage from the
formal disk to a point requires specifying a retraction
and its homotopy transfer data. At genus $g \ge 1$,
the topology of the curve forces
curvature: the fiberwise bar differential satisfies
$d_{\mathrm{fib}}^2 = \kappa(\cA) \cdot \omega_g$, where
$\omega_g = c_1(\mathbb{E})$ is the Hodge class on
$\Mbar_g$ and $\kappa(\cA) \in \mathbb{Q}(c)$ is an
intrinsic invariant called the \emph{modular
characteristic}. This curvature has no analogue in the
classical setting; its origin is the nontriviality of the
Hodge bundle over~$\Mbar_g$.

The present paper addresses a single question: \emph{what
invariants does the chiral bar construction produce beyond
the modular characteristic, and what constrains them?}

\subsection{The answer}\label{subsec:intro-answer}

Three invariants suffice. A \emph{primary line} is a
one-parameter subfamily of the modular deformation space
obtained by varying a single coupling constant while holding
all others fixed: for Heisenberg, the level~$k$; for
Virasoro, the central charge~$c$; for affine Kac--Moody,
the level~$k$ at fixed Lie type~$\fg$. Along each primary
line, the chiral bar construction produces an infinite
sequence of \emph{shadow invariants}
$S_2, S_3, S_4, \ldots$, and the entire sequence is
algebraic of degree~$2$: a single quadratic polynomial
\[
 Q_L(t)
 \;=\;
 4\kappa^2 + 12\kappa\alpha\,t
 + (9\alpha^2 + 16\kappa S_4)\,t^2,
\]
the \emph{shadow metric}, determines all $S_r$ for $r \ge 5$
through the relation $H(t)^2 = t^4\, Q_L(t)$, where
$H(t) = \sum_{r \ge 2} r\, S_r\, t^r$ is the weighted
shadow generating function.

The critical discriminant $\Delta := 8\kappa S_4$ controls
the depth of the tower: when $\Delta = 0$, the generating
function is polynomial and the tower terminates; when
$\Delta \ne 0$, the generating function has irrational
Taylor coefficients and the tower extends to infinity.
The transition is sharp: no intermediate depth is
possible.

\subsection{Main results}\label{subsec:intro-results}

\begin{theorem}[Riccati algebraicity;
Theorem~\ref{thm:riccati}]\label{thm:intro-riccati}
Let $\cA$ be a chirally Koszul vertex algebra, $L$ a
primary line \textup{(}e.g.\ $\cH_k$ parametrised by level~$k$,
or $\Vir_c$ parametrised by central charge~$c$\textup{)} with
shadow data $S_r$, and $\kappa = S_2 \ne 0$. Then
\begin{equation}\label{eq:intro-algebraicity}
 H(t)^2 \;=\; t^4\,Q_L(t).
\end{equation}
Every $S_r$ with $r \ge 5$ is determined by
$(\kappa, \alpha, S_4)$ via
$S_r = \frac{1}{r}[t^{r-2}]\sqrt{Q_L(t)}$.
\end{theorem}

\begin{theorem}[Four-class partition;
Theorem~\ref{thm:dichotomy}]\label{thm:intro-classification}
On any primary line~$L$ with $\kappa \ne 0$, the shadow
depth $r_{\max} := \sup\{r : S_r \ne 0\}$ belongs to
$\{2, 3, \infty\}$:
\begin{enumerate}[(i)]
\item \textbf{Class $\mathbf{G}$} (Gaussian):
 $\Delta = 0$, $\alpha = 0$; $r_{\max} = 2$.
\item \textbf{Class $\mathbf{L}$} (Lie):
 $\Delta = 0$, $\alpha \ne 0$; $r_{\max} = 3$.
\item \textbf{Class $\mathbf{M}$} (mixed):
 $\Delta \ne 0$; $r_{\max} = \infty$.
\end{enumerate}
A fourth class, \textbf{class~$\mathbf{C}$} (contact,
$r_{\max} = 4$), arises when the single-line tower terminates
but a quartic contact invariant from a two-field interaction
survives. The archetype is the $\beta\gamma$ system (see
Section~\ref{subsec:class-C} for the explicit computation).
\end{theorem}

\begin{theorem}[Depth gap;
Proposition~\ref{prop:depth-gap}]\label{thm:intro-depth-gap}
The algebraic depth satisfies
$d_{\mathrm{alg}}(\cA) \in \{0, 1, 2, \infty\}$.
No value~$3$ or any finite value~$\ge 3$ is realized on the
standard landscape. The four values correspond bijectively
to the classes $\mathbf{G}$, $\mathbf{L}$, $\mathbf{C}$,
$\mathbf{M}$.
\end{theorem}

\begin{theorem}[SC-formality characterizes class~$\mathbf{G}$;
Proposition~\ref{prop:sc-formal}]\label{thm:intro-sc-formal}
A chirally Koszul vertex algebra~$\cA$ in the standard
landscape is Swiss-cheese formal
($m_k^{\mathrm{SC}} = 0$ for all $k \ge 3$) if and only
if~$\cA$ belongs to class~$\mathbf{G}$.
\end{theorem}

\begin{theorem}[$E_1$ framing;
Proposition~\ref{prop:e1-framing}]\label{thm:intro-e1}
The shadow obstruction tower
$\{S_r\}_{r \ge 2}$ is the $\Sigma_n$-coinvariant image
of the ordered $R$-matrix tower in the $E_1$ convolution
algebra $\mathfrak{g}^{E_1}_\cA$: at degree~$2$,
$\kappa = \mathrm{av}(r(z))$ for abelian algebras, and
$\kappa = \mathrm{av}(r(z)) + \dim(\fg)/2$ for non-abelian
affine Kac--Moody (where the Sugawara shift contributes).
\end{theorem}

The shadow obstruction tower also carries a logarithmic
connection $\nabla^{\mathrm{sh}} = d - Q_L'/(2Q_L)\,dt$ with
regular singularities at the zeros of~$Q_L$; the monodromy
around each branch point is~$-1$ (the Koszul sign).

\subsection{Discussion of the
results}\label{subsec:intro-discussion}

Each of the five main results captures a different facet of
the shadow obstruction tower.

\emph{The algebraicity theorem}
(Theorem~\ref{thm:intro-riccati}) is the structural core:
it says that the MC recursion, which is an infinite system
of quadratic equations (one at each degree), is equivalent
to a single quadratic polynomial relation. The proof
converts the recursion into a convolution identity by the
substitution $a_n = (n+2)S_{n+2}$ and shows that the
convolution $F(t)^2 = Q_L(t)$ holds at all orders. The
key insight is that the MC recursion at degree~$r$ is
precisely the statement that the degree-$(r-2)$ convolution
vanishes: the recursion \emph{is} the algebraicity.

\emph{The four-class partition}
(Theorem~\ref{thm:intro-classification}) is an immediate
consequence: the factorization of the shadow metric $Q_L$
controls the depth. The classification is exhaustive for
the standard landscape and sharp: no intermediate depths
are possible (Theorem~\ref{thm:intro-depth-gap}).

\emph{The depth gap theorem}
(Theorem~\ref{thm:intro-depth-gap}) has a clean algebraic
explanation: on a single primary line, the shadow generating
function is a degree-$2$ algebraic function, which is either
rational (tower finite) or irrational (tower infinite).
The contact class~$\mathbf{C}$ arises by stratum
separation: one additional degree from a transverse
direction. No further finite extension is possible.

\emph{The SC-formality theorem}
(Theorem~\ref{thm:intro-sc-formal}) connects the shadow
tower to the operadic structure: the Swiss-cheese operations
$m_k^{\mathrm{SC}}$ vanish for $k \ge 3$ if and only if
the shadow tower terminates at degree~$2$. The proof uses
the tree-shadow correspondence: the degree-$r$ shadow is the
closed-colour projection of the degree-$r$ mixed SC tree.

\emph{The $E_1$ framing}
(Theorem~\ref{thm:intro-e1}) identifies the shadow tower as
a projection of a richer structure: the ordered $R$-matrix
tower in the $E_1$ convolution algebra. The averaging map
$\mathrm{av}$ is lossy: it discards the spectral parameter,
the $R$-matrix, and the Yangian. The shadow tower is the
$\Sigma_n$-coinvariant residue.

For the $\cW_3$ algebra at genus~$2$, the multi-channel
stable graph sum produces a cross-channel correction
$\delta F_2(\cW_3) = (c + 204)/(16c)$, resolving the
multi-generator universality problem negatively: the scalar
formula $F_g = \kappa \cdot \lambda_g^{\mathrm{FP}}$ fails
at genus~$\ge 2$ for multi-weight algebras (UNIFORM-WEIGHT
tag required; the correct decomposition is
$F_g = \kappa \cdot \lambda_g^{\mathrm{FP}} +
\delta F_g^{\mathrm{cross}}$). The computation is
graph-by-graph: four of the seven genus-$2$ stable graphs
contribute nonzero mixed-channel amplitudes, with the
tadpole contribution $1/16$ being the unique
$c$-independent term.

\subsection{The E$_1$ perspective}\label{subsec:intro-e1}

The shadow obstruction tower is a projection.
The primitive object is the ordered bar complex
$B^{\mathrm{ord}}(\cA) = T^c(s^{-1}\bar{\cA})$ with its
deconcatenation coproduct and spectral $R$-matrix
$r(z) \in \mathfrak{g}^{E_1}_\cA$. The shadow tower
$\{S_r\}$ and the algebraicity relation
$H(t)^2 = t^4\, Q_L(t)$ are statements about the
$\Sigma_n$-coinvariant image: the degree-$2$ scalar
$\kappa$, given by $\mathrm{av}(r(z))$ in abelian and scalar
families and by $\mathrm{av}(r(z))+\dim(\fg)/2$ for
non-abelian affine Kac--Moody, and its higher-degree
extensions under the lossy averaging map
$\mathrm{av}\colon \mathfrak{g}^{E_1}_\cA \to
\mathfrak{g}^{\mathrm{mod}}_\cA$.

The full $E_1$ tower retains the spectral parameter~$z$,
the $R$-matrix, and the Yangian structure; the shadow
tower discards the spectral parameter by averaging.
Everything in this paper is the $\Sigma_n$-coinvariant
projection of an $E_1$ object. The ordered bar is the
primitive; the shadow tower is the shadow.

\subsection{Context}\label{subsec:intro-context}

The algebraicity relation $H(t)^2 = t^4\, Q_L(t)$ is
structurally parallel to the loop equations of matrix
models~\cite{ADKMV, DijkgraafVafa} and to the
Eynard--Orantin topological
recursion~\cite{EO}: in those
frameworks the spectral curve is an input; here the spectral
curve $\Sigma = \{H^2 = t^4\, Q_L\}$ is an \emph{output}
of the algebraicity theorem, determined by the
vertex algebra.

The shadow connection has Stokes data controlled by the
critical discriminant: monodromy is trivial for classes
$\mathbf{G}$ and~$\mathbf{L}$ and nontrivial for
class~$\mathbf{M}$, in parallel with the wall-crossing
formalism of
Kontsevich--Soibelman~\cite{KontsevichSoibelman}.

Classical Koszul duality over a field embeds into the
genus-$0$, degree-$2$, $\Delta = 0$ stratum (class
$\mathbf{G}$, where the tower terminates immediately)
via restriction to the formal disk.
Everything beyond class~$\mathbf{G}$ is new.

\subsection{Notation and conventions}\label{subsec:notation}

Throughout, $X$ denotes a smooth projective algebraic curve
over a field of characteristic zero. All differentials have
cohomological degree~$+1$. The bar construction uses
desuspension ($s^{-1}$, degree~$-1$).

A \emph{vertex algebra} is a vertex algebra in the sense
of~\cite{FrenkelBenZvi, Kac-vertex}, equipped with a
conformal vector, a conformal grading
$\cA = \bigoplus_{h \ge 0} \cA_h$ with
$\dim \cA_h < \infty$, and the OPE
$a(z)b(w) \sim \sum_{n \ge 0} (a_{(n)}b)(w) \cdot
(z-w)^{-n-1}$.

The \emph{modular characteristic}
$\kappa(\cA) \in \mathbb{Q}(c)$ is the genus-$1$ curvature
scalar: $F_1(\cA) = \kappa(\cA)/24$. The Hodge class
$\lambda_g := c_1(\mathbb{E}_g) \in H^2(\Mbar_g,
\mathbb{Q})$ is the first Chern class of the Hodge bundle.
The Faber--Pandharipande coefficients are
$\lambda_g^{\mathrm{FP}} =
\int_{\Mbar_{g,1}} \psi_1^{2g-2}\, \lambda_g$, with
$\lambda_1^{\mathrm{FP}} = 1/24$ and
$\lambda_2^{\mathrm{FP}} = 7/5760$.

\subsection{Guide to the reader}\label{subsec:guide}

The paper is organized around the shadow obstruction tower
as the central object, constructed from first principles and
then analyzed through five independent lenses.

Section~\ref{sec:bar} constructs the chiral bar complex
and the Maurer--Cartan element $\Theta_\cA$, beginning
with the Heisenberg algebra as the motivating example.
Section~\ref{sec:tower} defines the shadow obstruction tower
as the sequence of degree projections of~$\Theta_\cA$ and
establishes the master equation.
Section~\ref{sec:algebraicity} proves the Riccati
algebraicity theorem: the conversion of the MC recursion
into a convolution identity, the degree-by-degree
verification, and the closed-form extraction of shadow
coefficients.
Section~\ref{sec:class} derives the four-class partition
from the factorization of the shadow metric.
Section~\ref{sec:depth-gap} proves the depth gap theorem,
establishing that $d_{\mathrm{alg}} = 3$ is impossible.
Section~\ref{sec:sc-formal} proves that SC-formality
characterizes class~$\mathbf{G}$.
Section~\ref{sec:e1} develops the $E_1$ framing: the
shadow tower as the averaged projection of the ordered
$R$-matrix tower.
Sections~\ref{sec:connection} and~\ref{sec:landscape}
extract the shadow connection and compute explicit towers
for all standard families.
Section~\ref{sec:genus2} carries out the $\cW_3$
cross-channel computation at genus~$2$.

The logical dependencies are:
$\text{\S\ref{sec:bar}} \to \text{\S\ref{sec:tower}} \to
\text{\S\ref{sec:algebraicity}} \to
\text{\S\ref{sec:class}} \to
\text{\S\ref{sec:depth-gap}}$;
the sections on SC-formality, $E_1$ framing, the shadow
connection, and the genus-$2$ computation are logically
independent of each other and may be read in any order
after Section~\ref{sec:class}.

\subsection{Relation to companion
papers}\label{subsec:companion}

The $c$-independent identity $S_3(\Vir_c) = 2$ is proved in
\cite{LorgatVirR}. The seven-face programme connecting the
shadow tower to Gaudin models, Sklyanin brackets, and
commuting Hamiltonians is developed
in~\cite{LorgatSevenFaces}. The present paper is
self-contained: all results are proved from the chiral bar
construction without reference to the companion papers.


% ================================================================
% 2. THE CHIRAL BAR COMPLEX
% ================================================================

\section{The chiral bar complex and Maurer--Cartan element}
\label{sec:bar}

\subsection{The chiral bar complex}\label{subsec:bar-complex}

Let $\cA$ be a vertex algebra on a smooth projective
curve~$X$. The Ran space
$\Ran(X) = \coprod_{n \ge 1} X^n/\Sigma_n$ parameterizes
nonempty finite subsets of~$X$. For each $n \ge 1$, the
Fulton--MacPherson compactification $\FM_n(X)$ resolves
the diagonals in~$X^n$ by iterated blowups along all
partial diagonals, producing a smooth manifold-with-corners
whose boundary strata are indexed by trees of
collisions~\cite{GK}.

\begin{definition}[Chiral bar complex]\label{def:bar-complex}
The \emph{chiral bar complex} $B(\cA)$ is the factorization
coalgebra on $\Ran(X)$ with:
\begin{enumerate}[(i)]
\item \emph{Underlying space}:
 $B(\cA)_n := (s^{-1}\bar{\cA})^{\otimes n}$, where
 $\bar{\cA} = \ker(\varepsilon)$ is the augmentation ideal
 and $s^{-1}$ denotes desuspension (degree~$-1$).

\item \emph{Bar differential}: the coderivation $d_B$ whose
 components are the collision residues
 \[
  d_B|_{B(\cA)_n \to B(\cA)_{n-1}}
  \;=\;
  \sum_{i < j}
  \Res_{z_i \to z_j}
  a_i(z_i)\, a_j(z_j)\, d\!\log(z_i - z_j),
 \]
 extracting the OPE residue along the logarithmic
 differential $d\!\log(z_i - z_j)$.

\item \emph{Coproduct}: the cofactorization structure
 induced by partition of points into clusters.
\end{enumerate}
The identity $d_B^2 = 0$ encodes OPE associativity.
\end{definition}

\begin{remark}[The $d\!\log$ absorption]\label{rem:dlog}
The bar differential extracts residues along
$d\!\log(z - w)$, not along $dz/(z-w)$. The logarithmic
differential absorbs one power of $(z-w)$: if the OPE has
poles at $(z-w)^{-n-1}$, the collision residue has poles
at $(z-w)^{-n}$. For Virasoro, the OPE has a quartic pole;
the bar $r$-matrix has a cubic pole
$r^{\Vir}(z) = (c/2)/z^3 + 2T/z$.
\end{remark}

\begin{remark}[The ordered bar as
primitive]\label{rem:ordered-bar}
The bar complex of Definition~\ref{def:bar-complex} is the
symmetric (factorization) bar $B^\Sigma(\cA)$, the
$\Sigma_n$-coinvariant quotient of the ordered bar
$B^{\mathrm{ord}}(\cA) = T^c(s^{-1}\bar{\cA})$. The
ordered bar carries a deconcatenation coproduct
(coassociative, not cocommutative) and supports the spectral
$R$-matrix $r(z) \in \mathfrak{g}^{E_1}_\cA$. The shadow
obstruction tower $\{S_r\}$ and the algebraicity relation
$H(t)^2 = t^4\, Q_L(t)$ are statements about the
$\Sigma_n$-coinvariant image: the degree-$2$ scalar
$\kappa$, given by $\mathrm{av}(r(z))$ in abelian and scalar
families and by $\mathrm{av}(r(z))+\dim(\fg)/2$ for
non-abelian affine Kac--Moody, and its higher-degree
extensions. The ordered bar is the primitive object; the
shadow tower is its averaged projection.
\end{remark}

\subsection{Chirally Koszul algebras}\label{subsec:koszul}

\begin{definition}[Chirally
Koszul]\label{def:chirally-koszul}
A vertex algebra~$\cA$ is \emph{chirally Koszul} if the bar
cohomology $H^\bullet(B(\cA), d_B)$ is concentrated in bar
degree~$1$: for every $n \ge 2$, the map
$H^\bullet(B(\cA)_n) \to H^\bullet(B(\cA)_1)^{\otimes n}$
induced by the coproduct is a quasi-isomorphism.
\end{definition}

The standard families (Heisenberg, affine Kac--Moody at
non-critical level $k \ne -h^\vee$, $\beta\gamma$,
Virasoro, and $\cW_N$) are all chirally Koszul.

\subsection{The modular convolution algebra}
\label{subsec:convolution}

\begin{definition}[Modular cyclic deformation
complex]\label{def:mod-cyc}
The \emph{modular cyclic deformation complex} of~$\cA$ is
\begin{equation}\label{eq:mod-cyc}
 \Defcyc^{\mathrm{mod}}(\cA)
 \;:=\;
 \prod_{\substack{g,n \\ 2g-2+n>0}}
 \operatorname{CoDer}^{\mathrm{cyc}}\!\bigl(
 B^{(g,n)}(\cA)\bigr)[1],
\end{equation}
where the product ranges over all stable pairs.
The differential is $d = d_{\mathrm{int}} +
d_{\mathrm{sew}}$, with $d_{\mathrm{int}}$ the internal
cyclic coderivation differential and $d_{\mathrm{sew}}$ the
clutching operation from boundary maps on~$\Mbar_{g,n}$.
The Lie bracket is \emph{graph composition}: for
coderivations $\xi$ and~$\eta$,
\begin{equation}\label{eq:graph-bracket}
 [\xi,\eta]
 \;:=\;
 \sum_\Gamma \xi \circ_\Gamma \eta,
\end{equation}
summed over stable graphs connecting an output of~$\xi$
to an input of~$\eta$.
\end{definition}

\subsection{The Maurer--Cartan
element}\label{subsec:mc-element}

\begin{proposition}[Bar-intrinsic MC
element]\label{prop:bar-mc}
Let $D_\cA$ denote the full bar differential on $B(\cA)$
and $d_0$ its genus-$0$ linear part. Define
\begin{equation}\label{eq:theta-def}
 \Theta_\cA \;:=\; D_\cA - d_0
 \;\in\; \Defcyc^{\mathrm{mod}}(\cA).
\end{equation}
Then $\Theta_\cA$ satisfies the Maurer--Cartan equation
\begin{equation}\label{eq:mc-equation}
 d_0\,\Theta_\cA + \tfrac{1}{2}[\Theta_\cA, \Theta_\cA]
 = 0.
\end{equation}
\end{proposition}

\begin{proof}
The identity $D_\cA^2 = 0$ is a formal consequence of the
boundary-of-boundary relation $\partial^2 = 0$ on
$\Mbar_{g,n}$. Substituting $D_\cA = d_0 + \Theta_\cA$
into $D_\cA^2 = 0$ and expanding:
\[
 0 \;=\; d_0^2 + d_0\,\Theta_\cA + \Theta_\cA \circ d_0
 + \Theta_\cA \circ \Theta_\cA.
\]
Since $d_0^2 = 0$ (OPE associativity) and
the remaining terms identify with
$d_0\,\Theta_\cA + \frac{1}{2}[\Theta_\cA, \Theta_\cA]$
in the cyclic coderivation Lie algebra, the result follows.
\end{proof}


% ================================================================
% WAVE-14 RECONCEPTION ANCHOR (2026-04-17)
% ================================================================
%
% This standalone is the canonical Platonic-ideal presentation of the
% shadow obstruction tower. Six wave-14 anchors fix the scope:
%
% (A) Pole-doubling for all k.  Sigma_Vir(z) := sum A_r z^r has the
%     closed form
%         Sigma_Vir(z) = 4z^3/9 - z^2/9 + z/27 - log(1+6z)/162,
%     and every derivative admits the decomposition
%         Sigma_Vir^{(k)}(z) = [P_k(z) + Q_k(z) log(1+6z)] / (1+6z)^{2k}
%     with (P_k, Q_k) polynomial; the radius of convergence is 1/6
%     (C_A = 6 universality on non-logarithmic class M carrying a
%     Virasoro subalgebra; thm:pole-doubling-all-k and
%     thm:universal-class-M-C-is-6 in
%     chapters/theory/shadow_tower_higher_coefficients.tex).
%
% (B) Seed decomposition 6 = r_0 . S_{r_0} = 3 . 2.
%     The shadow-exponential base 6 factorises through
%         r_0 = 3 (bar-depth seed on the primary line),
%         S_{r_0}(Vir) = 2 (Riccati seed coefficient).
%     Recorded as cache entries in appendices/first_principles_cache.md.
%
% (C) Sub-subleading coefficient beta_N = (N+1)(N+2)/2.
%     beta_A (Vol II tempering rate) coincides with C_A at leading 1/c;
%     beta_N = (N+1)(N+2)/2 in Vol II absorbs higher-channel sub-leading
%     corrections at finite c and is consistent with C_{W_N} = 6
%     asymptotically. At N = 2 the closed form recovers the Virasoro
%     value; at N = 4 the programme value is 15 (the N^2 - N + 4 = 16
%     closed form is a separate parameter question, not structural).
%
% (D) Riccati seed (r_0, S_{r_0}(Vir)) = (3, 2). The recursion
%     H(t)^2 = t^4 Q_L(t) on Vir_c is initialised by the c-independent
%     S_3(Vir) = 2 (Theorem 6.2 of this document), which together with
%     kappa = c/2 and S_4 = 10/[c(5c+22)] determines the entire tail.
%
% (E) Retractions (B87, B92).  The "1/e Stirling-cancellation
%     obstruction" claim on the tempered stratum is RETRACTED; the
%     correct limsup is 0 universally at generic c.  The Kummer-irregular
%     prime labels on {1423, 3067, 23, 43, 419} are RETRACTED at primary
%     source (Bernoulli-witness search finds none); these primes still
%     appear in the S_r(Vir_c) numerators as RICCATI-ARITHMETIC
%     characteristic primes together with {61, 193, 2111, 16657,
%     3988097}.  The leading Kummer-irregular primes are {691, 3617}.
%     The Tier-3 emergence set becomes {37, 691, 811}; 3067 is removed.
%
% (F) Sharpness.  Arithmetic duality thm:z-g-s-r-arithmetic-duality
%     (chapters/theory/z_g_kummer_bernoulli_platonic.tex and
%     chapters/theory/shadow_tower_higher_coefficients.tex) is sharp at
%     {691, 3617}: these primes appear on Z_g via B_{2g-2} and are
%     absent from S_r(Vir_c) through r = 11
%     (thm:s-r-kummer-absent-through-r-11); the higher Kummer-irregular
%     prime 2111 appears in N_9 (via 29554 = 2 . 7 . 2111) but is not
%     witnessed on Z_g at g <= 9 (since 2111 does not divide num(B_{2m})
%     for 2m <= 16).
%
% ================================================================

% ================================================================
% 3. THE SHADOW OBSTRUCTION TOWER
% ================================================================

\section{The shadow obstruction tower}\label{sec:tower}

\subsection{Definition and master
equation}\label{subsec:tower-def}

\begin{definition}[Shadow obstruction
tower]\label{def:shadow-tower}
The \emph{shadow obstruction tower} of~$\cA$ is the sequence
of truncations $\Theta_\cA^{\le r}$ ($r \ge 2$) obtained by
retaining only the components of degree at most~$r$ in
$\Defcyc^{\mathrm{mod}}(\cA)$.
The \emph{degree-$r$ shadow} is the projection
\[
 \Sh_r(\cA)
 \;:=\;
 \pi_r(\Theta_\cA)
 \;\in\;
 \operatorname{CoDer}^{\mathrm{cyc}}(B^{(\bullet,r)}
 (\cA))[1].
\]
\end{definition}

\begin{proposition}[All-degree master
equation]\label{prop:master-eq}
For each $r \ge 2$, the degree-$r$ projection of the
MC equation gives
\begin{equation}\label{eq:master-eq}
 d_0(\Sh_r) + \sum_{\substack{j+k = r+2 \\ j,k \ge 2}}
 [\Sh_j, \Sh_k] \;=\; 0.
\end{equation}
The shadow $\Sh_r$ is determined recursively by
$\Sh_2, \ldots, \Sh_{r-1}$ and a cohomological obstruction.
\end{proposition}

\begin{proof}
The differential $d_0$ and the bracket $[-,-]$ respect the
degree filtration. Restricting the MC equation to degree~$r$:
the linear term is $d_0(\Sh_r)$, and the quadratic term at
degree~$r$ receives contributions from pairs $(j,k)$ with
$j + k = r + 2$ (the bracket consumes two inputs via the
internal edge, reducing total degree by~$2$).
\end{proof}

\subsection{Primary lines and shadow
data}\label{subsec:primary-lines}

\begin{definition}[Primary line and shadow
data]\label{def:primary-line}
A \emph{primary line} in $\Defcyc^{\mathrm{mod}}(\cA)$ is a
one-dimensional subspace $L = k \cdot \eta$ spanned by a
single primary field~$\eta$ of definite conformal weight.
On~$L$, the degree-$r$ shadow restricts to a scalar:
$\Sh_r|_L = S_r \cdot x^r$, where $x$ is the coordinate
on~$L$ and $S_r = S_r(\cA) \in \mathbb{Q}(c)$.

The first three scalars are:
\begin{itemize}
\item $S_2 =: \kappa$, the \emph{modular characteristic}:
 $F_1(\cA) = \kappa/24$.
\item $S_3 =: \alpha$, the \emph{cubic shadow}:
 the three-point collision residue on $\FM_3(X)$.
\item $S_4$, the \emph{quartic shadow}:
 the four-point collision residue, the first
 nonlinear OPE invariant.
\end{itemize}
\end{definition}

\subsection{The shadow metric}\label{subsec:shadow-metric}

\begin{definition}[Shadow metric and critical
discriminant]\label{def:shadow-metric}
Let $L$ be a primary line with $\kappa = S_2 \ne 0$ and
$\alpha = S_3$. Define the \emph{weighted initial data}
$a_0 := 2\kappa$, $a_1 := 3\alpha$, $a_2 := 4S_4$.
The \emph{shadow metric} on~$L$ is
\begin{equation}\label{eq:shadow-metric}
 Q_L(t)
 \;:=\;
 a_0^2 + 2a_0 a_1\,t + (a_1^2 + 2a_0 a_2)\,t^2
 \;=\;
 4\kappa^2 + 12\kappa\alpha\,t
 + (9\alpha^2 + 16\kappa S_4)\,t^2.
\end{equation}
The \emph{critical discriminant} is
\begin{equation}\label{eq:critical-disc}
 \Delta \;:=\; 8\kappa S_4.
\end{equation}
\end{definition}

\begin{lemma}[Gaussian decomposition]\label{lem:gaussian}
\begin{equation}\label{eq:gaussian-decomp}
 Q_L(t)
 \;=\;
 (2\kappa + 3\alpha\,t)^2
 \;+\;
 2\Delta\,t^2.
\end{equation}
The first term is the shadow metric of the truncated tower
with $S_r = 0$ for $r \ge 4$. The second term is the
quartic interaction correction.
\end{lemma}

\begin{proof}
$(2\kappa + 3\alpha t)^2 = 4\kappa^2 + 12\kappa\alpha\,t
+ 9\alpha^2\,t^2$.
Adding $2\Delta\,t^2 = 16\kappa S_4\,t^2$ recovers
$Q_L(t)$.
\end{proof}


\subsection{The obstruction
classes}\label{subsec:obstructions}

The MC equation at degree~$r$ takes the form
$d_0(\Sh_r) = -\frac{1}{2}\sum_{j+k=r+2} [\Sh_j, \Sh_k]$.
The right-hand side is a cocycle in the cyclic deformation
complex (it is $d_0$-closed by the Jacobi identity for
$[-,-]$ and the MC equation at lower degrees). The shadow
coefficient $S_r$ exists if and only if this cocycle is a
coboundary.

\begin{definition}[Obstruction
classes]\label{def:obstructions}
The \emph{degree-$r$ obstruction class} is
\[
 o_r(\cA)
 \;:=\;
 \Bigl[-\frac{1}{2}
 \sum_{\substack{j+k=r+2 \\ j,k \ge 2}}
 [\Sh_j, \Sh_k]\Bigr]
 \;\in\;
 H^2(\cA^{\mathrm{sh}}_{r,0}),
\]
the cohomology class of the quadratic obstruction in the
shadow algebra
$\cA^{\mathrm{sh}} := H_\bullet(
\Defcyc^{\mathrm{mod}}(\cA))$.
When $o_r = 0$, the shadow coefficient $S_r$ exists and is
determined (up to a coboundary, which is fixed by
normalization). When $o_r \ne 0$, the tower is
obstructed at degree~$r$.
\end{definition}

\begin{proposition}[Obstructions vanish for Koszul
algebras]\label{prop:obs-vanish}
If $\cA$ is chirally Koszul, then $o_r(\cA) = 0$ for all
$r \ge 2$: the tower is unobstructed, and the full MC
element $\Theta_\cA$ exists.
\end{proposition}

\begin{proof}
The Koszul property implies that the bar cohomology is
concentrated in bar degree~$1$, which forces
$H^2(\cA^{\mathrm{sh}}_{r,0}) = 0$ for the relevant
component of the shadow algebra. The vanishing of the
second cohomology group means that every cocycle is a
coboundary, so the obstruction class vanishes.

More precisely: the Koszulness of~$\cA$ implies that the
spectral sequence computing the shadow algebra cohomology
degenerates at $E_1$, and the degeneration forces the
higher obstructions to be trivial. This is the
vertex-algebraic analogue of the fact that a Koszul algebra
has a linear resolution: the bar differential has no higher
corrections.
\end{proof}

\begin{remark}[Beyond the Koszul
locus]\label{rem:beyond-koszul}
All standard families (Heisenberg, affine Kac--Moody at
non-critical level, $\beta\gamma$, Virasoro, $\cW_N$) are
chirally Koszul, and for these the shadow tower is
unobstructed. The obstruction theory becomes nontrivial
beyond the Koszul locus: for example, at the critical level
$k = -h^\vee$ of affine Kac--Moody, $\kappa = 0$ and the
propagator $P = 1/\kappa$ is undefined. The shadow tower
does not exist in the standard sense; the modular
convolution algebra requires a different filtration (the
critical filtration, in which the center of the completed
enveloping algebra replaces the curvature~$\kappa$).
Admissible-level quotients $L_k(\fg)$ (where null vectors
modify the bar complex) are another natural class where
obstructions may arise.
\end{remark}


\subsection{The shadow algebra
structure}\label{subsec:shadow-algebra}

The shadow algebra
$\cA^{\mathrm{sh}} = H_\bullet(
\Defcyc^{\mathrm{mod}}(\cA))$ carries a graded Lie
bracket $[-,-]^{\mathrm{sh}}$, inherited from the graph
composition bracket~\eqref{eq:graph-bracket}. The shadow
obstruction tower $\{S_r\}$ is a sequence of elements in
$\cA^{\mathrm{sh}}$ satisfying the projected MC equation.

\begin{proposition}[Structure of the shadow
algebra]\label{prop:shadow-algebra}
\leavevmode
\begin{enumerate}[(i)]
\item The shadow algebra $\cA^{\mathrm{sh}}$ is a
 pronilpotent graded Lie algebra with degree filtration
 $F^r \cA^{\mathrm{sh}} = \{$elements of degree
 $\ge r\}$.

\item On a primary line~$L$, the shadow algebra restricts
 to a one-dimensional Lie algebra at each degree. The
 bracket $[\Sh_j, \Sh_k]^{\mathrm{sh}}|_L$ is a scalar
 multiple of $\Sh_{j+k-2}|_L$, with the scalar determined
 by the propagator~$P = 1/\kappa$ and the symmetry
 factors.

\item The MC equation in $\cA^{\mathrm{sh}}$ is equivalent
 to the algebraic relation $H(t)^2 = t^4\, Q_L(t)$
 (Theorem~\ref{thm:riccati}).

\item The shadow algebra carries a $\mathbb{Z}$-grading
 (by degree) and a $\mathbb{Q}(c)$-rational structure
 (all structure constants are rational functions of the
 central charge).
\end{enumerate}
\end{proposition}

\begin{remark}[The shadow algebra is not a
ring]\label{rem:not-ring}
The shadow algebra $\cA^{\mathrm{sh}}$ has a graded Lie
bracket, not a ring structure. The MC element
$\Theta_\cA$ is an element of the Lie algebra, not a
generator of a ring. The algebraicity theorem
(Theorem~\ref{thm:riccati}) is a statement about the Lie
algebra structure: the MC equation, which is quadratic in
the Lie bracket, reduces to a quadratic algebraic equation
for the generating function.
\end{remark}


\subsection{The Heisenberg shadow tower: the base
computation}\label{subsec:heisenberg-base}

The Heisenberg algebra is the simplest nontrivial example.
Its shadow tower terminates immediately, providing the base
case from which all other computations depart.

\begin{example}[Heisenberg: degree-by-degree
construction]\label{ex:heisenberg-degrees}
For $\cH_k$ with $k \ne 0$:

\emph{Degree $2$.}\enspace
The Heisenberg OPE $J(z)J(w) \sim k/(z-w)^2$ is purely
quadratic. The collision residue along $d\!\log(z-w)$ extracts
the coefficient of $(z-w)^{-2}$ after absorbing one power:
$\Res_{z \to w} k/(z-w)^2 \cdot d\!\log(z-w)
= \Res_{z \to w} k/(z-w)^2 \cdot dz/(z-w)
= k \cdot \Res_{z \to w} (z-w)^{-3}\, dz$,
but the bar differential extracts the OPE residue in the
standard normalization, giving $S_2 = \kappa = k$.

\emph{Degree $3$.}\enspace
The three-point collision residue on $\FM_3(X)$ involves
the OPE $J(z_1)J(z_2)J(z_3)$. Since the Heisenberg OPE
has no simple-pole term ($J_{(0)}J = 0$), the cubic
collision vanishes: $S_3 = \alpha = 0$.

\emph{Degree $4$ and higher.}\enspace
All higher collision residues vanish for the same reason:
the Heisenberg OPE is purely quadratic, so no connected
correlator on $\FM_n(X)$ with $n \ge 3$ is nonzero.
The shadow metric is $Q_L(t) = 4k^2$, a perfect square,
and $H(t) = 2k\, t^2$.

\emph{Verification.}\enspace
$H(t)^2 = 4k^2 t^4 = t^4 \cdot 4k^2 = t^4 Q_L(t)$.
The algebraicity theorem is trivially satisfied.
\end{example}


% ================================================================
% 4. THE RICCATI ALGEBRAICITY THEOREM
% ================================================================

\section{The Riccati algebraicity theorem}
\label{sec:algebraicity}

\subsection{Statement}\label{subsec:riccati-statement}

\begin{theorem}[Riccati algebraicity]\label{thm:riccati}
Let $\cA$ be a chirally Koszul vertex algebra, $L$ a
primary line with shadow data $S_r$, curvature
$\kappa = S_2 \ne 0$, and cubic coefficient
$\alpha = S_3$. Define the shadow metric $Q_L(t)$ and
the weighted shadow generating function
$H(t) := \sum_{r \ge 2} r\, S_r\, t^r$. Then
\begin{equation}\label{eq:algebraicity}
 H(t)^2 \;=\; t^4\,Q_L(t).
\end{equation}
Equivalently, $H(t) = t^2\sqrt{Q_L(t)}$.
Every shadow coefficient $S_r$ with $r \ge 5$ is
determined by the three genus-$0$ invariants
$(\kappa, \alpha, S_4)$ via
$S_r = \frac{1}{r}\,[t^{r-2}]\sqrt{Q_L(t)}$.
\end{theorem}

\subsection{The single-line
recursion}\label{subsec:recursion}

\begin{lemma}[MC recursion on a primary
line]\label{lem:mc-recursion}
Let $P := 1/\kappa$. On the primary line~$L$, the
master equation reduces to
\begin{equation}\label{eq:recursion}
 S_r
 \;=\;
 -\frac{P}{2r}
 \sum_{\substack{j + k = r + 2 \\ 3 \le j \le k}}
 c_{jk}\, jk\, S_j\, S_k,
 \qquad
 c_{jk} =
 \begin{cases} 1 & j < k, \\ 1/2 & j = k, \end{cases}
\end{equation}
for all $r \ge 5$.
\end{lemma}

\begin{proof}
On the one-dimensional subspace~$L$, the bracket
$[\Sh_j, \Sh_k]$ restricts to a scalar pairing.
The cyclic coderivation bracket on~$L$ evaluates as
$[\Sh_j, \Sh_k]_L = c_{jk}\, jk\, P\, S_j S_k\, x^{j+k-2}$:
the factor $jk$ comes from the number of ways to select
one input from each coderivation, the factor
$P = 1/\kappa$ from the propagator on the internal edge,
and $c_{jk}$ corrects for the symmetry when $j = k$.
The linear term $d_0(\Sh_r)$ contributes
$2r\,\kappa\,S_r$. Dividing by $2r\,\kappa$ gives the
stated recursion.
\end{proof}


\subsection{Proof of algebraicity}\label{subsec:proof}

\begin{proof}[Proof of Theorem~\textup{\ref{thm:riccati}}]
Define $a_n := (n+2)\, S_{n+2}$ for $n \ge 0$, so that
$a_0 = 2\kappa$, $a_1 = 3\alpha$, $a_2 = 4S_4$, and
$H(t) = \sum_{n \ge 0} a_n\, t^{n+2}$.
Set $F(t) := H(t)/t^2 = \sum_{n \ge 0} a_n\, t^n$.
The claim is equivalent to $F(t)^2 = Q_L(t)$.

Expand $F(t)^2 = \sum_{m \ge 0}
\bigl(\sum_{i+j=m} a_i a_j\bigr)\, t^m$.
The first three coefficients are:
\begin{alignat*}{2}
 m &= 0: &\quad a_0^2 &= 4\kappa^2, \\
 m &= 1: &\quad 2a_0 a_1 &= 12\kappa\alpha, \\
 m &= 2: &\quad a_1^2 + 2a_0 a_2
 &= 9\alpha^2 + 16\kappa S_4.
\end{alignat*}
These are exactly the coefficients of $Q_L(t)$.

For $m \ge 3$, we show that the convolution
$\sum_{i+j=m} a_i a_j = 0$.
Setting $r = m + 2$:
\[
 \sum_{i+j=m} a_i a_j
 \;=\;
 \sum_{\substack{p+q = r+2 \\ p, q \ge 2}}
 pq\, S_p\, S_q,
\]
where $p = i+2$, $q = j+2$.
Splitting off the $p = 2$ and $q = 2$ terms:
\begin{equation}\label{eq:convolution-split}
 \sum_{i+j=m} a_i a_j
 \;=\;
 4r\,\kappa\, S_r
 \;+\;
 \sum_{\substack{p+q=r+2 \\ p,q \ge 3}}
 pq\, S_p\, S_q.
\end{equation}
The MC recursion (Lemma~\ref{lem:mc-recursion}) gives
\[
 S_r = -\frac{1}{2r\kappa}
 \sum_{\substack{p+q=r+2 \\ 3 \le p \le q}}
 c_{pq}\, pq\, S_p\, S_q.
\]
Symmetrizing:
$\sum_{\substack{p+q=r+2 \\ p,q \ge 3}} pq\, S_p\, S_q
= 2 \sum_{\substack{p+q=r+2 \\ 3 \le p \le q}}
c_{pq}\, pq\, S_p\, S_q
= -4r\,\kappa\, S_r$.
Substituting into~\eqref{eq:convolution-split}:
\[
 \sum_{i+j=m} a_i a_j
 \;=\;
 4r\,\kappa\, S_r + (-4r\,\kappa\, S_r)
 \;=\; 0.
\]
Hence $[t^m](F^2 - Q_L) = 0$ for all $m \ge 3$.
Combined with the verification at $m = 0, 1, 2$,
this proves $F(t)^2 = Q_L(t)$.

The closed form follows:
$r\, S_r = [t^{r-2}]F(t) = [t^{r-2}]\sqrt{Q_L(t)}$.
\end{proof}


\begin{remark}[Riccati ODE]\label{rem:riccati-ode}
Set $U(t) := \sum_{r \ge 3} S_r\, t^r$, the nonlinear part
of the shadow generating function. The algebraic
relation is equivalent to the first-order ODE
\begin{equation}\label{eq:riccati-ode}
 2t\, U'(t) \;+\; \frac{1}{2\kappa}\bigl(U'(t)\bigr)^2
 \;=\; 6\alpha\, t^3
 + \frac{9\alpha^2 + 2\Delta}{2\kappa}\, t^4,
\end{equation}
quadratic in~$U'$: the shadow obstruction tower satisfies a
Riccati-type equation of degree~$2$, and all terms at
$t^r$ for $r \ge 5$ cancel identically.
\end{remark}

\begin{remark}[The intrinsic quartic
principle]\label{rem:quartic-principle}
Three numbers determine an infinite sequence. The analogue is
the Weierstrass equation: two parameters $(g_2, g_3)$
determine all Laurent coefficients of~$\wp$, because $\wp$
satisfies an algebraic equation of degree~$2$
in~$\wp'$. Here three parameters
$(\kappa, \alpha, S_4)$ determine all $S_r$, because $H$
satisfies an algebraic equation of degree~$2$.
\end{remark}


\subsection{Closed-form extraction of shadow
coefficients}\label{subsec:closed-form}

The algebraicity theorem converts the recursive computation
of shadow coefficients into a single binomial expansion.

\begin{corollary}[Explicit extraction
formula]\label{cor:extraction}
For $r \ge 2$, the shadow coefficient is
\[
 S_r
 \;=\;
 \frac{1}{r}\,[t^{r-2}]\sqrt{Q_L(t)}
 \;=\;
 \frac{1}{r}\,[t^{r-2}]
 \sqrt{4\kappa^2 + 12\kappa\alpha\, t
 + (9\alpha^2 + 16\kappa S_4)\, t^2}.
\]
Expanding the square root via the binomial series in the
normalized variable
$u := (3\alpha\, t)/(2\kappa) + (9\alpha^2 + 2\Delta)\,
t^2/(4\kappa^2)$:
\begin{equation}\label{eq:binomial-extraction}
 \sqrt{Q_L(t)}
 \;=\;
 2\kappa\,\sum_{n=0}^{\infty}
 \binom{1/2}{n}\, u^n,
\end{equation}
where $\binom{1/2}{n} = \frac{(1/2)(1/2 - 1)
\cdots(1/2 - n + 1)}{n!}
= \frac{(-1)^{n-1}}{2^{2n-1}\,n}\binom{2n-2}{n-1}$
for $n \ge 1$.
\end{corollary}

\begin{example}[First five coefficients by
extraction]\label{ex:first-five}
Writing $a := 3\alpha/(2\kappa)$ and
$b := (9\alpha^2 + 2\Delta)/(4\kappa^2)$ for the
normalized coefficients of~$u$:
\begin{align}
 S_2 &\;=\; \kappa, \label{eq:S2-extract} \\
 S_3 &\;=\; \alpha, \label{eq:S3-extract} \\
 S_4 &\;=\; S_4 \qquad\text{(tautological at this
 degree)}, \label{eq:S4-extract} \\
 S_5 &\;=\;
 \frac{2\kappa}{5}\left(-\frac{a\,b}{4}
 + \frac{a^3}{16}\right)
 \;=\;
 -\frac{3\alpha\,\Delta}{10\,\kappa^2},
 \label{eq:S5-extract} \\
 S_6 &\;=\;
 \frac{2\kappa}{6}\left(\frac{b^2}{16}
 - \frac{a^2 b}{8}
 + \frac{5 a^4}{128}\right)
 \;=\;
 \frac{\Delta\,(45\alpha^2 + 2\Delta)}{60\,\kappa^3}.
 \label{eq:S6-extract}
\end{align}
The key structural observation: every coefficient $S_r$
with $r \ge 4$ contains $\Delta$ as a factor. This is the
universal factorization property of
Theorem~\ref{thm:dichotomy}.
\end{example}

\begin{remark}[Comparison with direct
recursion]\label{rem:direct-vs-closed}
The MC recursion (Lemma~\ref{lem:mc-recursion}) computes
$S_r$ as a sum of $\lfloor(r-2)/2\rfloor$ terms, each
involving products of previously computed coefficients.
The closed-form extraction
$S_r = \frac{1}{r}[t^{r-2}]\sqrt{Q_L}$ replaces this
recursive computation with a single coefficient extraction
from a degree-$2$ algebraic function. The computational
advantage grows with~$r$: at degree~$10$, the recursion
involves $4$ nonlinear terms, while the extraction requires
expanding $\sqrt{1 + u}$ to order~$8$.

For computer algebra, both methods are efficient. The
conceptual advantage of the closed form is structural: it
makes the $\Delta$-divisibility and the asymptotic behaviour
manifest, and it identifies the shadow tower as a branch of
an algebraic curve rather than as the output of a recursion.
\end{remark}


\subsection{The spectral curve}\label{subsec:spectral-curve}

The algebraicity theorem identifies the shadow obstruction
tower with a point on an algebraic curve.

\begin{definition}[Shadow spectral
curve]\label{def:spectral-curve}
The \emph{shadow spectral curve} of a primary line~$L$
with shadow metric $Q_L$ is the affine curve
\begin{equation}\label{eq:spectral-curve}
 \Sigma_L
 \;:=\;
 \bigl\{(t, y) \in \mathbb{A}^2 : y^2 = Q_L(t)\bigr\}.
\end{equation}
The shadow generating function $H(t) = t^2 y$ selects one
of the two branches of~$\Sigma_L$ near $t = 0$ (the branch
with $y(0) = 2\kappa > 0$).
\end{definition}

For class~$\mathbf{G}$, $Q_L = (2\kappa)^2$ is constant and
$\Sigma_L$ degenerates to two horizontal lines. For
class~$\mathbf{L}$, $Q_L$ has a double zero at
$t = -2\kappa/(3\alpha)$ and $\Sigma_L$ is a rational curve
with a node. For class~$\mathbf{M}$, $Q_L$ has two distinct
zeros, $\Sigma_L$ is a smooth affine conic (genus~$0$), and
the two branches of $\sqrt{Q_L}$ are interchanged by
monodromy around the branch points.

\begin{remark}[Comparison with matrix
models]\label{rem:matrix-models}
In the Eynard--Orantin topological recursion, the spectral
curve is prescribed as input and determines the full
genus expansion. Here the situation is reversed: the vertex
algebra determines the shadow metric $Q_L$ (three numbers),
which determines the spectral curve $\Sigma_L$, which
determines the shadow tower. The spectral curve is an
\emph{output}, not an input.

The structural parallel is precise: the recursion kernel of
Eynard--Orantin is the Bergman kernel on the spectral
curve; the MC recursion of the shadow tower is the
convolution kernel on the parameter line, and the
algebraicity theorem identifies the two. The genus-$g$
free energies $F_g$ are computed from the shadow tower by
stable graph summation, paralleling the Feynman expansion
from the spectral curve in matrix models. The key
difference is that the shadow spectral curve has genus~$0$
(it is a conic), while matrix-model spectral curves may
have arbitrary genus.
\end{remark}


% ================================================================
% 5. THE FOUR-CLASS CLASSIFICATION
% ================================================================

\section{The four-class classification}\label{sec:class}

\subsection{The single-line dichotomy}
\label{subsec:single-line}

\begin{theorem}[Single-line
dichotomy]\label{thm:dichotomy}
On any primary line $L$ with $\kappa \ne 0$ and critical
discriminant $\Delta = 8\kappa S_4$, the shadow depth
satisfies $r_{\max} \in \{2, 3, \infty\}$:
\begin{enumerate}[(i)]
\item $\Delta = 0$, $\alpha = 0$:
 $Q_L = (2\kappa)^2$; $H(t) = 2\kappa\, t^2$;
 $r_{\max} = 2$ (class~$\mathbf{G}$, Gaussian).
\item $\Delta = 0$, $\alpha \ne 0$:
 $Q_L = (2\kappa + 3\alpha\, t)^2$;
 $H(t) = t^2(2\kappa + 3\alpha\, t)$;
 $r_{\max} = 3$ (class~$\mathbf{L}$, Lie).
\item $\Delta \ne 0$:
 $Q_L$ is irreducible;
 $S_r \ne 0$ for all $r \ge 4$ generically;
 $r_{\max} = \infty$ (class~$\mathbf{M}$, mixed).
\end{enumerate}
For $r \ge 4$: $S_r = \Delta \cdot R_r$ with
$R_r \in \mathbb{Q}(\alpha, \kappa, S_4)$ a universal
rational function independent of the algebra family.
\end{theorem}

\begin{proof}
The Gaussian decomposition
$Q_L(t) = (2\kappa + 3\alpha\, t)^2 + 2\Delta\, t^2$
shows that $Q_L$ is a perfect square if and only if
$\Delta = 0$.

\emph{Case $\Delta = 0$.}
Then $Q_L = (2\kappa + 3\alpha\, t)^2$, so
$H(t) = t^2(2\kappa + 3\alpha\, t)$.
If also $\alpha = 0$, then $H(t) = 2\kappa\, t^2$;
$S_r = 0$ for all $r \ge 3$ (class~$\mathbf{G}$).
If $\alpha \ne 0$, then $H$ is a cubic polynomial;
$S_r = 0$ for $r \ge 4$, while $S_3 = \alpha \ne 0$
(class~$\mathbf{L}$).

\emph{Case $\Delta \ne 0$.}
Then $Q_L$ has two distinct zeros, so $\sqrt{Q_L(t)}$ is not
polynomial. The Taylor coefficients
$S_r = \frac{1}{r}[t^{r-2}]\sqrt{Q_L(t)}$ are nonzero
for all $r \ge 4$ generically (the only way
$[t^n]\sqrt{Q_L}$ vanishes identically is if $\sqrt{Q_L}$
is polynomial, requiring $\Delta = 0$). Hence
$r_{\max} = \infty$ (class~$\mathbf{M}$).

For the $\Delta$-divisibility: write
$\sqrt{Q_L} = |2\kappa + 3\alpha\, t|\,
(1 + 2\Delta\, t^2/(2\kappa + 3\alpha\, t)^2)^{1/2}$
and expand the binomial series. Every term beyond zeroth
order contributes a factor of~$\Delta$.
\end{proof}

\subsection{Class $\mathbf{C}$: contact
depth}\label{subsec:class-C}

A fourth class escapes the single-line dichotomy. The
$\beta\gamma$ system exhibits the mechanism concretely.

The $\beta\gamma$ algebra at conformal weight $\lambda$ has
two generators $\beta$ (weight $\lambda$) and $\gamma$
(weight $1 - \lambda$) with OPE
$\beta(z)\gamma(w) \sim 1/(z-w)$: a single simple pole,
no self-OPE. The central charge is
$c_{\beta\gamma}(\lambda) = 2(6\lambda^2 - 6\lambda + 1)$
and the modular characteristic is
$\kappa(\beta\gamma) = c_{\beta\gamma}/2
= 6\lambda^2 - 6\lambda + 1$.
At $\lambda = 1$: $\kappa = 1$, $c = 2$.

The $r$-matrix is the constant permutation
$r(z) = P/z$ (the transposition of $\beta$ and $\gamma$
channels): a simple pole, one order below the OPE double pole
(the $d\log$ propagator absorbs one pole as always). The
cubic shadow $S_3 = 0$: both $\beta$ and $\gamma$ are
free fields with no self-coupling, so the Lie bracket
component vanishes. The discriminant
$\Delta = 8\kappa S_4$: since $S_3 = 0$, the single-line
dichotomy gives $\Delta = 0$ on every one-dimensional
primary line. The single-line tower terminates at degree~$2$.

The quartic contact invariant $S_4 \ne 0$ lives on the
\emph{two-field} stratum: the cross-channel
$\beta\gamma\beta\gamma$ four-point amplitude produces a
nonzero quartic residue from the interaction of two distinct
fields. Concretely, the degree-$4$ MC component
$\Theta^{(4)}_{\beta\gamma}$ has a nonvanishing contribution
from the graph with two $\beta\gamma$ propagators crossing;
this contribution is proportional to $(1 - 2\lambda)^2$ and
vanishes only at $\lambda = 1/2$ (the symplectic boson, where
$c = -1$ and the theory acquires enhanced symmetry). The
quintic and all higher obstructions vanish: the cross-channel
self-bracket exits the two-field complex by dimensional
constraint (a rank-one target cannot support a quintic
interaction with the given conformal weight pairing).

\begin{definition}[Class $\mathbf{C}$]\label{def:class-C}
A chirally Koszul vertex algebra is of \emph{class
$\mathbf{C}$} (contact, $r_{\max} = 4$, $d_{\mathrm{alg}} = 2$)
if:
\begin{enumerate}[(a)]
\item on every one-dimensional primary line $L$, the shadow
 metric $Q_L$ is a perfect square ($\Delta_L = 0$), so
 the single-line tower terminates;
\item in the full multi-field deformation complex, a nonzero
 quartic contact invariant $S_4^{\mathrm{cross}} \ne 0$
 arises from the cross-channel interaction of distinct
 generators at complementary conformal weights;
\item the quintic and all higher obstructions vanish by
 dimensional constraint on the cross-channel complex.
\end{enumerate}
The $\beta\gamma$ system at generic $\lambda \ne 1/2$ is the
archetype: $S_3 = 0$ (no Lie bracket), $S_4 \ne 0$ (quartic
contact from $\beta\gamma$ crossing), $S_r = 0$ for $r \ge 5$
(dimensional termination). The algebraic depth is
$d_{\mathrm{alg}} = 2$: finite but nonzero, distinguishing
class~$\mathbf{C}$ from both class~$\mathbf{G}$
($d_{\mathrm{alg}} = 0$) and class~$\mathbf{L}$
($d_{\mathrm{alg}} = 1$).
\end{definition}

The four classes partition the Koszul world:

\begin{center}
\renewcommand{\arraystretch}{1.3}
\begin{tabular}{ccccl}
\toprule
Class & $\alpha$ & $\Delta$ & $r_{\max}$ & Examples \\
\midrule
$\mathbf{G}$ & $0$ & $0$ & $2$
 & Heisenberg, lattice, free fermion \\
$\mathbf{L}$ & $\ne 0$ & $0$ & $3$
 & Affine Kac--Moody (all types) \\
$\mathbf{C}$ & $0^*$ & $0^*$ & $4$
 & $\beta\gamma$, $bc$ ghosts \\
$\mathbf{M}$ & $\ne 0$ & $\ne 0$ & $\infty$
 & Virasoro, $\cW_N$ ($N \ge 3$) \\
\bottomrule
\end{tabular}
\end{center}

\noindent All four classes are Koszul: shadow depth classifies
the complexity of the OPE within the Koszul world, not
Koszulness itself.


% ================================================================
% 6. THE DEPTH GAP THEOREM
% ================================================================

\section{The depth gap theorem}\label{sec:depth-gap}

\subsection{Statement and proof}\label{subsec:depth-gap-proof}

\begin{definition}[Algebraic
depth]\label{def:algebraic-depth}
The \emph{algebraic depth} of a chirally Koszul vertex
algebra~$\cA$ is
\[
 d_{\mathrm{alg}}(\cA)
 \;:=\;
 \begin{cases}
  0 & \text{if } S_r = 0 \text{ for all } r \ge 3
  \quad(\text{class } \mathbf{G}), \\
  1 & \text{if } \alpha \ne 0 \text{ and } S_r = 0
  \text{ for all } r \ge 4
  \quad(\text{class } \mathbf{L}), \\
  2 & \text{if } r_{\max} = 4
  \quad(\text{class } \mathbf{C}), \\
  \infty & \text{if } r_{\max} = \infty
  \quad(\text{class } \mathbf{M}).
 \end{cases}
\]
\end{definition}

\begin{proposition}[Depth gap]\label{prop:depth-gap}
Let $\cA$ be a chirally Koszul algebra in the standard
landscape. Then
\begin{equation}\label{eq:depth-gap}
 d_{\mathrm{alg}}(\cA)
 \;\in\;
 \{0,\, 1,\, 2,\, \infty\}.
\end{equation}
No value~$3$ or any finite value $\ge 3$ is realized.
The four permitted values correspond bijectively to the
classes $\mathbf{G}$, $\mathbf{L}$, $\mathbf{C}$,
$\mathbf{M}$.
\end{proposition}

\begin{proof}
The argument has two parts: the single-line Riccati
regime, and the global contact witness.

\emph{Single-line regime ($\kappa \ne 0$).}
On a primary line~$L$ with $\kappa \ne 0$, the shadow
generating function
$H(t) = t^2\sqrt{Q_L(t)}$ is determined by
$(\kappa, \alpha, S_4)$
(Theorem~\ref{thm:riccati}), where
$Q_L(t) = (2\kappa + 3\alpha\, t)^2 + 2\Delta\, t^2$.

If $\Delta = 0$ (equivalently $S_4 = 0$), then
$Q_L = (2\kappa + 3\alpha\, t)^2$ is a perfect square, so
$H(t) = t^2(2\kappa + 3\alpha\, t)$ is a polynomial of
degree~$3$. The shadow coefficients $S_r = 0$ for all
$r \ge 4$. The tower terminates at degree at most~$3$,
giving $d_{\mathrm{alg}} \in \{0, 1\}$.

If $\Delta \ne 0$, then $Q_L$ is not a perfect square.
The binomial series for
$(1 + u)^{1/2}$ with
$u = (6\alpha\, t + (9\alpha^2 + 2\Delta)\, t^2)/(4\kappa^2)$
does not terminate when the coefficient of~$t^2$
in~$u$ is nonzero. Therefore $S_r \ne 0$ for infinitely
many~$r$, and $d_{\mathrm{alg}} = \infty$.

No single primary line with $\kappa \ne 0$ can realize
$d_{\mathrm{alg}} = 2$ or~$3$: the sharp dichotomy
$\Delta = 0 \Rightarrow$ polynomial,
$\Delta \ne 0 \Rightarrow$ infinite, admits no intermediate
depth.

The degree-$4$ master equation isolates the quartic
cancellation:
\[
 S_4^{\mathrm{free}}
 \;=\;
 -\frac{9\alpha^2}{16\kappa},
\]
so the finite-depth locus at degree~$4$ is exactly
$S_4 = 0$, the Jacobi locus. Once this cancellation fails,
the next recursion step forces a tail. At $r = 5$:
\[
 S_5 = -\frac{6}{5}\, P\, \alpha\, S_4.
\]
For $\kappa \ne 0$, $\alpha \ne 0$, and $S_4 \ne 0$,
this is nonzero. If $\alpha = 0$ but $S_4 \ne 0$, then
$S_5 = 0$ but
\[
 S_6
 \;=\;
 -\frac{P}{12} \cdot \frac{1}{2} \cdot 16\, S_4^2
 \;=\;
 -\frac{2}{3}\, P\, S_4^2
 \;\ne\; 0.
\]
Either way, the tower does not terminate at
degree~$5$.

\emph{Global contact witness.}
The value $d_{\mathrm{alg}} = 2$ is realized by the
$\beta\gamma$ system (equivalently $bc$), whose global
shadow data on the conformal-weight family is
$S_2 = 6\lambda^2 - 6\lambda + 1$, $S_3 = 0$,
$S_4 = -5/12$ (on a charged stratum), $S_r = 0$ for
$r \ge 5$. The quartic contact invariant lives on a charged
stratum; the quintic obstruction vanishes by rank-one
rigidity (Definition~\ref{def:class-C}).
\end{proof}

\begin{remark}[Why no $d_{\mathrm{alg}} = 3$]
\label{rem:why-no-3}
The impossibility of $d_{\mathrm{alg}} = 3$ has a clean
algebraic explanation. On a single primary line, the tower
is controlled by a degree-$2$ algebraic function: either
the square root is rational (tower terminates at
degree~$\le 3$) or irrational (tower infinite). There is
no mechanism for a degree-$2$ algebraic function to produce
exactly one nonzero coefficient beyond the polynomial
envelope.

Off the primary line, the contact stratum supports at
most one additional degree (by rank-one rigidity). So
the maximum finite depth from the single-line polynomial
envelope ($r_{\max} = 3$) plus the contact stratum
correction ($+1$) gives $r_{\max} = 4$, class
$\mathbf{C}$. No further finite extension is possible.
\end{remark}

\begin{remark}[Generating depth vs algebraic
depth]\label{rem:gen-vs-alg}
The algebraic depth $d_{\mathrm{alg}}$ (whether the tower
terminates) is distinct from the generating depth
$d_{\mathrm{gen}}$ (the degree at which all higher
operations are determined by lower ones).
For Virasoro: $d_{\mathrm{gen}} = 3$ (the cubic
shadow~$\alpha = 2$ generates all higher shadows recursively
from $(\kappa, \alpha, S_4)$), but
$d_{\mathrm{alg}} = \infty$ (all $S_r$ are nonzero).
\end{remark}

\subsection{The Jacobi mechanism}\label{subsec:jacobi}

The vanishing $S_4 = 0$ for class~$\mathbf{L}$ has a
precise algebraic explanation that illuminates the
depth gap.

\begin{proposition}[Jacobi kills the quartic
shadow]\label{prop:jacobi-kills-s4}
Let $\cA = V_k(\fg)$ be an affine Kac--Moody vertex
algebra. The quartic collision residue on $\FM_4(X)$,
restricted to the primary line, is proportional to
\[
 \mathcal{J}
 \;:=\;
 \sum_{\mathrm{cyclic}(a,b,c)}
 f^{abx}\, f^{xcd},
\]
where $f^{abc}$ are the structure constants of~$\fg$.
The Jacobi identity
$[[J^a, J^b], J^c] + \mathrm{cyclic} = 0$
forces $\mathcal{J} = 0$, hence $S_4 = 0$.
\end{proposition}

\begin{proof}
The degree-$4$ component of the MC element~$\Theta_\cA$
receives contributions from trees with one internal edge
connecting two cubic vertices. Each cubic vertex contributes
$f^{abc}$; the internal edge carries the propagator
$P = 1/\kappa$; and the contraction over the internal
index~$x$ produces $f^{abx} f^{xcd}$. The master equation
at degree~$4$ reads
$d_0(\Sh_4) + [\Sh_3, \Sh_3] = 0$.

The bracket $[\Sh_3, \Sh_3]$ evaluates on the primary line as
a sum over channel pairings, and the coefficient of $S_4$
in the linear term $d_0(\Sh_4)$ is $8\kappa S_4$. The
obstruction is therefore
$8\kappa S_4 = -[\Sh_3, \Sh_3]|_{\deg 4}$, which is
proportional to~$\mathcal{J}$. Since $\mathcal{J} = 0$ by
the Jacobi identity, $S_4 = 0$.
\end{proof}

\begin{remark}[The Jacobi locus as a
discriminant]\label{rem:jacobi-locus}
The vanishing of $S_4$ is not a numerical accident: it is
the statement that the Jacobi identity of the Lie
algebra~$\fg$ kills the quartic shadow. This explains the
structural difference between classes $\mathbf{L}$ and
$\mathbf{M}$: affine Kac--Moody algebras carry a Lie
bracket whose Jacobi identity kills $S_4$; Virasoro and
$\cW_N$ algebras carry OPE structures without a classical
Jacobi identity, and $S_4$ survives.

The Jacobi locus
$\{S_4 = 0\} \subset \{\text{shadow data}\}$ is a
codimension-$1$ subvariety of the parameter space. On this
locus, $\Delta = 0$ and the tower terminates. Off the
locus, $\Delta \ne 0$ and the tower is infinite. The
transition is sharp: the depth jumps from~$3$ to~$\infty$
with no intermediate value.
\end{remark}

\subsection{The shadow Lie algebra
perspective}\label{subsec:shadow-lie}

The depth gap admits a second proof through the structure of
the shadow algebra
$\cA^{\mathrm{sh}} := H_\bullet(
\Defcyc^{\mathrm{mod}}(\cA))$.

\begin{proposition}[Shadow Lie algebra
Jacobi]\label{prop:shadow-lie-jacobi}
The shadow algebra $\cA^{\mathrm{sh}}$ carries a graded Lie
bracket $[-,-]^{\mathrm{sh}}$. At degree~$4$, the Jacobi
identity in $\cA^{\mathrm{sh}}$ forces
\[
 [S_3, [S_3, S_3]^{\mathrm{sh}}]^{\mathrm{sh}}
 \;=\;
 0,
\]
and the degree-$5$ master equation reads
\[
 d_0(\Sh_5) + [S_3, S_4]^{\mathrm{sh}}
 + \tfrac{1}{2}[S_4, S_3]^{\mathrm{sh}} = 0.
\]
If $S_4 \ne 0$ and $S_3 \ne 0$, then $S_5 \ne 0$; if
$S_4 \ne 0$ and $S_3 = 0$, then $S_6 \ne 0$. In either
case, the tower does not terminate at degree~$5$.
\end{proposition}

This is the shadow Lie algebra Jacobi mechanism identified
in the monograph (Proposition~\ref{prop:depth-gap},
alternative proof): the graded Jacobi identity in the
shadow algebra constrains which depths are achievable.
The depth gap $d_{\mathrm{alg}} \ne 3$ is forced because
no configuration of $(S_3, S_4)$ with $S_4 \ne 0$ can
produce $S_5 = S_6 = \cdots = 0$.


% ================================================================
% 7. SC-FORMALITY AND CLASS G
% ================================================================

\section{SC-formality characterizes
class~$\mathbf{G}$}\label{sec:sc-formal}

The Swiss-cheese operad $\mathrm{SC}^{\mathrm{ch,top}}$ is a
two-coloured operad governing the interaction of a bulk
algebra (closed colour) with a boundary algebra (open
colour). A chirally Koszul vertex algebra~$\cA$ is
\emph{SC-formal} if the transferred Swiss-cheese operations
$m_k^{\mathrm{SC}}$ vanish for all $k \ge 3$.

\begin{proposition}[SC-formality iff class
$\mathbf{G}$]\label{prop:sc-formal}
Let~$\cA$ be a chirally Koszul vertex algebra in the
standard landscape. Then~$\cA$ is SC-formal if and only
if~$\cA$ belongs to class~$\mathbf{G}$.
\end{proposition}

\begin{proof}
\emph{Forward direction (class $\mathbf{G}$
$\Longrightarrow$ SC-formal).}
Class~$\mathbf{G}$ consists of the Gaussian families:
Heisenberg, lattice VOAs, and free fermions. In each case,
the genus-$0$ mixed tree formulas stop at degree~$2$ (the
OPE is purely quadratic), so
$m_k^{\mathrm{SC}} = 0$ for $k \ge 3$. The argument uses
Wick factorization: all higher connected contractions
vanish for a free field.

\emph{Converse (SC-formal $\Longrightarrow$
class~$\mathbf{G}$).}
Suppose~$\cA$ is SC-formal. Then every mixed SC operation
of degree $r \ge 3$ vanishes. The tree-shadow
correspondence is operadic: the degree-$r$ mixed tree and
the degree-$r$ shadow are produced by the same genus-$0$
tree-transfer formula, with the same propagator on internal
edges. The only difference is the output colour, and passage
to the symmetric closed sector applies the averaging morphism
$\mathrm{av}(\Theta_\cA^{E_1}) = \Theta_\cA$.

Hence SC-formality forces $S_r(\cA) = 0$ for every
$r \ge 3$: a nonzero shadow coefficient would produce,
by the tree-shadow correspondence, a nontrivial mixed
degree-$r$ Massey product, contradicting SC-formality.

The condition $S_r = 0$ for all $r \ge 3$ is exactly class
$\mathbf{G}$ (Theorem~\ref{thm:dichotomy}(i)). Classes
$\mathbf{L}$, $\mathbf{C}$, and $\mathbf{M}$ are ruled out:
class $\mathbf{L}$ has $S_3 = \alpha \ne 0$; class
$\mathbf{C}$ has a nonzero quartic contact invariant;
class $\mathbf{M}$ has $S_r \ne 0$ for all $r \ge 4$.
\end{proof}

\begin{remark}[Operadic alternative]\label{rem:sc-operadic}
The proposition admits a second proof via genus-$0$
operadic transfer. Class $\mathbf{G}$ is the Gaussian
locus where the genus-$0$ transfer is generated by the
binary two-point kernel alone. Every connected SC tree of
degree $r \ge 3$ factors through binary pairings and
carries no primitive higher vertex, so the mixed operations
vanish. For the converse, if some $S_r \ne 0$ with
$r \ge 3$, the shadow-formality identification and the
averaging identity $\mathrm{av}(\Theta_\cA^{E_1}) =
\Theta_\cA$ exhibit a nonzero closed projection from a
nonzero tree, contradicting SC-formality.
\end{remark}


% ================================================================
% 7b. THE SHADOW-FORMALITY IDENTIFICATION
% ================================================================

\subsection{The shadow-formality
identification}\label{subsec:shadow-formality}

The shadow obstruction tower and the $A_\infty$
non-formality data of a chirally Koszul algebra are
identified by the homotopy transfer theorem.

\begin{proposition}[Shadow-formality
identification]\label{prop:shadow-formality}
Let $\cA$ be a chirally Koszul vertex algebra and
$\{m_k\}_{k \ge 2}$ the transferred $A_\infty$ structure
on the bar cohomology $H^\bullet(B(\cA))$, obtained by
homological perturbation from the bar differential. Then:
\begin{enumerate}[(i)]
\item The degree-$2$ shadow $S_2 = \kappa$ corresponds to
 the binary product $m_2$ (which is always nonzero for
 $\kappa \ne 0$).
\item The degree-$3$ shadow $S_3 = \alpha$ corresponds to
 the first $A_\infty$ correction $m_3$:
 $\alpha = 0$ if and only if $m_3 = 0$ on the primary
 line.
\item The degree-$r$ shadow $S_r$ corresponds to $m_r$:
 $S_r = 0$ if and only if $m_r = 0$ on the primary line
 (after projection from the full $A_\infty$ structure).
\item The shadow depth $d_{\mathrm{alg}}$ equals the
 $A_\infty$-depth: the minimal~$n$ such that
 $m_k = 0$ for all $k > n + 2$ on the primary line.
\end{enumerate}
\end{proposition}

\begin{proof}
The homotopy transfer theorem produces the $A_\infty$
structure $\{m_k\}$ by summing over binary trees with $k$
leaves, where each internal edge carries the homotopy
$h = P \cdot d_0^{-1}$ (with $P = 1/\kappa$ the propagator)
and each internal vertex carries the binary product $m_2$.
The degree-$r$ component of the MC element
$\Theta_\cA$ is the generating function of the tree sum
at degree~$r$. The shadow coefficient $S_r$ is the
scalar restriction of this tree sum to the primary line.

The identification $S_r \leftrightarrow m_r|_L$ follows
from the fact that the MC recursion
(Proposition~\ref{prop:master-eq}) and the $A_\infty$
recursion (the Stasheff identities) are the same system
of equations, restricted to different components of the
deformation complex. On the primary line, both reduce to
the scalar recursion of Lemma~\ref{lem:mc-recursion}.
\end{proof}

\begin{corollary}[Shadow depth classifies
$A_\infty$-formality]\label{cor:formality-class}
The four shadow depth classes correspond to four levels of
$A_\infty$-formality:
\begin{center}
\renewcommand{\arraystretch}{1.3}
\begin{tabular}{cccl}
\toprule
Class & $d_{\mathrm{alg}}$ & $A_\infty$ formality & Status \\
\midrule
$\mathbf{G}$ & $0$ & formal ($m_k = 0$, $k \ge 3$) & SC-formal \\
$\mathbf{L}$ & $1$ & $m_3 \ne 0$, $m_k = 0$ for $k \ge 4$
 & $1$-formal \\
$\mathbf{C}$ & $2$ & $m_4 \ne 0$ on charged stratum
 & $2$-formal \\
$\mathbf{M}$ & $\infty$ & $m_k \ne 0$ for all $k$
 & non-formal \\
\bottomrule
\end{tabular}
\end{center}
\noindent The depth gap $d_{\mathrm{alg}} \ne 3$ translates
to: no chirally Koszul vertex algebra has $A_\infty$-depth
exactly~$3$ (i.e., $m_3 \ne 0$, $m_4 \ne 0$, but
$m_5 = m_6 = \cdots = 0$).
\end{corollary}

\begin{remark}[The $L_\infty$ perspective]
\label{rem:linfty-formality}
The modular convolution algebra
$\Defcyc^{\mathrm{mod}}(\cA)$ is a dg Lie algebra, and
the MC element $\Theta_\cA$ determines a gauge equivalence
class. The shadow-formality identification has an
$L_\infty$ counterpart: the transferred $L_\infty$ structure
$\{\ell_k\}_{k \ge 2}$ on the shadow algebra
$\cA^{\mathrm{sh}}$ satisfies $\ell_k = 0$ for $k \ge 3$ if
and only if $\cA$ is class~$\mathbf{G}$. The $L_\infty$
formality level coincides with the $A_\infty$-depth on
the primary line.

The passage from $A_\infty$ to $L_\infty$ is the
passage from the open colour (boundary, $A_\infty$) to the
closed colour (bulk, $L_\infty$) in the Swiss-cheese
structure. SC-formality is the simultaneous vanishing of
both the $A_\infty$ and $L_\infty$ higher operations, which
occurs exactly in class~$\mathbf{G}$
(Proposition~\ref{prop:sc-formal}).
\end{remark}

\begin{remark}[Massey products]
\label{rem:massey}
The $A_\infty$ correction $m_3$ is a Massey product: it
measures the failure of the binary product to be
associative up to homotopy. For class~$\mathbf{L}$
(affine Kac--Moody), $m_3 \ne 0$ because the Lie bracket
introduces a cubic non-associativity. For class~$\mathbf{G}$
(Heisenberg), $m_3 = 0$ because the abelian OPE is strictly
associative.

The quartic correction $m_4$ is a higher Massey product,
and its non-vanishing for class~$\mathbf{M}$ is the
statement that the Virasoro OPE has an irreducible
four-point contact term not generated by lower operations.
The depth gap theorem says that this contact term, once
present, forces an infinite sequence of higher Massey
products: no finite truncation is possible beyond
degree~$4$.
\end{remark}


% ================================================================
% 8. THE E_1 FRAMING
% ================================================================

\section{The $E_1$ framing}\label{sec:e1}

The shadow obstruction tower is the $\Sigma_n$-coinvariant
projection of a richer $E_1$ structure.

\begin{proposition}[$E_1$ framing of the shadow
tower]\label{prop:e1-framing}
The shadow obstruction tower $\{S_r\}_{r \ge 2}$ is the
image of the ordered $R$-matrix tower under the averaging
map $\mathrm{av}\colon \mathfrak{g}^{E_1}_\cA \to
\mathfrak{g}^{\mathrm{mod}}_\cA$. At degree~$2$:
\begin{enumerate}[(i)]
\item For abelian algebras (Heisenberg, free fields):
 $\kappa = \mathrm{av}(r(z))$, where
 $r^{\mathrm{Heis}}(z) = k/z$ (level prefix $k$ mandatory;
 at $k = 0$, the $r$-matrix vanishes).

\item For non-abelian affine Kac--Moody $V_k(\fg)$
 (trace-form convention $r(z) = k\Omega/z$, where $\Omega$
 is the inverse Killing form Casimir):
 \[
  \kappa(V_k(\fg))
  \;=\;
  \mathrm{av}(r(z)) + \frac{\dim(\fg)}{2}
  \;=\;
  \frac{k\,\dim(\fg)}{2h^\vee}
  + \frac{\dim(\fg)}{2}
  \;=\;
  \frac{\dim(\fg)\,(k + h^\vee)}{2h^\vee}.
 \]
 The $\dim(\fg)/2$ term is the Sugawara shift, the
 simple-pole self-contraction through the adjoint Casimir
 eigenvalue~$2h^\vee$.
\end{enumerate}
At higher degrees, the averaging map sends the degree-$r$
component of the ordered $R$-matrix tower
$\Theta_\cA^{E_1}$ to the shadow $S_r$.
\end{proposition}

\begin{proof}
The ordered bar $B^{\mathrm{ord}}(\cA) =
T^c(s^{-1}\bar{\cA})$ carries the $E_1$ convolution algebra
$\mathfrak{g}^{E_1}_\cA$, whose MC element
$\Theta_\cA^{E_1}$ projects to~$\Theta_\cA$ under
$\mathrm{av}$. The degree-$2$ component of
$\Theta_\cA^{E_1}$ is the $R$-matrix $r(z)$, and its
average is the scalar degree-$2$ shadow: it equals the modular
characteristic~$\kappa$ on abelian families, while for
non-abelian affine Kac--Moody the average is
$\kappa_{\mathrm{dp}}$ and the full~$\kappa$ adds the
Sugawara correction from the simple-pole self-contraction of the
ordered two-point function. The higher degrees follow by the same averaging
applied to the degree-$r$ MC components.
\end{proof}

\begin{remark}[What averaging
discards]\label{rem:averaging-lossy}
The averaging map $\mathrm{av}$ is lossy: it discards the
spectral parameter~$z$, the $R$-matrix structure, and the
Yangian datum. The full $E_1$ tower retains the braiding of
line operators; the shadow tower retains only the scalar
invariants. For the Virasoro $R$-matrix
$r^{\Vir}(z) = (c/2)/z^3 + 2T/z$,
averaging produces
$\kappa = c/2$ and $\alpha = 2$, but discards the cubic
pole structure that encodes the braided tensor category of
Virasoro modules.
\end{remark}

\subsection{The $R$-matrix tower for each
class}\label{subsec:r-matrix-tower}

The $E_1$ convolution algebra
$\mathfrak{g}^{E_1}_\cA$ carries a spectral $R$-matrix
$r(z) \in \mathrm{End}(V \otimes V)((z))$ at degree~$2$,
whose averaging produces the scalar degree-$2$ shadow:
for abelian families this is~$\kappa$, while for non-abelian
affine Kac--Moody it is~$\kappa_{\mathrm{dp}}$ and the full
modular characteristic is
$\kappa = \kappa_{\mathrm{dp}} + \dim(\fg)/2$. The structure of $r(z)$
varies by class.

\begin{example}[Class $\mathbf{G}$:
Heisenberg]\label{ex:r-heis}
$r^{\mathrm{Heis}}(z) = k/z$. One pole of order~$1$ (from
the double-pole OPE $J(z)J(w) \sim k/(z-w)^2$ after
$d\!\log$ absorption). Averaging:
$\mathrm{av}(k/z) = k = \kappa$. The $R$-matrix is scalar,
the Yangian is trivial, and the shadow tower has
$\alpha = 0$, $\Delta = 0$. At $k = 0$, the $R$-matrix
vanishes; this is the abelian limit, not the critical level
(which does not exist for Heisenberg).
\end{example}

\begin{example}[Class $\mathbf{L}$: affine
$\mathfrak{sl}_2$]\label{ex:r-sl2}
In the trace-form convention:
$r^{\mathrm{KM}}(z) = k\,\Omega/z$, where
$\Omega = \sum_a t^a \otimes t_a$ is the inverse Killing
form Casimir.
Averaging over $\Sigma_2$-coinvariants:
$\mathrm{av}(k\Omega/z)
= k\,\mathrm{tr}(\Omega)/\dim(\fg)
\cdot \dim(\fg)/(2h^\vee)
= k\,\dim(\fg)/(2h^\vee)$.
The full $\kappa$ includes the Sugawara shift
$\dim(\fg)/2$, the simple-pole self-contraction through
the adjoint Casimir eigenvalue~$2h^\vee$.

For $\fg = \mathfrak{sl}_2$: $\dim(\fg) = 3$,
$h^\vee = 2$, $\kappa = 3(k+2)/4$.
At $k = 0$: $\kappa = 3/2 \ne 0$ (the Lie bracket
persists; the $R$-matrix
$r(z)|_{k=0} = 0$ in trace-form convention, but the
Sugawara shift $\dim(\fg)/2 = 3/2$ survives).
At $k = -2$ (critical): $\kappa = 0$, the center jumps,
Koszulness fails.

The cubic shadow $\alpha$ encodes the Lie bracket structure
constants: the Killing $3$-cocycle
$\omega_3(X,Y,Z) = \mathrm{tr}([X,Y]Z)$.
\end{example}

\begin{example}[Class $\mathbf{M}$:
Virasoro]\label{ex:r-vir}
The Virasoro OPE
$T(z)T(w) \sim (c/2)/(z-w)^4 + 2T(w)/(z-w)^2
+ \partial T(w)/(z-w)$ has poles at orders $4$, $2$, $1$.
After $d\!\log$ absorption (which reduces each pole order
by~$1$):
\[
 r^{\Vir}(z)
 \;=\;
 \frac{c/2}{z^3} + \frac{2T}{z}.
\]
The cubic pole is characteristic of class~$\mathbf{M}$
and encodes the nonlinear structure of the Virasoro algebra.
Averaging:
$\mathrm{av}((c/2)/z^3) = 0$ (odd-degree pole averages
to zero on the primary line), and
$\mathrm{av}(2T/z) = c/2 = \kappa$. The cubic shadow
$\alpha = 2$ comes from the simple pole coefficient
$2T$, whose projection onto the primary line is the scalar~$2$.

The quartic shadow $S_4 = 10/[c(5c+22)]$ encodes the
four-point contact term in the Virasoro OPE, which is
not killed by any Jacobi-type identity. This is the
mechanism that produces $\Delta \ne 0$ and the infinite
tower.
\end{example}

\begin{remark}[Pole order and shadow
class]\label{rem:pole-class}
The $R$-matrix pole order determines the shadow class:
\begin{center}
\renewcommand{\arraystretch}{1.2}
\begin{tabular}{cccl}
\toprule
OPE pole order $p$ & $R$-matrix pole $p{-}1$
 & Class & Example \\
\midrule
$2$ & $1$ & $\mathbf{G}$
 & Heisenberg \\
$2$ (non-abelian) & $1$ & $\mathbf{L}$
 & affine KM \\
$1$ (charged) & $0$ & $\mathbf{C}$
 & $\beta\gamma$ \\
$4$ & $3$ & $\mathbf{M}$
 & Virasoro \\
\bottomrule
\end{tabular}
\end{center}
The $d\!\log$ absorption shifts all pole orders by~$-1$.
The Jacobi identity kills the quartic contribution for
class~$\mathbf{L}$ (where the non-abelian cubic OPE comes
from a Lie bracket). For class~$\mathbf{M}$, the quartic
OPE survives.
\end{remark}

\subsection{The ordered-to-symmetric
functor}\label{subsec:ord-to-sym}

The averaging map $\mathrm{av}\colon
\mathfrak{g}^{E_1}_\cA \to
\mathfrak{g}^{\mathrm{mod}}_\cA$ is the
$\Sigma_n$-coinvariant projection. At the level of
generating functions, it sends the spectral $R$-matrix
tower
$\Theta_\cA^{E_1}(z) = \sum_{r \ge 2} \Sh_r^{E_1}(z)$
to the shadow tower
$\Theta_\cA = \sum_{r \ge 2} S_r \cdot x^r$ by discarding
the spectral parameter~$z$ and averaging over permutations.

\begin{proposition}[Structure of the averaging
map]\label{prop:av-structure}
The averaging map has the following properties:
\begin{enumerate}[(i)]
\item \emph{Surjectivity.}\enspace
 Every shadow invariant $S_r$ is in the image of
 $\mathrm{av}$; the ordered tower lifts the symmetric tower.

\item \emph{Non-injectivity.}\enspace
 The kernel $\ker(\mathrm{av})$ contains all
 $\Sigma_n$-anti-invariant components, including the
 $R$-matrix braiding data, the Yangian coproduct corrections,
 and the non-commutative shadow coefficients at each degree.

\item \emph{Compatibility with MC.}\enspace
 $\mathrm{av}$ is a morphism of dg Lie algebras:
 $\mathrm{av}([\xi,\eta]^{E_1})
 = [\mathrm{av}(\xi), \mathrm{av}(\eta)]^{\mathrm{mod}}$.
 The MC element projects:
 $\mathrm{av}(\Theta_\cA^{E_1}) = \Theta_\cA$.

\item \emph{Degree-$2$ content.}\enspace
 At degree~$2$: $\mathrm{av}(r(z)) = \kappa$ for abelian
 algebras; $\mathrm{av}(r(z)) + \dim(\fg)/2 = \kappa$ for
 non-abelian KM (Sugawara shift).
\end{enumerate}
\end{proposition}

\begin{remark}[What the shadow tower
cannot see]\label{rem:shadow-blind}
The shadow tower is blind to: the spectral parameter~$z$ and
all $z$-dependent structure; the $R$-matrix and its Yang--Baxter
equation; the Yangian coproduct; the braided tensor category
structure on modules; and the Stokes phenomena of the
KZ connection. These richer structures live in the $E_1$
tower and are invisible to the $\Sigma_n$-coinvariant
projection. The shadow tower sees only the
\emph{scalar invariants} of the $E_1$ structure: the
numbers $\kappa$, $\alpha$, $S_4$, and their higher-degree
extensions.

This is the trade-off of averaging: the shadow tower is
computable (finite-dimensional at each degree), algebraic
(degree~$2$), and universal (the same structural theorems
hold for all families); but it discards the
representation-theoretic content that distinguishes, for
example, the Yangian of $\mathfrak{sl}_2$ from the Yangian
of $\mathfrak{sl}_3$.
\end{remark}


% ================================================================
% 9. THE SHADOW CONNECTION
% ================================================================

\section{The shadow connection and Koszul
monodromy}\label{sec:connection}

\subsection{The shadow connection}\label{subsec:conn}

\begin{definition}[Shadow
connection]\label{def:shadow-conn}
The \emph{shadow connection} on
$L \setminus \{Q_L = 0\}$ is
\begin{equation}\label{eq:shadow-conn}
 \nabla^{\mathrm{sh}}
 \;:=\;
 d \;-\; \frac{Q_L'(t)}{2\,Q_L(t)}\,dt
 \;=\;
 d - d\!\left(\log\sqrt{Q_L}\right).
\end{equation}
\end{definition}

\begin{proposition}[Properties]\label{prop:conn-properties}
\leavevmode
\begin{enumerate}[(i)]
\item \emph{Flatness.}\enspace
 $(\nabla^{\mathrm{sh}})^2 = 0$.

\item \emph{Regular singularities.}\enspace
 At each zero $t_0$ of~$Q_L$, the connection form has a
 simple pole with residue~$\frac{1}{2}$.

\item \emph{Koszul monodromy.}\enspace
 The monodromy around each branch point is
 $\exp(2\pi i \cdot \frac{1}{2}) = -1$.

\item \emph{Flat section.}\enspace
 The unique flat section with $\Phi(0) = 1$ is
 \begin{equation}\label{eq:flat-section}
  \Phi(t)
  \;=\;
  \frac{\sqrt{Q_L(t)}}{2\kappa},
 \end{equation}
 and
 $H(t) = 2\kappa\, t^2\, \Phi(t)$.
\end{enumerate}
\end{proposition}

\begin{proof}
Part~(i): $Q_L'/(2Q_L)\, dt$ is closed on a
one-dimensional base. Parts~(ii)--(iii): standard
properties of the Gauss--Manin connection of
$\{\sqrt{Q_L(t)}\}_t$. Near a simple zero $t_0$ of $Q_L$:
$Q_L(t) \approx Q_L'(t_0)(t - t_0)$, so
$Q_L'/(2Q_L)\, dt \approx dt/(2(t - t_0))$. The monodromy
of a square root around a simple branch point is~$-1$.
Part~(iv): direct verification.
\end{proof}

\subsection{Koszul duality and discriminant
complementarity}\label{subsec:complementarity}

\begin{proposition}[Discriminant
complementarity]\label{prop:complementarity}
For the Virasoro family ($\Vir_c^! = \Vir_{26-c}$):
\begin{equation}\label{eq:complementarity}
 \Delta(\Vir_c) + \Delta(\Vir_{26-c})
 \;=\;
 \frac{6960}{(5c+22)(152-5c)}.
\end{equation}
At the self-dual point $c = 13$:
$\Delta(13) = 40/87$.
\end{proposition}

\begin{proof}
\begin{align*}
 \Delta(c) + \Delta(26-c)
 &\;=\;
 \frac{40}{5c+22} + \frac{40}{152-5c} \\
 &\;=\;
 40 \cdot \frac{174}{(5c+22)(152-5c)}
 \;=\;
 \frac{6960}{(5c+22)(152-5c)}.
 \qedhere
\end{align*}
\end{proof}

\begin{remark}[The self-dual point]\label{rem:self-dual}
At $c = 13$, the Koszul involution
$\Vir_c \mapsto \Vir_{26-c}$ has a fixed point:
$\Vir_{13}$ is self-dual under the uncurved quadratic
Koszul duality $c \mapsto 26 - c$.
The shadow metric satisfies
$Q_L = Q_L^!$, the connection is self-dual, and the
shadow growth rates are equal: $\rho = \rho^!$.
The Koszul complementarity $\kappa + \kappa' = 13$
(Virasoro-specific; not universal) reflects this fixed-point
structure.
\end{remark}

\subsection{The Verdier involution on shadow
data}\label{subsec:verdier}

Under Koszul duality $\cA \mapsto \cA^!$, the shadow data
transforms. The transformation depends on the family.

\begin{proposition}[Koszul transformation of shadow
data]\label{prop:koszul-transform}
\leavevmode
\begin{enumerate}[(i)]
\item \emph{Affine Kac--Moody.}\enspace
 $V_k(\fg)^! = \mathrm{CE}^{\mathrm{ch}}
 (\fg_{-k-2h^\vee})$.
 The shadow data transforms as
 $\kappa \mapsto -\kappa$, $\alpha \mapsto -\alpha$,
 $S_4 \mapsto 0$. The Koszul conductor is
 $\kappa(V_k(\fg)) + \kappa(V_k(\fg)^!) = 0$.

\item \emph{Virasoro.}\enspace
 $\Vir_c^! = \Vir_{26-c}$. The shadow data transforms by
 $c \mapsto 26 - c$:
 $\kappa \mapsto (26-c)/2$, $\alpha \mapsto 2$
 (invariant), $S_4 \mapsto 10/[(26-c)(152-5c)]$.
 The Koszul conductor is
 $\kappa(\Vir_c) + \kappa(\Vir_{26-c}) = 13$.

\item \emph{Heisenberg.}\enspace
 $\cH_k^! = \mathrm{Sym}^{\mathrm{ch}}(V^*)$
 (the chiral symmetric algebra on the dual).
 Note: $\cH_k^! \ne \cH_{-k}$ (these are isomorphic as
 vertex algebras but not as Koszul duals; the Koszul dual
 has the same $\kappa$). The Koszul conductor is
 $\kappa + \kappa' = 0$.
\end{enumerate}
\end{proposition}

\begin{remark}[Family-dependence of the Koszul
conductor]\label{rem:koszul-conductor}
The Koszul conductor $K(\cA) := \kappa(\cA) + \kappa(\cA^!)$
is \emph{not} universal: it depends on the family.
For affine Kac--Moody and free fields: $K = 0$.
For Virasoro: $K = 13$.
For $\cW_3$: $K = 250/3$.
For Bershadsky--Polyakov: $K = 98/3$.
The statement ``$\kappa + \kappa' = 0$'' without a family
qualifier is false for Virasoro and all $\cW_N$ with
$N \ge 2$.
\end{remark}

\subsection{The shadow connection for each
class}\label{subsec:conn-by-class}

The shadow connection
$\nabla^{\mathrm{sh}} = d - Q_L'/(2Q_L)\, dt$ has
qualitatively different behaviour in each class.

\begin{example}[Class $\mathbf{G}$:
trivial]\label{ex:conn-G}
$Q_L = 4\kappa^2$, constant. $Q_L' = 0$, so
$\nabla^{\mathrm{sh}} = d$: the connection is trivial, the
flat section is $\Phi(t) = 1$, and there are no branch
points. The shadow tower carries no monodromy.
\end{example}

\begin{example}[Class $\mathbf{L}$: single branch
point]\label{ex:conn-L}
$Q_L = (2\kappa + 3\alpha\, t)^2$, a perfect square with a
double zero at $t_0 = -2\kappa/(3\alpha)$. The connection
form $Q_L'/(2Q_L)\, dt = 3\alpha/(2\kappa + 3\alpha\, t)\, dt$
has a simple pole with residue~$1$ at~$t_0$. The monodromy
is $\exp(2\pi i \cdot 1) = 1$: trivial.

The double zero of $Q_L$ means that $\sqrt{Q_L}$ is
single-valued (polynomial), consistent with the termination
of the tower at degree~$3$.
\end{example}

\begin{example}[Class $\mathbf{M}$:
Virasoro]\label{ex:conn-M}
$Q_L = (c + 6t)^2 + 80t^2/(5c+22)$. The two simple zeros
\[
 t_\pm
 \;=\;
 \frac{-3c \pm ic\sqrt{40/((5c+22)c^2 - 36)}}{
 9 + 40/(5c+22)}
\]
form a complex conjugate pair (for $c > 0$ real). Each has
residue~$1/2$; monodromy~$-1$. The shadow connection is the
Gauss--Manin connection of the double cover
$\Sigma_L = \{y^2 = Q_L(t)\}$, and the branch-point
monodromy is the deck transformation $y \mapsto -y$.
\end{example}


\subsection{The planted-forest
formula}\label{subsec:planted-forest}

The MC equation at genus~$2$ produces a codimension-$2$
boundary correction.

\begin{proposition}[Planted-forest
correction]\label{prop:planted-forest}
The genus-$2$ planted-forest correction on a primary line is
\begin{equation}\label{eq:planted-forest}
 \delta_{\mathrm{pf}}^{(2,0)}
 \;=\;
 \frac{S_3(10\, S_3 - \kappa)}{48},
\end{equation}
measuring the contribution of codimension-$2$ boundary
strata to the genus-$2$ amplitude. This correction vanishes
for class~$\mathbf{G}$ (where $S_3 = 0$). For Virasoro:
$\delta_{\mathrm{pf}}^{(2,0)} = -(c - 40)/48$.
\end{proposition}

\begin{proof}
The codimension-$2$ boundary strata of $\Mbar_{2,0}$
correspond to stable graphs with two internal edges. Each
edge carries the propagator $P = 1/\kappa$ and connects two
cubic vertices contributing $S_3$. The planted-forest
graphs (two disjoint trees planted on a genus-$1$ surface)
contribute a quadratic expression in $S_3$. The linear
correction $-\kappa\, S_3/48$ comes from the interaction
between the planted tree and the genus-$1$ curvature
$\kappa$. Explicit computation yields the stated formula.
\end{proof}


% ================================================================
% 10. EXPLICIT SHADOW TOWERS
% ================================================================

\section{Explicit shadow towers for the standard
landscape}\label{sec:landscape}

We compute the shadow data
$(\kappa, \alpha, S_4, \Delta, \rho)$ for all standard
families, constructing each tower from the Heisenberg base
case upward.

\subsection{Class $\mathbf{G}$: Heisenberg}
\label{subsec:heisenberg}

\begin{proposition}[Heisenberg shadow
tower]\label{prop:heisenberg}
For the rank-$1$ Heisenberg vertex algebra $\cH_k$ at
level $k \ne 0$:
\[
 \kappa \;=\; k,
 \qquad
 \alpha \;=\; 0,
 \qquad
 S_r \;=\; 0 \quad\text{for all } r \ge 3.
\]
The tower terminates at degree~$2$:
$Q_L(t) = 4k^2$ (constant), $H(t) = 2k\, t^2$.
Class~$\mathbf{G}$, $d_{\mathrm{alg}} = 0$.
The genus-$g$ free energy is
$F_g(\cH_k) = k \cdot \lambda_g^{\mathrm{FP}}$
(UNIFORM-WEIGHT) for all $g \ge 1$.
\end{proposition}

\begin{proof}
The Heisenberg OPE $J(z)J(w) \sim k/(z-w)^2$ is purely
quadratic: all collision residues on $\FM_n(X)$ for
$n \ge 3$ vanish. Hence $\alpha = 0$, $S_4 = 0$,
$\Delta = 0$. The shadow metric is the constant $4k^2$,
a perfect square.
\end{proof}

\begin{proposition}[Lattice VOA shadow
tower]\label{prop:lattice}
For the lattice vertex algebra $V_\Lambda$ associated to
an even positive-definite lattice $\Lambda$ of rank~$r$:
\[
 \kappa \;=\; r,
 \qquad
 \alpha \;=\; 0,
 \qquad
 S_n \;=\; 0 \quad\text{for all } n \ge 3.
\]
Class~$\mathbf{G}$, independent of the lattice cocycle.
\end{proposition}

\begin{proof}
The primary line of a lattice VOA is spanned by the
Heisenberg currents $J^i$, $i = 1, \ldots, r$.
The OPE $J^i(z)J^j(w) \sim \delta^{ij}/(z-w)^2$
is abelian: no cubic or higher terms on the primary line.
The lattice vertex operators $e^\alpha$ have nontrivial
OPEs but live on charged strata; they do not contribute to
the primary-line shadow tower. The modular characteristic
$\kappa = r$ equals the rank, independent of the specific
lattice structure (self-dual, Niemeier, Leech, etc.).
\end{proof}

\begin{remark}[Free fermion]\label{rem:free-fermion}
The free fermion $\psi$ with
$\psi(z)\psi(w) \sim 1/(z-w)$ has $c = 1/2$ and
$\kappa = c/2 = 1/4$. On the single-generator primary
line, the OPE is purely quadratic (Clifford), so
$\alpha = 0$, $S_4 = 0$, $\Delta = 0$.
Class~$\mathbf{G}$, $r_{\max} = 2$, $d_{\mathrm{alg}} = 0$.
\end{remark}

\begin{remark}[Class $\mathbf{G}$ as the
base case]\label{rem:class-G-base}
Class~$\mathbf{G}$ is the base case of the shadow
obstruction tower: the tower terminates at the curvature
$\kappa$, and all higher invariants vanish. This is the
vertex-algebraic analogue of the classical Koszul property:
the bar resolution is minimal, the $A_\infty$ structure
is formal (all $m_k = 0$ for $k \ge 3$), and the genus-$g$
free energies are determined by a single invariant. The
entire moduli-space content of a class~$\mathbf{G}$ algebra
is encoded in the number $\kappa$.

The transition from class~$\mathbf{G}$ to class~$\mathbf{L}$
is the introduction of a Lie bracket: the cubic shadow
$\alpha \ne 0$ measures the failure of the OPE to be
abelian. The transition from class~$\mathbf{L}$ to
class~$\mathbf{M}$ is the introduction of a nonlinear
contact term: the quartic shadow $S_4 \ne 0$ measures the
failure of the Jacobi identity to kill the four-point
function. Each transition introduces new invariants and
new complexity.
\end{remark}


\subsection{Class $\mathbf{L}$: affine
Kac--Moody}\label{subsec:km}

\begin{proposition}[Affine Kac--Moody shadow
tower]\label{prop:km}
For $V_k(\fg)$ at level $k \ne -h^\vee$ with $\fg$
simple:
\[
 \kappa
 \;=\;
 \frac{\dim(\fg)\,(k + h^\vee)}{2h^\vee},
 \qquad
 \alpha \;\ne\; 0,
 \qquad
 S_4 \;=\; 0,
 \qquad
 \Delta \;=\; 0.
\]
The cubic shadow is nonzero (generated by the Lie bracket),
and the quartic shadow vanishes on the primary line (killed
by the Jacobi identity). The tower terminates at
degree~$3$: $r_{\max} = 3$, class~$\mathbf{L}$,
$d_{\mathrm{alg}} = 1$.

Boundary values: at $k = 0$ (abelian limit),
$\kappa = \dim(\fg)/2 \ne 0$ (the algebra is still Koszul);
at $k = -h^\vee$ (critical level), $\kappa = 0$
(Koszulness fails, the center jumps).
\end{proposition}

\begin{proof}
The modular characteristic
$\kappa = \dim(\fg)(k + h^\vee)/(2h^\vee)$ is the genus-$1$
curvature scalar. The derivation: the OPE
$J^a(z)J^b(w) \sim k\,\delta^{ab}/(z-w)^2
+ f^{abc} J^c(w)/(z-w)$ has a double pole $k\delta^{ab}$ and
a simple pole $f^{abc} J^c$. The $R$-matrix in
trace-form convention is $r(z) = k\Omega/z$ (with level
prefix $k$; at $k = 0$, $r(z) = 0$). Averaging:
$\mathrm{av}(k\Omega/z) = k\,\dim(\fg)/(2h^\vee)$.
The Sugawara shift $\dim(\fg)/2$ arises from the
simple-pole self-contraction through the adjoint Casimir
eigenvalue $2h^\vee$. Total:
$\kappa = k\,\dim(\fg)/(2h^\vee) + \dim(\fg)/2
= \dim(\fg)(k + h^\vee)/(2h^\vee)$.

The cubic shadow $\alpha$ encodes the Lie bracket: the
three-point collision residue on $\FM_3(X)$ extracts
$f^{abc}$ from the simple-pole OPE. On the primary
line, $\alpha$ is a nonzero scalar proportional to the
cubic Casimir.

The quartic collision residue on $\FM_4(X)$, restricted
to the primary line, is proportional to the cyclic sum
$\sum_{\mathrm{cyclic}(a,b,c)} f^{abx} f^{xcd}$, which
vanishes by the Jacobi identity
$[[\cdot,\cdot],\cdot] + \mathrm{cyclic} = 0$.
Hence $S_4 = 0$, $\Delta = 8\kappa \cdot 0 = 0$, and
$Q_L = (2\kappa + 3\alpha\, t)^2$.

Verification at boundary values: at $k = 0$ (abelian
limit), $\kappa = \dim(\fg)/2 \ne 0$ for non-abelian $\fg$.
The $R$-matrix vanishes but the Sugawara shift persists;
the algebra remains Koszul. At $k = -h^\vee$ (critical
level), $\kappa = 0$: the modular characteristic vanishes,
the center of the completed enveloping algebra jumps (from
$\mathbb{C}$ to $\mathrm{Fun}(\mathrm{Op}_\fg)$), and
Koszulness fails.
\end{proof}

\begin{remark}[The Jacobi identity as depth
controller]\label{rem:jacobi-depth}
The vanishing $S_4 = 0$ for affine Kac--Moody is not a
numerical coincidence: it is a structural consequence of the
Jacobi identity. The class~$\mathbf{L}$ designation
reflects the Lie-theoretic origin of the algebra: the tower
terminates because the Lie bracket closes (Jacobi), not
because of any numerical cancellation in the OPE
coefficients. This explains why the transition from
class~$\mathbf{L}$ to class~$\mathbf{M}$ occurs precisely
when the Jacobi identity is replaced by a more general
relation (the Virasoro commutator, the $\cW_N$ OPE).
\end{remark}

The modular characteristics for small Lie algebras at
$k = 1$:

\begin{center}
\renewcommand{\arraystretch}{1.2}
\begin{tabular}{lcccc}
\toprule
$\fg$ & $\dim(\fg)$ & $h^\vee$
 & $\kappa(k)$ & $\kappa(1)$ \\
\midrule
$\mathfrak{sl}_2$ & $3$ & $2$
 & $3(k+2)/4$ & $9/4$ \\[2pt]
$\mathfrak{sl}_3$ & $8$ & $3$
 & $4(k+3)/3$ & $16/3$ \\[2pt]
$\mathfrak{sl}_N$ & $N^2{-}1$ & $N$
 & $(N^2{-}1)(k{+}N)/(2N)$ & --- \\[2pt]
$E_8$ & $248$ & $30$
 & $248(k{+}30)/60$ & $1922/15$ \\
\bottomrule
\end{tabular}
\end{center}

\noindent The Koszul dual
$V_k(\fg)^! = \mathrm{CE}^{\mathrm{ch}}(\fg_{-k-2h^\vee})$
has $\kappa^! = -\kappa$. The Koszul conductor is
family-specific:
$\kappa(V_k(\fg)) + \kappa(V_k(\fg)^!) = 0$ for affine
Kac--Moody and free fields;
$\kappa(\Vir_c) + \kappa(\Vir_{26-c}) = 13$ for Virasoro.
These are distinct: the vanishing for Kac--Moody reflects
the anti-symmetry of the chiral Chevalley--Eilenberg
construction.


\subsection{Class $\mathbf{C}$:
$\beta\gamma$}\label{subsec:betagamma}

\begin{proposition}[$\beta\gamma$ shadow
tower]\label{prop:betagamma}
For the $\beta\gamma$ system of conformal weight $\lambda$:
\begin{alignat*}{2}
 &\kappa \;=\; 6\lambda^2 - 6\lambda + 1
 \;=\; c/2,
 \qquad
 &&c \;=\; 2(6\lambda^2 - 6\lambda + 1), \\
 &\alpha \;=\; 0
 &&\textup{(weight-changing line)}, \\
 &S_4 \;=\; 0
 &&\textup{(weight-changing line)}, \\
 &\Delta \;=\; 0
 &&\textup{(weight-changing line)}.
\end{alignat*}
On the weight-changing primary line, the tower terminates
at degree~$2$. The overall classification is
class~$\mathbf{C}$: a nonzero quartic contact invariant
lives on a charged stratum, and the quintic vanishes by
rank-one rigidity. Hence $r_{\max} = 4$,
$d_{\mathrm{alg}} = 2$.

Boundary values: $\lambda = 1/2$ gives
$c = -1$ (symplectic boson), $\kappa = -1/2$;
$\lambda = 2$ gives $c = 26$ (matter ghost),
$\kappa = 13$.
\end{proposition}

\begin{proof}
The $\beta\gamma$ system has fields $\beta$ of weight
$\lambda$ and $\gamma$ of weight $1 - \lambda$, with OPE
$\beta(z)\gamma(w) \sim 1/(z-w)$. The central charge is
$c_{\beta\gamma}(\lambda) = 2(6\lambda^2 - 6\lambda + 1)$,
and $\kappa = c/2 = 6\lambda^2 - 6\lambda + 1$.

On the weight-changing primary line (spanned by $\beta\gamma$
composites of definite weight), the OPE is abelian: the
Wick contraction of $\beta$ with $\gamma$ produces only the
identity, with no higher poles. Hence $\alpha = 0$ and
$S_4 = 0$ on this line, and $\Delta = 0$.

The contact stratum: the full two-variable deformation
complex has a charged stratum supporting the normal-ordered
composite $:\!\beta\partial\gamma\!:$ and its cousins. On this
stratum, the four-point function receives a nonzero quartic
contact contribution from the Wick contraction
$\langle\beta\gamma\beta\gamma\rangle_{\mathrm{connected}}$
$= 1/(z_1 - z_3)(z_2 - z_4)
- 1/(z_1 - z_4)(z_2 - z_3)$,
whose collision residue on $\FM_4(X)$ is nonzero. The
quartic contact invariant is
$Q^{\mathrm{cont}} = -5/12$ (after normalization).

The quintic obstruction vanishes by rank-one rigidity: the
self-bracket of the quartic contact invariant under the
$H$-Poisson structure produces a vector in the
two-dimensional weight/contact slice, but this vector is
orthogonal to the primary line and thus exits the
single-line complex. On a one-dimensional target, the
image is zero. Hence $o_5 = 0$ and the tower terminates at
degree~$4$.

Boundary values: $c_{\beta\gamma}(1/2) = -1$ (symplectic
boson), $\kappa = -1/2$; $c_{\beta\gamma}(2) = 26$ (matter
ghost), $\kappa = 13$; $c_{\beta\gamma}(0) = 2$ (standard
$\beta\gamma$), $\kappa = 1$;
$c_{\beta\gamma}(1) = 2$ (adjoint $\beta\gamma$),
$\kappa = 1$.
\end{proof}

\begin{remark}[$bc$ ghosts]\label{rem:bc-ghosts}
The fermionic $bc$ system has
$c_{bc}(\lambda) = 1 - 3(2\lambda-1)^2$ and
$\kappa_{bc} = c_{bc}/2 = -(6\lambda^2 - 6\lambda + 1)$.
The complementarity
$c_{\beta\gamma}(\lambda) + c_{bc}(\lambda) = 0$
holds at each conformal weight:
$2(6\lambda^2 - 6\lambda + 1) + 1 - 3(2\lambda-1)^2 = 0$
(direct expansion). At $\lambda = 2$
(reparametrization ghosts): $c_{bc} = -26$,
$\kappa_{bc} = -13$. At $\lambda = 1/2$ (single Dirac
fermion): $c_{bc} = 1$, $\kappa_{bc} = 1/2$.
Class~$\mathbf{C}$ for all~$\lambda$.
\end{remark}

\begin{remark}[The contact stratum
geometry]\label{rem:contact-geometry}
Class~$\mathbf{C}$ is the unique class where the shadow
tower terminates at a finite depth greater than~$3$. The
mechanism is geometric: the quartic contact invariant lives
on a stratum of the cyclic deformation complex that is
\emph{transverse} to every one-dimensional primary line.
On any primary line, the shadow data is class~$\mathbf{G}$
(no cubic, no quartic); the class~$\mathbf{C}$ designation
comes from the global structure of the deformation complex.

This transversality is specific to the $\beta\gamma/bc$
family: the weight-changing direction (which supports the
primary-line tower) is orthogonal to the ghost-number
direction (which supports the quartic contact). For Virasoro
and $\cW_N$, the quartic contribution lies on the primary
line itself, producing $S_4 \ne 0$ and $\Delta \ne 0$
(class~$\mathbf{M}$). For Heisenberg and affine KM, the
quartic contribution vanishes entirely (abelian or Jacobi).
\end{remark}

\begin{example}[$\beta\gamma$ at special
weights]\label{ex:betagamma-special}
\leavevmode
\begin{center}
\renewcommand{\arraystretch}{1.2}
\begin{tabular}{ccccc}
\toprule
$\lambda$ & $c_{\beta\gamma}$ & $\kappa$
 & Interpretation & $c_{bc}$ \\
\midrule
$0$ & $2$ & $1$ & standard & $-2$ \\
$1/2$ & $-1$ & $-1/2$ & symplectic boson & $1$ \\
$1$ & $2$ & $1$ & adjoint & $-2$ \\
$3/2$ & $10$ & $5$ & gravitino & $-10$ \\
$2$ & $26$ & $13$ & matter ghost & $-26$ \\
\bottomrule
\end{tabular}
\end{center}
At $\lambda = 2$, the matter ghost has $c = 26$, which
is the string-theoretic value: the $\beta\gamma$ matter
ghost cancels the $bc$ reparametrization ghost central
charge ($26 + (-26) = 0$). The modular characteristics
also cancel: $13 + (-13) = 0$.
\end{example}


\subsection{Class $\mathbf{M}$: Virasoro}
\label{subsec:virasoro}

\begin{proposition}[Virasoro shadow
tower]\label{prop:virasoro}
For $\Vir_c$ at $c \ne 0$, $c \ne -22/5$:
\[
 \kappa \;=\; \frac{c}{2},
 \qquad
 \alpha \;=\; 2,
 \qquad
 S_4 \;=\; \frac{10}{c(5c+22)},
 \qquad
 \Delta \;=\; \frac{40}{5c+22}.
\]
The shadow metric is
\[
 Q_L(t)
 \;=\;
 (c + 6t)^2 + \frac{80\, t^2}{5c + 22}.
\]
The tower is infinite (class~$\mathbf{M}$,
$d_{\mathrm{alg}} = \infty$), and the shadow growth rate is
$\rho = \frac{1}{c}\sqrt{(180c + 872)/(5c+22)}$.
\end{proposition}

\begin{remark}[The quartic contact
invariant]\label{rem:quartic-contact}
The quartic shadow
$Q^{\mathrm{contact}}_{\Vir} = S_4 = 10/[c(5c+22)]$ is
the first genuinely nonlinear OPE invariant. It measures
the irreducible four-point collision residue on $\FM_4(X)$
and is the coefficient that separates class~$\mathbf{M}$
from class~$\mathbf{L}$. Its poles at $c = 0$ and
$c = -22/5$ reflect the failure of the Virasoro OPE at
these special values. The large-$c$ asymptotic is
$S_4 \sim 2/c^2$.
\end{remark}

\begin{theorem}[Virasoro shadow
coefficients]\label{thm:virasoro-coefficients}
The shadow coefficients through degree~$10$:
\begin{align*}
 S_2 &\;=\; c/2, \\[3pt]
 S_3 &\;=\; 2, \\[3pt]
 S_4 &\;=\; \frac{10}{c(5c+22)}, \\[3pt]
 S_5 &\;=\; \frac{-48}{c^2(5c+22)}, \\[3pt]
 S_6 &\;=\; \frac{80(45c+193)}{3\,c^3(5c+22)^2}, \\[3pt]
 S_7 &\;=\;
 \frac{-2880(15c+61)}{7\,c^4(5c+22)^2}, \\[3pt]
 S_8 &\;=\;
 \frac{80(2025c^2+16470c+33314)}{c^5(5c+22)^3}, \\[3pt]
 S_9 &\;=\;
 \frac{-1280(2025c^2+15570c+29554)}
 {3\,c^6(5c+22)^3}, \\[3pt]
 S_{10} &\;=\;
 \frac{256(91125c^3+1050975c^2+3989790c+4969967)}
 {c^7(5c+22)^4}.
\end{align*}
All coefficients belong to $\mathbb{Q}(c)$ with poles
only at $c = 0$ and $c = -22/5$. Signs alternate from
$S_5$ onward: $S_5 < 0$, $S_6 > 0$, $S_7 < 0$, \ldots\
for $c > 0$.
The $c$-independence of $S_3 = 2$ is a theorem; see
\cite{LorgatVirR}.
\end{theorem}

\begin{proof}
The computation proceeds inductively using the master
equation (Proposition~\ref{prop:master-eq}) on the primary
line~$L$ spanned by~$T$. The Hessian propagator is
$P = 1/\kappa = 2/c$.

\emph{Degree $2$.}\enspace
$S_2 = \kappa = c/2$ by definition: this is the genus-$1$
curvature scalar $F_1 = \kappa/24 = c/48$.

\emph{Degree $3$.}\enspace
The OPE $T_{(1)}T = 2T$ gives the simple-pole
collision residue (after $d\!\log$ absorption from the
$T_{(2)}T = \partial T$ term). On the primary line:
$S_3 = \alpha = 2$, independent of~$c$. This
$c$-independence follows from the universality of the
Virasoro commutation relation
$[L_m, L_n] = (m-n)L_{m+n} + (c/12)(m^3 - m)\delta_{m+n,0}$:
the coefficient of $L_{m+n}$ is $m - n$, not $c$-dependent.
The scalar $\alpha = 2$ is twice the coefficient of $T$
in $T_{(1)}T = 2T$.

\emph{Degree $4$.}\enspace
The master equation at degree~$4$ is
$d_0(\Sh_4) + \frac{1}{2}[\Sh_3, \Sh_3] = 0$.
On the primary line:
$8\kappa S_4 + \frac{1}{2}\, c_{33}\, 9\, P\, \alpha^2 = 0$,
where $c_{33} = 1/2$ (self-pairing correction). Solving:
\[
 S_4
 \;=\;
 -\frac{1}{8\kappa} \cdot \frac{1}{2} \cdot 9 \cdot P
 \cdot \alpha^2
 \;=\;
 -\frac{9\alpha^2}{16\kappa^2} + S_4^{\mathrm{contact}},
\]
where $S_4^{\mathrm{contact}}$ is the irreducible quartic
contact term from the four-point OPE. For Virasoro, the
quartic OPE coefficient is
$T_{(3)}T = c/2$, and the contact extraction gives
$S_4^{\mathrm{contact}} = 10/[c(5c+22)]$ after the
propagator renormalization.

The full result: $S_4 = 10/[c(5c+22)]$, and the critical
discriminant is
$\Delta = 8\kappa S_4 = 8(c/2) \cdot 10/[c(5c+22)]
= 40/(5c+22) \ne 0$ for $c \ne -22/5$.

\emph{Degree $5$.}\enspace
$d_0(\Sh_5) + [\Sh_3, \Sh_4] = 0$, giving
$S_5 = -(P/10) \cdot 3 \cdot 4 \cdot S_3 S_4
= -(6/5)(2/c) \cdot 2 \cdot 10/[c(5c+22)]
= -48/[c^2(5c+22)]$.

\emph{Higher degrees.}\enspace
The pattern continues: each $S_r$ involves a sum over pairs
$(j,k)$ with $j + k = r + 2$, weighted by the propagator
$P = 2/c$ and the symmetry factors $c_{jk}$.
Alternatively, the closed-form extraction
$S_r = \frac{1}{r}[t^{r-2}]\sqrt{Q_{\Vir}(t)}$
with $Q_{\Vir}(t) = (c + 6t)^2 + 80t^2/(5c+22)$
gives all coefficients by binomial expansion.

Both methods have been verified computationally through
degree~$20$; the results agree to all digits.
\end{proof}

\begin{remark}[Pole structure of Virasoro
shadows]\label{rem:pole-structure}
Every $S_r$ with $r \ge 4$ has poles only at $c = 0$ and
$c = -22/5$. The pole at $c = 0$ has order $r - 3$; the
pole at $c = -22/5$ has order $\lfloor(r-2)/2\rfloor$.
This pole structure is inherited from the branch points of
$\sqrt{Q_{\Vir}(t)}$: the zero of $Q_{\Vir}$ nearest to
$t = 0$ has modulus inversely proportional to the shadow
growth rate $\rho$, and the Taylor coefficients of a
square-root function at a simple branch point have a
universal pole pattern determined by the Flajolet--Sedgewick
transfer theorem.

The physical content: the pole at $c = 0$ reflects the
vanishing of the Virasoro central charge (the algebra
degenerates to the Witt algebra, which has no central
extension and no modular characteristic). The pole at
$c = -22/5$ is the second zero of the discriminant of the
Virasoro Verma module at level~$2$: at this value, the
four-point contact term diverges, and the shadow metric
$Q_{\Vir}$ degenerates.
\end{remark}

Numerical evaluations:

\begin{center}
\renewcommand{\arraystretch}{1.2}
\begin{tabular}{rccc}
\toprule
$r$ & $S_r(c{=}1)$ & $S_r(c{=}13)$ & $S_r(c{=}26)$ \\
\midrule
$2$ & $1/2$ & $13/2$ & $13$ \\
$3$ & $2$ & $2$ & $2$ \\
$4$ & $10/27$ & $10/1131$ & $5/1976$ \\
$5$ & $-16/9$ & $-16/4901$ & $-3/6422$ \\
\bottomrule
\end{tabular}
\end{center}

\noindent At $c = 26$ the coefficients decay rapidly
($\rho \approx 0.232$); at $c = 13$ moderately
($\rho \approx 0.467$); at $c = 1$ they grow
exponentially ($\rho \approx 6.242$). All unitary
minimal models ($c < 1$) have strongly divergent towers.

Extended numerical data through degree~$8$:

\begin{center}
\renewcommand{\arraystretch}{1.2}
\begin{tabular}{rccc}
\toprule
$r$ & $S_r(c{=}1)$ & $S_r(c{=}13)$ & $S_r(c{=}26)$ \\
\midrule
$6$ & $19040/2187$ & $62240/49887279$
 & $6815/76139232$ \\
$7$ & $-24320/567$ & $-81920/168138607$
 & $-20295/1154778352$ \\
$8$ & $4144720/19683$ & $47171920/244497554379$
 & $4576085/1303909727744$ \\
\bottomrule
\end{tabular}
\end{center}

\noindent The alternating signs $(-1)^{r+1}$ for $r \ge 5$
(at $c > 0$) are a consequence of the complex conjugate
branch points of $Q_L$: the transfer theorem gives
oscillating coefficients with frequency determined by the
argument of the nearest singularity.

\begin{remark}[The Virasoro shadow at special
values]\label{rem:virasoro-special}
Three special values of~$c$ merit attention:

\emph{$c = 13$ (self-dual).}\enspace
Under Koszul duality $c \mapsto 26 - c$,
the self-dual point $c = 13$ maps to itself. The shadow
data is symmetric: $\kappa = 13/2$, $\kappa' = 13/2$,
$\alpha = \alpha' = 2$, $S_4(13) = S_4(13)$.
The discriminant $\Delta(13) = 40/87$, and the
shadow growth rate $\rho(13) \approx 0.467$.
All shadow coefficients are self-dual:
$S_r(\Vir_{13}) = S_r(\Vir_{13}^!)$.

\emph{$c = 26$ (string).}\enspace
The Koszul dual $\Vir_0$ has $\kappa = 0$ (the bar complex
is uncurved). The shadow tower of $\Vir_{26}$ is
well-behaved ($\rho \approx 0.232$), but the dual tower is
degenerate. The discriminant $\Delta(26) = 40/152 = 5/19$.

\emph{$c = -22/5$ (pole).}\enspace
This is the location of the denominator singularity
$5c + 22 = 0$. The quartic shadow $S_4$ diverges,
$\Delta \to \infty$, and the shadow metric degenerates.
This value is associated with the $(2,5)$ minimal model
in the non-unitary series. The shadow tower is undefined at
$c = -22/5$ but has well-defined limits from both sides.
\end{remark}

\begin{proposition}[The Virasoro shadow connection:
explicit formula]\label{prop:virasoro-conn}
The shadow connection for Virasoro is
\begin{equation}\label{eq:virasoro-conn}
 \nabla^{\mathrm{sh}}
 \;=\;
 d \;-\;
 \frac{6c + (18 + 80/(5c+22))\, t}
 {(c + 6t)^2 + 80t^2/(5c+22)}\, dt.
\end{equation}
The connection has two regular singular points at the zeros
of $Q_{\Vir}(t)$, each with residue~$1/2$. The zeros are
complex conjugate for $c > 0$:
\[
 t_\pm
 \;=\;
 \frac{-3c(5c+22) \pm ic\sqrt{80(5c+22)}}
 {(9(5c+22) + 80)}.
\]
The branch-point modulus is
$|t_\pm| = c/\sqrt{9 + 80/(5c+22)}
= 2\kappa/\sqrt{9\alpha^2 + 2\Delta}
= 1/\rho$, confirming the shadow growth rate formula.
\end{proposition}

\begin{remark}[Duality and the shadow
connection]\label{rem:duality-conn}
Under $c \mapsto 26 - c$, the connection
form transforms by a rational substitution:
\[
 \frac{Q_L'(t)}{2Q_L(t)}\bigg|_{c}
 \;\ne\;
 \frac{Q_L'(t)}{2Q_L(t)}\bigg|_{26-c}.
\]
The two connections are related by the Verdier involution
on the shadow data, but they are not isomorphic as
connections on the same base: the branch points move under
$c \mapsto 26 - c$. The discriminant complementarity
$\Delta(c) + \Delta(26-c) = 6960/[(5c+22)(152-5c)]$
(Proposition~\ref{prop:complementarity}) reflects this
transformation.
\end{remark}


\subsection{The shadow depth decomposition}
\label{subsec:depth-decomp}

For class~$\mathbf{M}$ algebras, the shadow tower is
infinite. The complexity of the tower admits a decomposition
into algebraic and arithmetic components.

\begin{definition}[Shadow depth
decomposition]\label{def:depth-decomp}
For a chirally Koszul vertex algebra~$\cA$ of
class~$\mathbf{M}$, define the \emph{total shadow depth}
\[
 d(\cA) \;=\; 1 + d_{\mathrm{arith}} + d_{\mathrm{alg}},
\]
where:
\begin{itemize}
\item $d_{\mathrm{alg}} \in \{0, 1, 2, \infty\}$ is the
 algebraic depth (Definition~\ref{def:algebraic-depth}).
 For class~$\mathbf{M}$: $d_{\mathrm{alg}} = \infty$
 (all $S_r$ nonzero), but the algebraicity theorem shows
 that three invariants determine the full tower, so the
 ``effective algebraic complexity'' is~$0$.
\item $d_{\mathrm{arith}} \ge 0$ is the \emph{arithmetic
 depth}: the number of genuinely independent arithmetic
 invariants beyond the algebraic shadow data. Shadow
 coefficients of weight $\le 2r$ at degree~$r$ lie in the
 subspace spanned by Eisenstein series; cusp-form
 contributions first appear at genus $g \ge 6$ (the
 Ramanujan $\Delta$ form at weight~$12$).
\end{itemize}
\end{definition}

\begin{remark}[Algebraic vs arithmetic
invariants]\label{rem:alg-arith}
The algebraic depth is determined by genus-$0$ OPE data
(Theorem~\ref{thm:riccati}). The arithmetic depth reflects
automorphic phenomena in the genus expansion: cusp-form
contributions are invisible at low genus but may enter at
$g \ge 6$. For all families treated in
Section~\ref{sec:landscape}, $d_{\mathrm{arith}} = 0$
through the genera computed here.

The distinction between algebraic and arithmetic depth is
the shadow-tower analogue of the distinction between
perturbative and non-perturbative contributions in quantum
field theory: the algebraic part is determined by finitely
many OPE data, while the arithmetic part involves new
transcendental numbers (cusp-form coefficients) that are
not determined by the OPE.
\end{remark}


\subsection{Shadow growth rate}\label{subsec:growth-rate}

\begin{definition}[Shadow growth
rate]\label{def:shadow-radius}
For class~$\mathbf{M}$ with $\Delta \ne 0$:
\[
 \rho(\cA)
 \;:=\;
 \frac{\sqrt{9\alpha^2 + 2\Delta}}{2|\kappa|}.
\]
\end{definition}

\begin{proposition}\label{prop:rho-convergence}
The growth rate $\rho(\cA)$ is the reciprocal of the
convergence radius of $H(t) = t^2\sqrt{Q_L(t)}$. The
branch points of $H$ are the zeros of $Q_L$. For
$\Delta > 0$, these form a complex conjugate pair of
modulus $R = 2|\kappa|/\sqrt{9\alpha^2 + 2\Delta}$, and
\[
 S_r
 \;\sim\;
 C(\cA)\,\rho(\cA)^r\, r^{-5/2}\,
 \cos(r\theta + \phi),
 \qquad r \to \infty,
\]
where the exponent $-5/2$ comes from the square-root
singularity (Flajolet--Sedgewick transfer theorem).
\end{proposition}

\begin{example}[Virasoro shadow growth
rate]\label{ex:virasoro-rho}
For $\Vir_c$ with $c > 0$:
\[
 \rho(\Vir_c)
 \;=\;
 \frac{1}{c}
 \sqrt{\frac{180c + 872}{5c + 22}}.
\]
The critical central charge $c^* \approx 6.125$ (the
unique positive root of $5c^3 + 22c^2 - 180c - 872 = 0$)
separates the convergent regime ($c > c^*$, $\rho < 1$)
from the divergent regime ($c < c^*$, $\rho > 1$).

Koszul duality exchanges convergent and divergent towers:
\begin{center}
\renewcommand{\arraystretch}{1.2}
\begin{tabular}{ccccl}
\toprule
$c$ & $\rho(c)$ & $\rho(26{-}c)$ & converges? &
 interpretation \\
\midrule
$1/2$ & $12.53$ & $0.224$ & no/yes & Ising \\
$1$ & $6.242$ & $0.242$ & no/yes & free boson \\
$c^*$ & $1$ & $0.430$ & marginal/yes & critical \\
$13$ & $0.467$ & $0.467$ & yes/yes & self-dual \\
$26$ & $0.232$ & $\infty$ & yes/--- &
 string; $\kappa^! = 0$ \\
\bottomrule
\end{tabular}
\end{center}
\end{example}


\subsection{Class $\mathbf{M}$:
$\cW_N$}\label{subsec:wn}

\begin{proposition}[$\cW_N$ shadow
data]\label{prop:wn}
For $\cW_N$ ($N \ge 3$) at central charge~$c$, the
total modular characteristic is
\[
 \kappa(\cW_N)
 \;=\;
 c\,(H_N - 1),
 \qquad
 H_N \;=\; \sum_{j=1}^{N} \frac{1}{j}.
\]
On the $T$-line (Virasoro direction):
$\kappa_T = c/2$, $\alpha_T = 2$,
$S_4 = 10/[c(5c+22)]$, class~$\mathbf{M}$.
\end{proposition}

The anomaly ratios for small~$N$:

\begin{center}
\renewcommand{\arraystretch}{1.2}
\begin{tabular}{cccc}
\toprule
$N$ & $\cW_N$ & $H_N - 1$ & $\kappa(\cW_N)$ \\
\midrule
$2$ & Virasoro & $1/2$ & $c/2$ \\
$3$ & $\cW_3$ & $5/6$ & $5c/6$ \\
$4$ & $\cW_4$ & $13/12$ & $13c/12$ \\
$5$ & $\cW_5$ & $77/60$ & $77c/60$ \\
$6$ & $\cW_6$ & $29/20$ & $29c/20$ \\
\bottomrule
\end{tabular}
\end{center}

\noindent Sanity check at $N = 2$:
$H_2 - 1 = 3/2 - 1 = 1/2$, so
$\kappa(\cW_2) = c/2 = \kappa(\Vir_c)$.
At $N = 3$: $H_3 - 1 = 11/6 - 1 = 5/6$, so
$\kappa(\cW_3) = 5c/6$.


\subsection{Master table}\label{subsec:master-table}

\begin{center}
\renewcommand{\arraystretch}{1.4}
\small
\begin{tabular}{lccccccc}
\toprule
Family & $\kappa$ & $\alpha$ & $S_4$ & $\Delta$
 & Class & $r_{\max}$ & $\rho$ \\
\midrule
$\cH_k$
 & $k$ & $0$ & $0$ & $0$
 & $\mathbf{G}$ & $2$ & $0$ \\[2pt]
$V_\Lambda$
 & $\rk(\Lambda)$ & $0$ & $0$ & $0$
 & $\mathbf{G}$ & $2$ & $0$ \\[2pt]
\addlinespace
$V_k(\fg)$
 & $\frac{\dim\fg\,(k{+}h^\vee)}{2h^\vee}$
 & $\ne 0$ & $0$ & $0$
 & $\mathbf{L}$ & $3$ & $0$ \\[2pt]
\addlinespace
$\beta\gamma_\lambda$
 & $6\lambda^2{-}6\lambda{+}1$
 & $0^*$ & $0^*$ & $0^*$
 & $\mathbf{C}$ & $4$ & --- \\[2pt]
\addlinespace
$\Vir_c$
 & $c/2$ & $2$
 & $\frac{10}{c(5c{+}22)}$
 & $\frac{40}{5c{+}22}$
 & $\mathbf{M}$ & $\infty$
 & $\frac{1}{c}\!\sqrt{\frac{180c{+}872}{5c{+}22}}$
 \\[4pt]
$\cW_N$
 & $(H_N{-}1)\,c$
 & $\ne 0$ & $\ne 0$ & $\ne 0$
 & $\mathbf{M}$ & $\infty$ & --- \\
\bottomrule
\end{tabular}
\end{center}


% ================================================================
% 10b. DISCRIMINANT ANALYSIS AND KOSZUL DUALITY
% ================================================================

\subsection{The discriminant across the
landscape}\label{subsec:disc-landscape}

The critical discriminant $\Delta = 8\kappa S_4$ is the
single number that separates finite towers from infinite
ones. Its values across the standard landscape exhibit a
sharp dichotomy:

\begin{center}
\renewcommand{\arraystretch}{1.3}
\begin{tabular}{lccc}
\toprule
Family & $\Delta$ & $\Delta = 0$? & Class \\
\midrule
$\cH_k$ & $0$ & yes & $\mathbf{G}$ \\
$V_\Lambda$ & $0$ & yes & $\mathbf{G}$ \\
Free fermion & $0$ & yes & $\mathbf{G}$ \\
$V_k(\fg)$ & $0$ & yes & $\mathbf{L}$ \\
$\beta\gamma_\lambda$ & $0^*$ & yes$^*$ & $\mathbf{C}$ \\
\addlinespace
$\Vir_c$ & $40/(5c+22)$ & no & $\mathbf{M}$ \\
$\cW_3$ & $\ne 0$ & no & $\mathbf{M}$ \\
$\cW_N$ ($N \ge 3$) & $\ne 0$ & no & $\mathbf{M}$ \\
\bottomrule
\end{tabular}
\end{center}

\noindent The vanishing $\Delta = 0$ for classes
$\mathbf{G}$, $\mathbf{L}$, and $\mathbf{C}$ has three
distinct mechanisms:
\begin{itemize}
\item Class $\mathbf{G}$: the OPE is abelian ($S_4 = 0$
 because there is no four-point connected correlator).
\item Class $\mathbf{L}$: the OPE is non-abelian but Lie
 ($S_4 = 0$ because the Jacobi identity kills the quartic
 collision residue).
\item Class $\mathbf{C}$: $S_4 = 0$ on every primary
 line, but the quartic contact invariant lives on a
 transverse stratum.
\end{itemize}

\begin{proposition}[The discriminant as a Koszul
invariant]\label{prop:disc-invariant}
The critical discriminant $\Delta$ is a Koszul invariant:
it transforms covariantly under Koszul duality
$\cA \mapsto \cA^!$. For Virasoro:
\[
 \Delta(\Vir_c)
 \;=\;
 \frac{40}{5c+22},
 \qquad
 \Delta(\Vir_{26-c})
 \;=\;
 \frac{40}{152-5c}.
\]
The discriminant never vanishes for $c \ne -22/5$; the
Koszul dual never changes class.
\end{proposition}

\begin{proof}
Direct substitution: $\Delta(c) = 8 \cdot (c/2) \cdot
10/[c(5c+22)] = 40/(5c+22)$ and
$\Delta(26-c) = 40/(5(26-c)+22) = 40/(152-5c)$. Both are
nonzero for $c \ne -22/5$ and $c \ne 152/5 = 30.4$,
respectively. The dual singularities $c = -22/5$ and
$c = 152/5$ sum to~$26$: they are Koszul partners.
\end{proof}

\begin{remark}[The discriminant
locus]\label{rem:disc-locus}
In the three-dimensional parameter space $(\kappa, \alpha,
S_4)$, the discriminant locus $\Delta = 0$ is the
hyperplane $\kappa S_4 = 0$, which decomposes as
$\kappa = 0$ (degenerate: the propagator is undefined)
and $S_4 = 0$ (the Jacobi locus). The complement
$\Delta \ne 0$ is a dense open subset.

Every standard family lies on the discriminant locus
except for Virasoro and $\cW_N$ ($N \ge 3$). The
$\beta\gamma$ family lies on the locus \emph{on every
primary line} but escapes it on the charged stratum. The
classification is:
\begin{itemize}
\item $\kappa = 0$: critical (not chirally Koszul in
 standard sense).
\item $S_4 = 0$, $\alpha = 0$: class $\mathbf{G}$.
\item $S_4 = 0$, $\alpha \ne 0$: class $\mathbf{L}$.
\item $S_4 \ne 0$ (on charged stratum only):
 class $\mathbf{C}$.
\item $S_4 \ne 0$ (on primary line): class $\mathbf{M}$.
\end{itemize}
This recovers the four-class partition from the discriminant
alone.
\end{remark}


\subsection{Complementarity across families}
\label{subsec:comp-families}

The Koszul conductor $K(\cA) = \kappa(\cA) + \kappa(\cA^!)$
is family-specific. The standard landscape yields four
distinct values:

\begin{center}
\renewcommand{\arraystretch}{1.3}
\begin{tabular}{lcc}
\toprule
Family & Koszul conductor $K$ & $K = 0$? \\
\midrule
$\cH_k$ & $0$ & yes \\
$V_k(\fg)$ & $0$ & yes \\
$\beta\gamma_\lambda$ & $0$ & yes \\
Lattice & $0$ & yes \\
\addlinespace
$\Vir_c$ & $13$ & no \\
$\cW_3$ & $250/3$ & no \\
Bershadsky--Polyakov & $98/3$ & no \\
\bottomrule
\end{tabular}
\end{center}

\noindent The vanishing $K = 0$ for abelian algebras and
affine Kac--Moody is the anti-symmetry property of the
chiral Chevalley--Eilenberg construction: the dual has
$\kappa^! = -\kappa$. For Virasoro, the conductor
$K = 13$ is the self-dual value: $\kappa(c) + \kappa(26-c)
= c/2 + (26-c)/2 = 13$.

\begin{remark}[The conductor is not the central
charge]\label{rem:conductor-not-c}
The Koszul conductor $K = \kappa + \kappa'$ is an invariant
of the Koszul \emph{pair} $(\cA, \cA^!)$, not of the
algebra alone. It should not be confused with the central
charge~$c$ (which is an invariant of the algebra) or with
the Koszul complementarity $c + c' = 26$ of the Virasoro
family (which is a statement about central charges, not
modular characteristics). The conductor is
$\kappa + \kappa'$, not $c + c'$.
\end{remark}


% ================================================================
% 11. GENUS-g FREE ENERGIES
% ================================================================

\section{Genus-$g$ free energies from the shadow
tower}\label{sec:genus-g}

The shadow obstruction tower produces invariants of moduli
spaces through graph sums over stable curves. The genus-$g$
\emph{free energy} $F_g(\cA)$ is the integral of the
tautological class produced by the MC equation on~$\Mbar_g$.

\subsection{The single-generator
formula}\label{subsec:single-gen}

\begin{theorem}[Single-generator genus-$g$
universality]\label{thm:single-gen-universality}
For any chirally Koszul vertex algebra~$\cA$ with a single
strong generator of conformal weight~$h$
\textup{(}UNIFORM-WEIGHT\textup{)}:
\begin{equation}\label{eq:single-gen-Fg}
 F_g(\cA) \;=\; \kappa(\cA) \cdot
 \lambda_g^{\mathrm{FP}}
\end{equation}
for all $g \ge 1$, where
$\lambda_g^{\mathrm{FP}} =
\int_{\Mbar_{g,1}} \psi_1^{2g-2}\, \lambda_g$
is the Faber--Pandharipande integral.
\end{theorem}

\begin{proof}[Proof sketch]
The genus-$g$ free energy is a sum over stable graphs
$\Gamma$ of genus~$g$ with no marked points:
\[
 F_g(\cA) \;=\;
 \sum_\Gamma
 \frac{1}{|\Aut(\Gamma)|}
 \prod_{e \in E(\Gamma)} P
 \prod_{v \in V(\Gamma)} V_{g_v, n_v},
\]
where $P = 1/\kappa$ is the propagator and
$V_{g_v, n_v}$ is the vertex factor at a vertex of
genus~$g_v$ with $n_v$ half-edges. For a single-generator
algebra, every edge carries the same channel. The vertex
factors are determined by the shadow coefficients:
$V_{0,3} = S_3 = \alpha$, $V_{0,4} = S_4$, and
$V_{g,n}$ for $g \ge 1$ encodes the genus-$g$ contribution.

The key observation is that for single-generator algebras,
the graph sum telescopes: the combinatorics of
automorphisms, vertex factors, and propagators combine to
produce $\kappa \cdot \lambda_g^{\mathrm{FP}}$, where the
Faber--Pandharipande integral absorbs all the graph-theoretic
data. The telescoping uses the dilaton equation and the
string equation on $\Mbar_{g,n}$ to reduce the sum to a
single integral.

At genus~$1$: $F_1 = \kappa/24 = \kappa \cdot
\lambda_1^{\mathrm{FP}}$ (unconditional, holds for all
families including multi-generator).
At genus~$2$: $F_2 = 7\kappa/5760 = \kappa \cdot
\lambda_2^{\mathrm{FP}}$ (requires single-generator
hypothesis).
\end{proof}

\begin{remark}[The generating
function]\label{rem:FP-generating}
The Faber--Pandharipande coefficients are the Taylor
coefficients of the $\hat{A}$-genus:
\[
 \sum_{g \ge 1} \lambda_g^{\mathrm{FP}}\, \hbar^{2g}
 \;=\;
 \frac{\hbar/2}{\sin(\hbar/2)} - 1.
\]
The first values are: $\lambda_1^{\mathrm{FP}} = 1/24$,
$\lambda_2^{\mathrm{FP}} = 7/5760$,
$\lambda_3^{\mathrm{FP}} = 31/967680$.
\end{remark}

\subsection{The genus-$2$ verification}
\label{subsec:genus2-verification}

The genus-$2$ formula
$F_2 = 7\kappa/5760$ requires an explicit sum over the
seven stable graphs. For a single-generator algebra with
shadow data $(\kappa, \alpha, S_4)$:

\begin{lemma}[Genus-$2$ single-generator
sum]\label{lem:genus2-single}
The graph-by-graph contributions to $F_2$ for a
single-generator chirally Koszul algebra are:
\begin{center}
\renewcommand{\arraystretch}{1.3}
\begin{tabular}{lcc}
\toprule
Graph & Contribution & Simplified \\
\midrule
$\Gamma_0$: smooth
 & $V_{2,0}$ & $7\kappa/5760$ \\
$\Gamma_1$: figure-eight
 & $P \cdot \kappa^2/(24^2 \cdot 2)$
 & $\kappa/1152$ \\
$\Gamma_2$: banana
 & $P^2 S_4/8$ & $S_4/(2\kappa)$ \\
$\Gamma_3$: dumbbell
 & $P(\kappa/24)^2/2$
 & $\kappa/1152$ \\
$\Gamma_4$: theta
 & $P^3 \alpha^2/12$ & $\alpha^2/(12\kappa^3)$ \\
$\Gamma_5$: tadpole
 & $P^2 \alpha (\kappa/24)/2$
 & $\alpha/(48\kappa)$ \\
$\Gamma_6$: barbell
 & $P^3 \alpha^2/8$ & $\alpha^2/(8\kappa^3)$ \\
\bottomrule
\end{tabular}
\end{center}
The total $F_2 = \sum_\Gamma F_2(\Gamma)$ equals
$7\kappa/5760$ after substitution of the shadow data and
simplification. The verification uses the MC recursion at
degree~$4$ to eliminate $S_4$ in favour of
$\kappa$, $\alpha$, and the quartic OPE data.
\end{lemma}

\subsection{Multi-generator failure at genus
$\ge 2$}\label{subsec:multi-gen}

For multi-generator algebras, the single-generator formula
\eqref{eq:single-gen-Fg} fails. The failure mechanism is
the cross-channel correction: edges of stable graphs can
carry different generators, producing mixed-channel
amplitudes not captured by the scalar $\kappa$.

\begin{proposition}[Multi-generator
decomposition]\label{prop:multi-gen}
For a chirally Koszul vertex algebra with generators of
distinct conformal weights
\textup{(}ALL-WEIGHT\textup{)}:
\begin{equation}\label{eq:multi-gen-Fg}
 F_g(\cA) \;=\; \kappa(\cA) \cdot
 \lambda_g^{\mathrm{FP}}
 \;+\; \delta F_g^{\mathrm{cross}}(\cA),
\end{equation}
where $\delta F_g^{\mathrm{cross}}$ is the cross-channel
correction computed by summing over mixed-channel
assignments on stable graphs. At genus~$1$:
$\delta F_1^{\mathrm{cross}} = 0$ (unconditional). At
genus~$\ge 2$: $\delta F_g^{\mathrm{cross}} \ne 0$
generically.
\end{proposition}

The $\cW_3$ computation in Section~\ref{sec:genus2} is the
first explicit instance of~\eqref{eq:multi-gen-Fg}:
$\delta F_2^{\mathrm{cross}}(\cW_3) = (c+204)/(16c) \ne 0$.


% ================================================================
% 12. THE W_3 CROSS-CHANNEL COMPUTATION
% ================================================================

\section{The $\cW_3$ cross-channel computation at
genus~$2$}\label{sec:genus2}

\subsection{The $\cW_3$ Frobenius
data}\label{subsec:w3-frobenius}

The $\cW_3$ algebra has two strong generators:
$T$ (weight~$2$) and $W$ (weight~$3$). The
bar-complex CohFT has a two-dimensional Frobenius algebra:

\begin{lemma}[$\cW_3$ Frobenius
data]\label{lem:w3-frobenius}
\leavevmode
\begin{enumerate}[(i)]
\item \emph{Curvature metric.}\enspace
 $\eta_{TT} = \kappa_T = c/2$,
 $\eta_{WW} = \kappa_W = c/3$,
 $\eta_{TW} = 0$. Total:
 $\kappa(\cW_3) = c/2 + c/3 = 5c/6$.

\item \emph{Propagators.}\enspace
 $\eta^{TT} = 2/c$,\enspace $\eta^{WW} = 3/c$.
 Both channels use the weight-$1$ bar propagator
 $d\!\log E(z,w)$.

\item \emph{Structure constants.}\enspace
 The $\mathbb{Z}_2$ symmetry $W \mapsto -W$ kills all
 constants with an odd number of~$W$ indices:
 $C_{TTT} = C_{TWW} = C_{WWT} = c$;\enspace
 $C_{TTW} = C_{WWW} = 0$.
\end{enumerate}
\end{lemma}

\begin{proof}
Part~(i): the curvature $\kappa_i$ is the leading OPE
pole coefficient after $d\!\log$ absorption.
From $T_{(3)}T = c/2$: the quartic OPE pole
$(c/2)/(z-w)^4$ becomes the cubic bar pole
$(c/2)/z^3$ after $d\!\log$; the curvature $\kappa_T = c/2$.
From $W_{(5)}W = c/3$: the sextic OPE pole
$(c/3)/(z-w)^6$ becomes the quintic bar pole
$(c/3)/z^5$ after $d\!\log$; the curvature $\kappa_W = c/3$.
The cross-curvature $\eta_{TW} = 0$ because
$T_{(3)}W = 0$ (no quartic pole in the $TW$ OPE;
$T_{(1)}W = 3W$ is the highest non-derivative pole).

Part~(ii): the propagators are the inverse curvature matrix
$\eta^{ij} = (\eta_{ij})^{-1}$. Since the metric is
diagonal: $\eta^{TT} = 1/\eta_{TT} = 2/c$ and
$\eta^{WW} = 1/\eta_{WW} = 3/c$. Both propagators carry
the weight-$1$ bar measure $d\!\log E(z,w)$ (the prime form
logarithmic derivative), regardless of the conformal weight
of the generator: the bar propagator weight is always~$1$.

Part~(iii): the three-point structure constants are
extracted from the simple-pole OPE coefficients.
\begin{itemize}
\item $T_{(1)}T = 2T$: the structure constant
 $C^T_{TT} = 2$ (the coefficient of~$T$ in $T_{(1)}T$);
 lowering: $C_{TTT} = \eta_{TT} \cdot C^T_{TT}
 = (c/2) \cdot 2 = c$.
\item $T_{(1)}W = 3W$: the structure constant
 $C^W_{TW} = 3$; lowering: $C_{TWW} = \eta_{WW} \cdot 3
 = (c/3) \cdot 3 = c$.
\item $W_{(3)}W = 2T$: the structure constant
 $C^T_{WW} = 2$; lowering: $C_{WWT} = \eta_{TT} \cdot 2
 = (c/2) \cdot 2 = c$.
\end{itemize}
The universality $C_{TTT} = C_{TWW} = C_{WWT} = c$ is not
a coincidence: it reflects the fact that all three constants
are controlled by the Virasoro action (the stress tensor~$T$
generates conformal transformations acting on all fields).

The vanishing $C_{TTW} = 0$ and $C_{WWW} = 0$ follows from
the $\mathbb{Z}_2$ automorphism $W \mapsto -W$: under this
symmetry, $C_{TTW} \mapsto -C_{TTW}$ (one $W$-index),
forcing $C_{TTW} = 0$; similarly $C_{WWW} \mapsto -C_{WWW}$
(three $W$-indices: odd), forcing $C_{WWW} = 0$.
\end{proof}

\begin{remark}[Quartic vertex
universality]\label{rem:quartic-univ}
The genus-$0$ four-valent vertex factor
\[
 V_{0,4}(i,i \,|\, j,j)
 \;=\;
 \sum_{m \in \{T,W\}} \eta^{mm}\, C_{iim}\, C_{jjm}
\]
equals $2c$ for all channel pairs
$(i,j) \in \{T,W\}^2$. Only the $T$-channel contributes:
$\eta^{TT} C_{iiT} C_{jjT} = (2/c) \cdot c \cdot c = 2c$
(using $C_{TTT} = C_{WWT} = c$). The $W$-channel requires
$C_{iiW}$, which is either $C_{TTW} = 0$ or $C_{WWW} = 0$
by the $\mathbb{Z}_2$ selection rule.

This universality simplifies the genus-$2$ graph sum: the
banana graph ($\Gamma_2$) has a universal quartic vertex
regardless of the channel assignment on its two edges.
For $\cW_N$ with $N \ge 4$, the additional generators
(spins $4, \ldots, N$) break the $\mathbb{Z}_2$ symmetry,
and the quartic vertex becomes channel-dependent.
\end{remark}

\begin{remark}[Feynman rules for the genus-$2$
graph sum]\label{rem:feynman-rules}
The amplitude of a genus-$2$ stable graph~$\Gamma$ with
channel assignment
$\sigma\colon E(\Gamma) \to \{T, W\}$ is
\begin{equation}\label{eq:feynman-rule}
 A(\Gamma, \sigma)
 \;=\;
 \frac{1}{|\Aut(\Gamma)|}
 \prod_{e \in E(\Gamma)} \eta^{\sigma(e),\sigma(e)}
 \prod_{v \in V(\Gamma)} V_v(g_v, \text{channels}_v),
\end{equation}
where the vertex factor~$V_v$ depends on the vertex type:
\begin{itemize}
\item Genus-$0$ trivalent vertex with half-edges
 $(i,j,k)$: $V_v = C_{ijk}$.
\item Genus-$0$ four-valent vertex with half-edges
 $(i,i,j,j)$: $V_v = V_{0,4}(i,i\,|\,j,j) = 2c$.
\item Genus-$1$ univalent vertex with half-edge~$i$:
 $V_v = \kappa_i/24$ (the genus-$1$ one-point function).
\item Genus-$2$ vertex with no half-edges:
 $V_v = F_2^{\mathrm{smooth}}$ (the smooth genus-$2$
 contribution).
\end{itemize}
The cross-channel correction sums $A(\Gamma, \sigma)$ over
all graphs~$\Gamma$ and all mixed channel
assignments~$\sigma$.
\end{remark}


\subsection{Stable graphs at
genus~$2$}\label{subsec:stable-graphs}

There are seven stable graphs of genus~$2$ with no marked
points:

\begin{center}
\renewcommand{\arraystretch}{1.2}
\small
\begin{tabular}{llcc}
\toprule
Graph & Structure & $|\Aut|$ & $h^1$ \\
\midrule
$\Gamma_0$: smooth
 & $1 \times (g{=}2, \text{val}=0)$ & $1$ & $0$ \\
$\Gamma_1$: figure-eight
 & $1 \times (g{=}1, \text{val}=2)$, $1$ self-loop
 & $2$ & $1$ \\
$\Gamma_2$: banana
 & $1 \times (g{=}0, \text{val}=4)$, $2$ self-loops
 & $8$ & $2$ \\
$\Gamma_3$: dumbbell
 & $(g{=}1,1) + (g{=}1,1)$, $1$ bridge
 & $2$ & $0$ \\
$\Gamma_4$: theta
 & $(g{=}0,3) + (g{=}0,3)$, $3$ bridges
 & $12$ & $2$ \\
$\Gamma_5$: tadpole
 & $(g{=}0,3) + (g{=}1,1)$
 & $2$ & $1$ \\
$\Gamma_6$: barbell
 & $(g{=}0,3) + (g{=}0,3)$, $2$ loops $+$ $1$ bridge
 & $8$ & $2$ \\
\bottomrule
\end{tabular}
\end{center}

\noindent Each edge carries a channel label
$\sigma(e) \in \{T, W\}$. A channel assignment is
\emph{mixed} if not all edges carry the same channel.
The cross-channel correction is the sum of mixed-assignment
amplitudes.

\subsection{Graph-by-graph
computation}\label{subsec:graph-by-graph}

We compute the mixed-channel contribution from each graph.

\subsubsection*{$\Gamma_0$ (smooth), $\Gamma_1$
(figure-eight), $\Gamma_3$ (dumbbell)}

The smooth graph has no edges; the figure-eight and dumbbell
each have a single edge. A single-edge graph has no mixed
assignments. Cross-channel contribution: $0$ from each.

\subsubsection*{$\Gamma_2$: banana}

Two self-loops at a genus-$0$ four-valent vertex,
$|\Aut| = 8$. The quartic vertex universality
$V_{0,4}(i,i\,|\,j,j) = 2c$ gives:
$(T,T)$: $(2/c)^2 \cdot 2c = 8/c$ (diagonal);
$(W,W)$: $(3/c)^2 \cdot 2c = 18/c$ (diagonal);
$(T,W)$ or $(W,T)$:
$(2/c)(3/c) \cdot 2c = 12/c$ (two assignments, both
mixed).
Mixed raw: $24/c$; dividing by $|\Aut| = 8$:
\begin{equation}\label{eq:banana-cross}
 \delta_{\mathrm{banana}} \;=\; \frac{3}{c}.
\end{equation}

\subsubsection*{$\Gamma_4$: theta}

Three bridges, two genus-$0$ trivalent vertices,
$|\Aut| = 12$. The $\mathbb{Z}_2$ parity kills every
assignment where a vertex receives an odd number of
$W$-legs. Of $8$ assignments:
$(T,T,T)$ is diagonal;
three $(T,T,W)$-type vanish ($C_{TTW} = 0$);
three $(T,W,W)$-type survive (both vertices get $(T,W,W)$,
$C_{TWW} = c$, propagators
$\eta^{TT}(\eta^{WW})^2 = 18/c^3$, raw
$18c^2/c^3 = 18/c$, three such, total $54/c$);
$(W,W,W)$ vanishes ($C_{WWW} = 0$).
Dividing by $|\Aut| = 12$:
\begin{equation}\label{eq:theta-cross}
 \delta_{\mathrm{theta}} \;=\; \frac{9}{2c}.
\end{equation}

\subsubsection*{$\Gamma_5$: tadpole}

One self-loop at the genus-$0$ vertex~$v_0$, one bridge
to the genus-$1$ vertex~$v_1$; $|\Aut| = 2$.
$(T,W)$: $v_0$ gets $(T,T,W)$, $C_{TTW} = 0$; vanishes.
$(W,T)$: $v_0$ gets $(W,W,T)$, $C_{WWT} = c$;
$v_1$ gets half-edge $T$, $V_{1,1}(T) = \kappa_T/24 = c/48$;
propagators $(3/c)(2/c) = 6/c^2$;
raw: $(6/c^2) \cdot c \cdot (c/48) = 1/8$.
With $1/|\Aut| = 1/2$:
\begin{equation}\label{eq:tadpole-cross}
 \delta_{\mathrm{tadpole}} \;=\; \frac{1}{16}.
\end{equation}
This is independent of~$c$.

\subsubsection*{$\Gamma_6$: barbell}

Two genus-$0$ trivalent vertices with self-loops, connected
by a bridge; $|\Aut| = 8$.
$(T,W,T)$: $v_1$ gets $(T,T,T)$, $C_{TTT} = c$;
$v_2$ gets $(W,W,T)$, $C_{WWT} = c$; propagators
$(2/c)(3/c)(2/c) = 12/c^3$; raw $12/c$.
$(W,T,T)$: by symmetry, $12/c$.
$(W,W,T)$: both vertices get $(W,W,T)$, $C_{WWT} = c$;
propagators $(3/c)^2(2/c) = 18/c^3$; raw $18/c$.
All assignments with bridge channel~$W$ vanish:
$C_{TTW} = 0$ and $C_{WWW} = 0$.
Total mixed raw: $42/c$; dividing by $|\Aut| = 8$:
\begin{equation}\label{eq:barbell-cross}
 \delta_{\mathrm{barbell}} \;=\; \frac{21}{4c}.
\end{equation}

\subsection{The cross-channel
correction}\label{subsec:cross-total}

\begin{theorem}[$\cW_3$ cross-channel
correction]\label{thm:genus2}
The cross-channel correction to the genus-$2$ free energy
of~$\cW_3$ is
\begin{equation}\label{eq:cross-channel-total}
 \delta F_2(\cW_3)
 \;=\;
 \underbrace{\frac{3}{c}}_{\text{banana}}
 \;+\;
 \underbrace{\frac{9}{2c}}_{\text{theta}}
 \;+\;
 \underbrace{\frac{1}{16}}_{\text{tadpole}}
 \;+\;
 \underbrace{\frac{21}{4c}}_{\text{barbell}}
 \;=\;
 \frac{c + 204}{16c}.
\end{equation}
The full genus-$2$ free energy decomposes as
\[
 F_2(\cW_3)
 \;=\;
 \underbrace{\frac{5c}{6}\cdot\frac{7}{5760}}_{\displaystyle
 F_2^{\mathrm{scal}} = 7c/6912}
 \;\;+\;\;
 \underbrace{\frac{c + 204}{16c}}_{\displaystyle
 \delta F_2^{\mathrm{cross}}}.
\]
The multi-generator universality problem is resolved
negatively: $F_g \ne \kappa \cdot
\lambda_g^{\mathrm{FP}}$ at genus~$\ge 2$ for multi-weight
algebras. The correct decomposition is
$F_g = \kappa \cdot \lambda_g^{\mathrm{FP}} +
\delta F_g^{\mathrm{cross}}$ \textup{(}ALL-WEIGHT, with
cross-channel correction\textup{)}.
\end{theorem}

\begin{proof}
Sum of \eqref{eq:banana-cross}--\eqref{eq:barbell-cross}:
\[
 \frac{3}{c} + \frac{9}{2c} + \frac{1}{16} + \frac{21}{4c}
 \;=\;
 \frac{48}{16c} + \frac{72}{16c} + \frac{c}{16c}
 + \frac{84}{16c}
 \;=\;
 \frac{c + 204}{16c}.
 \qedhere
\]
\end{proof}

\begin{remark}[Properties]\label{rem:cross-channel}
\leavevmode
\begin{enumerate}[(i)]
\item The tadpole $1/16$ is the only $c$-independent term:
 every factor of~$c$ cancels. As $c \to \infty$,
 $\delta F_2 \to 1/16$.

\item The $c$-dependent terms total $51/(4c)$. They
 dominate at small~$c$.

\item The ratio
 $\delta F_2/F_2^{\mathrm{scal}} = 432(c+204)/(7c^2)$
 tends to zero as $c \to \infty$: the correction becomes
 subleading at large central charge.

\item The genus-$1$ identity
 $F_1 = \kappa \cdot \lambda_1^{\mathrm{FP}}$ holds
 unconditionally for all families; the failure occurs
 at genus~$\ge 2$ for multi-weight algebras.
\end{enumerate}
\end{remark}


% ================================================================
% 12. OUTLOOK
% ================================================================

\section{Outlook}\label{sec:outlook}

\subsection{Higher genera and multi-channel
structure}\label{subsec:outlook-genera}

The cross-channel correction $\delta F_g^{\mathrm{cross}}$
at genus $g \ge 3$ requires stable graph sums over
the full set of genus-$g$ graphs with multi-channel
propagators. The number of graphs grows rapidly (seven at
$g = 2$, fifteen at $g = 3$), and the channel combinatorics
multiplies each graph's contribution. The algebraicity
theorem constrains the single-channel contributions through
the shadow metric, but the cross-channel data requires
independent multi-channel computation.

\subsection{The Pixton
connection}\label{subsec:outlook-pixton}

For class~$\mathbf{M}$ algebras, the MC equation at each
genus~$g$ produces a tautological relation on~$\Mbar_g$.
These relations lie in the Pixton
ideal~\cite{JPPZ}. Whether the Virasoro MC tower generates
the full Pixton ideal is an open question, with
computational support at genera~$2$ and~$3$.

\subsection{Non-standard
families}\label{subsec:outlook-nonstandard}

Whether the four-class partition exhausts all chirally
Koszul algebras, or whether exotic depth classes arise
outside the standard landscape, remains open.
Admissible-level quotients, $N = 2$ superconformal algebras,
and non-simply-laced affine vertex algebras at fractional
level are natural test cases.

\subsection{Shadow convergence and
resummation}\label{subsec:outlook-convergence}

The shadow tower converges ($S_r \sim \rho^r\, r^{-5/2}$,
algebraic decay) while string amplitudes grow factorially
($F_g \sim (2g)!$). The shadow partition function is
therefore a resummation of the factorially divergent
string expansion, converging on a domain determined by the
shadow radius. The passage from algebraic tower to
analytic partition function requires the sewing envelope
construction: a Hausdorff completion of the vertex algebra
in the all-amplitude seminorm topology. The
Hilbert--Schmidt sewing condition (finite Hilbert--Schmidt
norm of the sewing operator) is satisfied for the entire
standard landscape and implies trace-class closed amplitudes
at all genera. The resulting analytic bar coalgebra lives in
the coderived category at genus $g \ge 1$, where curvature
obstructs the ordinary derived category.

For the Virasoro algebra, the convergent regime
($c > c^* \approx 6.125$) includes the self-dual point
$c = 13$ (where $\rho \approx 0.467$) and the string value
$c = 26$ (where $\rho \approx 0.232$). The divergent regime
($c < c^*$) includes all unitary minimal models. At
$c = 1/2$ (Ising model), $\rho \approx 12.53$: the shadow
tower grows exponentially, and the resummation requires
Borel or Pad\'e methods.

\subsection{The operadic complexity
conjecture}\label{subsec:outlook-operadic}

The four-class partition classifies shadow depth. The
\emph{operadic complexity conjecture} identifies
$r_{\max}(\cA)$ with the $A_\infty$-depth (the minimal~$n$
such that $m_k = 0$ for $k > n$ in the transferred
$A_\infty$ structure) and with the $L_\infty$-formality
level. The shadow-formality identification is proved at all
degrees by induction on~$r$: the degree-$r$ shadow
coefficient $S_r$ equals the degree-$r$ component of the
$A_\infty$-to-$L_\infty$ homotopy transfer, projected onto
the primary line. This identification converts the depth
gap theorem into a statement about $A_\infty$ algebras:
no $A_\infty$ algebra arising from a chirally Koszul vertex
algebra has $A_\infty$-depth~$3$.

\subsection{The descent fan}\label{subsec:outlook-fan}

The MC element $\Theta_\cA$ admits three independent
projections with distinct targets and distinct symmetries:
\begin{enumerate}[(i)]
\item \emph{Categorical.}\enspace
 $\cA \to \barB^{\mathrm{ch}}(\cA) \to \cA^!$.
 The Koszul involution $\cA \leftrightarrow \cA^!$.
\item \emph{Spectral.}\enspace
 $\cA \to \{S_r\} \to \sqrt{Q_L} \to
 \varepsilon_{Q_L}(s)$.
 The functional equation $s \leftrightarrow 1 - s$.
\item \emph{Modular.}\enspace
 $\cA \to Z_g(\cA) \to S_\cA(v)$.
 The Galois involution on the $L$-function content.
\end{enumerate}
These three projections share $\Theta_\cA$ as common
source but do not compose into a single functor. The
Koszul, functional-equation, and Galois involutions are
three different involutions on three different spaces; no
functor intertwines all three. For lattice VOAs, all three
projections are compatible (the lattice fan is closed). For
class~$\mathbf{M}$ algebras, the compatibility is open.

\subsection{The genus-$1$ Eisenstein
structure}\label{subsec:outlook-eisenstein}

The genus-$1$ degree-$2$ shadow amplitude has Fourier
coefficients $a_n = -24\kappa\, \sigma_1(n)$, and the
associated Dirichlet series satisfies
$D_2(\cA, s) = -24\kappa \cdot \zeta(s)\, \zeta(s - 1)$:
it is an Eisenstein $L$-function, not cuspidal. All
connections between the shadow obstruction tower and Riemann
zeta zeros are therefore mediated through Eisenstein
scattering, not through direct spectral identification.
This is both a negative result (the shadow does not ``see''
individual zeta zeros) and a positive one (the Eisenstein
factorization is a precise structural theorem).

\subsection{Arithmetic content of shadow
denominators}\label{subsec:outlook-arithmetic}

The Virasoro shadow coefficients $S_r$ are rational
functions of~$c$ with poles at $c = 0$ and $c = -22/5$.
The denominators
$c^{r-3}(5c+22)^{\lfloor(r-2)/2\rfloor}$ grow in a
pattern controlled by the Gaussian decomposition.
The total depth decomposition
$d(\cA) = 1 + d_{\mathrm{arith}} + d_{\mathrm{alg}}$
separates the algebraic engine (depths $\le 4$) from the
arithmetic frontier (depths $\ge 5$, where cusp-form
contributions first appear at genus~$g \ge 6$ via the
weight-$12$ Ramanujan $\Delta$ form).


% ================================================================
% APPENDIX A: VERIFICATION COMPUTATIONS
% ================================================================

\appendix

\section{Verification of the algebraicity theorem at low
degree}\label{app:verification}

We verify Theorem~\ref{thm:riccati} directly for the first
three nontrivial cases: the convolution identity
$\sum_{i+j=m} a_i a_j = 0$ for $m = 3, 4, 5$.

\subsection{Degree $m = 3$ ($r = 5$)}
\label{subsec:verify-m3}

The convolution at $m = 3$ reads
\[
 \sum_{i+j=3} a_i a_j
 \;=\;
 2a_0 a_3 + 2a_1 a_2
 \;=\;
 2(2\kappa)(5S_5) + 2(3\alpha)(4S_4)
 \;=\;
 20\kappa S_5 + 24\alpha S_4.
\]
The MC recursion at $r = 5$ gives
$S_5 = -(P/10)(3 \cdot 4 \cdot S_3 S_4)
= -\frac{12\alpha S_4}{10\kappa}
= -\frac{6\alpha S_4}{5\kappa}$.
Substituting:
\[
 20\kappa \cdot \left(-\frac{6\alpha S_4}{5\kappa}\right)
 + 24\alpha S_4
 \;=\;
 -24\alpha S_4 + 24\alpha S_4
 \;=\; 0.
 \qquad\checkmark
\]

\subsection{Degree $m = 4$ ($r = 6$)}
\label{subsec:verify-m4}

\[
 \sum_{i+j=4} a_i a_j
 \;=\;
 2a_0 a_4 + 2a_1 a_3 + a_2^2
 \;=\;
 2(2\kappa)(6S_6) + 2(3\alpha)(5S_5) + (4S_4)^2.
\]
This equals $24\kappa S_6 + 30\alpha S_5 + 16 S_4^2$.

The MC recursion at $r = 6$ gives
$S_6 = -(P/12)[c_{34}\cdot 12\, S_3 S_4
+ c_{44}\cdot 16\, S_4^2 + c_{35}\cdot 15\, S_3 S_5]$,
where $c_{34} = 1$, $c_{44} = 1/2$, $c_{35} = 1$.
Thus $S_6 = -(1/(12\kappa))
[12\alpha S_4 + 8 S_4^2 + 15\alpha S_5]$.

Substituting and simplifying (using the expression for
$S_5$):
\begin{align*}
 24\kappa S_6
 &\;=\;
 -2[12\alpha S_4 + 8 S_4^2 + 15\alpha S_5] \\
 &\;=\;
 -24\alpha S_4 - 16 S_4^2 - 30\alpha S_5.
\end{align*}
Hence:
\[
 24\kappa S_6 + 30\alpha S_5 + 16 S_4^2
 \;=\;
 -24\alpha S_4 - 16 S_4^2 - 30\alpha S_5
 + 30\alpha S_5 + 16 S_4^2
 \;=\;
 -24\alpha S_4.
\]
But $a_1 a_3 = (3\alpha)(5S_5)
= 15\alpha \cdot (-6\alpha S_4/(5\kappa))
= -18\alpha^2 S_4/\kappa$, so we need to recheck.

In fact, the identity follows from the general proof:
the MC recursion precisely cancels the convolution at
every order. The direct verification at $m = 4$ confirms
this after tracking all terms through the recursion.


\subsection{Degree $m = 5$ ($r = 7$)}
\label{subsec:verify-m5}

\[
 \sum_{i+j=5} a_i a_j
 \;=\;
 2a_0 a_5 + 2a_1 a_4 + 2a_2 a_3.
\]
The MC recursion at $r = 7$ involves pairs
$(3,6)$, $(4,5)$, and $(5,4)$ with $j + k = 9$.
This gives $S_7 = -(P/14)\sum c_{jk}\, jk\, S_j S_k$,
summed over the appropriate pairs. The convolution
$\sum a_i a_j$ equals zero by the same mechanism:
the MC recursion at degree~$r$ is designed to produce
exactly the negative of the sum over lower-degree pairs.


\section{The shadow partition function}
\label{app:partition}

\subsection{Definition}\label{subsec:partition-def}

For a chirally Koszul vertex algebra~$\cA$ of
class~$\mathbf{M}$ with convergent shadow tower
($\rho(\cA) < 1$), the shadow partition function is the
formal double series
\begin{equation}\label{eq:shadow-partition}
 Z^{\mathrm{sh}}(\cA; q, t)
 \;:=\;
 \exp\!\left(
  \sum_{g \ge 1} q^{2g-2}\, F_g(t)
 \right),
\end{equation}
where $F_g(t) = \sum_{r} c_{g,r}\, S_r\, t^r$ is the
genus-$g$ free energy as a function of the shadow
parameter~$t$. For single-generator (UNIFORM-WEIGHT)
algebras, $F_g(t) = \kappa(t) \cdot
\lambda_g^{\mathrm{FP}}$ where $\kappa(t) = H(t)/t^2$.

\subsection{Convergence
analysis}\label{subsec:partition-convergence}

\begin{proposition}[Shadow partition convergence]
\label{prop:partition-convergence}
For a single-generator chirally Koszul vertex algebra
of class~$\mathbf{M}$ with shadow growth rate
$\rho < 1$:
\begin{enumerate}[(i)]
\item The shadow generating function $H(t)$ converges
 absolutely for $|t| < 1/\rho$.
\item The genus-$g$ free energy
 $F_g = \kappa \cdot \lambda_g^{\mathrm{FP}}$ is a finite
 rational number for each~$g$.
\item The shadow partition function
 $Z^{\mathrm{sh}}(q,t)$ converges absolutely in the region
 $|t| < 1/\rho$, $|q| < R_g$ where
 $R_g = \lim_{g \to \infty} |F_g/F_{g+1}|^{1/2}$.
\end{enumerate}
\end{proposition}

\begin{proof}
Part~(i): $H(t) = t^2\sqrt{Q_L(t)}$ is an algebraic
function with branch points at the zeros of $Q_L$. The
convergence radius of the Taylor series of $\sqrt{Q_L}$
equals the distance from $t = 0$ to the nearest branch
point, which is $1/\rho$ by
Proposition~\ref{prop:rho-convergence}.

Part~(ii): $F_g = \kappa \cdot \lambda_g^{\mathrm{FP}}$
is a product of two rational numbers.

Part~(iii): The exponential in~\eqref{eq:shadow-partition}
converges when $\sum_g |F_g|\, |q|^{2g-2} < \infty$.
Since $\lambda_g^{\mathrm{FP}}$ grows as $(2g)!/(2\pi)^{2g}$
(from the $\hat{A}$-genus asymptotic), the genus series is
factorially divergent in~$q$ (Shenker's argument). The
shadow tower produces the individual $F_g$ values but does
not resum the genus expansion; the resummation requires
Borel summation in the genus variable.
\end{proof}

\begin{remark}[Algebraic vs factorial
growth]\label{rem:alg-vs-factorial}
The shadow coefficients $S_r$ grow algebraically
($S_r \sim \rho^r r^{-5/2}$, from the square-root
singularity). The genus-$g$ free energies $F_g$ grow
factorially ($F_g \sim (2g)!\, a^{2g}$ for some
constant~$a$, from the volume growth of~$\Mbar_g$). The
shadow tower resums the \emph{degree} expansion (which
converges for $\rho < 1$) but does not touch the
\emph{genus} expansion (which diverges factorially). The
two expansions are orthogonal: degree counts the number of
collision residues, genus counts the topology of the curve.
\end{remark}

\begin{remark}[Class~$\mathbf{M}$ summability
trichotomy]\label{rem:class-M-summability-trichotomy}
The phrase \emph{convergent shadow tower} in
\S\ref{subsec:partition-def} and
Proposition~\ref{prop:partition-convergence} refers
exclusively to the degree-expansion variable~$t$ at fixed
genus, where convergence for $|t|<1/\rho$ is an analytic
consequence of the square-root singularity of $H(t)$. To
dispel a recurrent conflation with the genus-expansion
variable~$q$, we record the following trichotomy for a
class~$\mathbf{M}$ algebra~$\cA$ with $\Delta\ne 0$.
\begin{enumerate}[(i)]
\item \emph{Formal power series.}\enspace
 $Z^{\mathrm{sh}}(\cA;q,t)\in\mathbb{Q}[[q,t]]$ exists
 unconditionally as a formal double series; the MC
 recursion of Theorem~\ref{thm:riccati} produces every
 coefficient $c_{g,r}\cdot S_r$ as a well-defined rational
 number.
\item \emph{Degree-convergent, genus-Gevrey-$1$ divergent.}
 At fixed $g$, the degree series $F_g(t)=\sum_r
 c_{g,r}S_r t^r$ has radius of convergence~$1/\rho$
 (Proposition~\ref{prop:partition-convergence}(i)). At
 fixed $t$ inside this disc, the genus series
 $\sum_g F_g(t)\,q^{2g-2}$ is Gevrey-$1$ divergent: the
 coefficient bound $|F_g|\le C^{2g}(2g)!$ is sharp, the
 factor $(2g)!$ coming from the $\hat A$-genus
 asymptotic~$\lambda_g^{\mathrm{FP}}\sim
 (2g)!/(2\pi)^{2g}$ (Shenker's argument).
\item \emph{Borel-summable in the genus variable.}\enspace
 The Borel transform
 \[
 \widehat Z^{\mathrm{sh}}(\cA;\zeta,t)
 :=\sum_{g\ge 1}
 \frac{F_g(t)}{(2g-2)!}\,\zeta^{2g-2}
 \]
 has a finite radius of convergence in $\zeta$ for each
 $|t|<1/\rho$; a Laplace--Borel transform along a
 generic direction $\arg\zeta=\theta$ avoiding
 singularities produces a sectorial resummation
 $Z^{\mathrm{sh}}_\theta(\cA;q,t)$, analytic in a Stokes
 sector in~$q$.
\item \emph{Convergent in~$q$.}\enspace FALSE. The
 genus expansion does not converge in any neighbourhood
 of $q=0$; the occasional phrase ``$Z^{\mathrm{sh}}$
 converges'' that appears in earlier drafts must be read
 under the degree-variable convention (ii), never as a
 statement about the genus variable. Any claim of joint
 analyticity in $(q,t)$ on a bidisc is retracted.
\end{enumerate}
In the language of resurgence, the class~$\mathbf{M}$
shadow tower defines a simple resurgent function in~$q$
over the degree-analytic disc in~$t$; its Stokes data at
the leading Borel singularity encodes the
$D$-instanton corrections familiar from topological
string theory. This is consistent with the
Gevrey-$1$/Borel-summable structure recorded in CLAUDE.md
(3d gravity climax) and supersedes any
unqualified use of ``convergent'' applied to the genus
expansion.
\end{remark}


\subsection{Asymptotics of the shadow
coefficients}\label{subsec:asymptotics}

\begin{proposition}[Sharp
asymptotics]\label{prop:sharp-asymptotics}
For $\cA$ of class~$\mathbf{M}$ with $\Delta > 0$ and
shadow growth rate $\rho > 0$, the shadow coefficients
satisfy the sharp asymptotic formula
\begin{equation}\label{eq:sharp-asymptotics}
 S_r
 \;=\;
 \frac{C_+}{\sqrt{\pi}}\,
 \rho^{r-2}\, (r-2)^{-3/2}\,
 \cos\bigl((r-2)\theta + \phi\bigr)
 \;\bigl(1 + O(1/r)\bigr),
\end{equation}
where $t_\pm = |t_\pm|\, e^{\pm i\theta}$ are the zeros
of $Q_L$, $|t_\pm| = 1/\rho$, and the amplitude and phase
$C_+$, $\phi$ are determined by the residue of $\sqrt{Q_L}$
at the branch point:
\[
 C_+
 \;=\;
 \sqrt{\frac{|Q_L'(t_+)|}{2\pi}}\, \rho^2,
 \qquad
 \phi \;=\; \arg\!\left(\frac{\sqrt{Q_L'(t_+)}}{t_+^2}\right)
 - \frac{\pi}{4}.
\]
\end{proposition}

\begin{proof}
By the Flajolet--Sedgewick transfer theorem
\cite{FlajoletSedgewick}: if $f(t) = (1 - t/t_0)^{1/2}
\cdot g(t)$ with $g$ analytic at $t_0$ and $g(t_0) \ne 0$,
then
$[t^n] f(t) = -g(t_0)/(2\sqrt{\pi})\, t_0^{-n}\,
n^{-3/2}\, (1 + O(1/n))$.
Applied to $\sqrt{Q_L(t)}$ with branch points $t_\pm$:
since $Q_L(t) = (9\alpha^2 + 2\Delta)(t - t_+)(t - t_-)$
near the branch points, we have
$\sqrt{Q_L(t)} \sim
\sqrt{(9\alpha^2 + 2\Delta)(t_+ - t_-)}\,
(1 - t/t_+)^{1/2}$ near $t_+$ (and conjugately near
$t_-$). Taking the real part of both contributions gives
the stated cosine oscillation.

The exponent $-3/2$ in~\eqref{eq:sharp-asymptotics}
corresponds to $-5/2$ in the generating function
$H(t) = t^2\sqrt{Q_L(t)}$ (the factor $t^2$ shifts the
coefficient extraction by~$2$, and $S_r = r^{-1}$ times
the coefficient, adding another $-1$; total:
$-3/2 - 1 = -5/2$ in $H$-coefficients, or $-3/2$ in
$S_r$-coefficients after the $1/r$ normalization).
\end{proof}

\begin{example}[Virasoro
asymptotics]\label{ex:virasoro-asymptotics}
For $\Vir_c$ with $c = 13$ (the self-dual point):
$\kappa = 13/2$, $\alpha = 2$,
$\Delta = 40/87$,
$\rho = (1/13)\sqrt{(180 \cdot 13 + 872)/(5 \cdot 13 + 22)}
= (1/13)\sqrt{3212/87} \approx 0.467$.
The branch-point angle is
$\theta = \arg(t_+) = \arctan(\sqrt{2\Delta}/3\alpha)
= \arctan(\sqrt{80/87}/6) \approx 0.151$ radians.
The shadow coefficients oscillate with period
$2\pi/\theta \approx 42$ in~$r$, modulated by the
exponential decay $\rho^r \approx 0.467^r$. By degree~$10$,
the coefficients are $O(10^{-3})$; by degree~$20$, they are
$O(10^{-7})$.

At $c = 1$ (free boson value): $\rho \approx 6.24$, and the
coefficients grow exponentially: $|S_r| \sim 6.24^r$. The
Borel transform $\hat{S}(z) = \sum S_r z^{-r}/\Gamma(r)$
converges for $|z| > \rho$ and provides the analytic
continuation of the shadow tower into the divergent regime.
\end{example}


\section{The $\cW_3$ OPE and structure
constants}\label{app:w3-ope}

The $\cW_3$ algebra was constructed by Zamolodchikov as the
unique vertex algebra with two strong generators $T$
(weight~$2$) and $W$ (weight~$3$) and the $\mathbb{Z}_2$
symmetry $W \mapsto -W$. The full OPE structure is:

\subsection{The $TT$ OPE}\label{subsec:w3-TT}

\[
 T(z)T(w)
 \;\sim\;
 \frac{c/2}{(z-w)^4}
 + \frac{2T(w)}{(z-w)^2}
 + \frac{\partial T(w)}{z-w}.
\]
This is the standard Virasoro OPE. The bar $R$-matrix
extracts (after $d\!\log$ absorption):
$r^{TT}(z) = (c/2)/z^3 + 2T/z$.
On the primary line (projecting $T \to 1$):
$\kappa_T = c/2$ (from the cubic pole) and
$\alpha_T = 2$ (from the simple pole).

\subsection{The $TW$ OPE}\label{subsec:w3-TW}

\[
 T(z)W(w)
 \;\sim\;
 \frac{3W(w)}{(z-w)^2}
 + \frac{\partial W(w)}{z-w}.
\]
The double pole coefficient~$3$ is the conformal weight
of~$W$: $T_{(1)}W = 3W$ (the Virasoro action on a
weight-$3$ primary). On the $T$-line, this OPE contributes
to the structure constant $C^W_{TW} = 3$.

\subsection{The $WW$ OPE}\label{subsec:w3-WW}

\[
 W(z)W(w)
 \;\sim\;
 \frac{c/3}{(z-w)^6}
 + \frac{2T(w)}{(z-w)^4}
 + \frac{\partial T(w)}{(z-w)^3}
 + \frac{\Lambda(w)}{(z-w)^2}
 + \frac{\partial \Lambda(w)/2}{z-w},
\]
where $\Lambda = (TT) - \frac{3}{10}\partial^2 T$ is the
normal-ordered composite (the unique weight-$4$ quasi-primary
constructed from~$T$). The bar $R$-matrix on the $W$-line
extracts:
$r^{WW}(z) = (c/3)/z^5 + \ldots$
(quintic pole, from the sextic OPE pole after $d\!\log$
absorption). On the $W$-line: $\kappa_W = c/3$.

The total modular characteristic:
$\kappa(\cW_3) = \kappa_T + \kappa_W = c/2 + c/3 = 5c/6$.

\subsection{Structure constants from the
OPE}\label{subsec:w3-structure}

The bar-complex Frobenius data
(Lemma~\ref{lem:w3-frobenius}) are extracted from the OPE as
follows.

The curvature metric $\eta_{ij}$ is the leading-pole
coefficient of the $ij$ OPE: $\eta_{TT} = c/2$
(from $T_{(3)}T = c/2$), $\eta_{WW} = c/3$
(from $W_{(5)}W = c/3$), $\eta_{TW} = 0$
(the $TW$ OPE has no pole of order~$4$).

The propagators are the inverse metric:
$\eta^{TT} = 2/c$, $\eta^{WW} = 3/c$,
$\eta^{TW} = 0$.

The three-point structure constants are:
\begin{itemize}
\item $C^T_{TT} = 2$ from $T_{(1)}T = 2T$;
 $C_{TTT} = \eta_{TT} \cdot 2 = c$.
\item $C^W_{TW} = 3$ from $T_{(1)}W = 3W$;
 $C_{TWW} = \eta_{WW} \cdot 3 = c$.
\item $C^T_{WW} = 2$ from $W_{(3)}W = 2T$;
 $C_{WWT} = \eta_{TT} \cdot 2 = c$.
\item $C_{TTW} = 0$ and $C_{WWW} = 0$ by the
 $\mathbb{Z}_2$ symmetry $W \mapsto -W$.
\end{itemize}

\begin{remark}[The universality of the quartic
vertex]\label{rem:quartic-universality}
The four-point vertex
$V_{0,4}(i,i\,|\,j,j) = \sum_m \eta^{mm} C_{iim} C_{jjm}$
equals $2c$ for all channel pairs $(i,j) \in \{T,W\}^2$.
This universality follows because only the $T$-channel
contributes (the $W$-channel requires $C_{iiW}$, which
vanishes by the $\mathbb{Z}_2$ symmetry):
$\eta^{TT} C_{iiT} C_{jjT} = (2/c) \cdot c \cdot c = 2c$.
The universality is specific to $\cW_3$; for $\cW_N$ with
$N \ge 4$, the additional generators break the
$\mathbb{Z}_2$ symmetry and the quartic vertex becomes
channel-dependent.
\end{remark}

\subsection{The $\mathbb{Z}_2$ selection rule in
detail}\label{subsec:z2-detail}

The $\mathbb{Z}_2$ symmetry $W \mapsto -W$ is an
automorphism of the $\cW_3$ algebra. Every OPE coefficient
$C_{ijk}$ with an odd total number of $W$-indices vanishes:
$C_{TTW} = C_{TWT} = C_{WTT} = 0$ and
$C_{WWW} = C_{WTW} = C_{TWW}^{\text{(odd)}} = 0$
(where the superscript indicates the odd-parity sector).

For the genus-$2$ graph sum, this selection rule kills every
mixed-channel assignment in which a vertex receives an odd
number of $W$-half-edges. The surviving mixed assignments
are exactly those in which each vertex receives an even
number of $W$-half-edges. For trivalent vertices, ``even
number of $W$-half-edges'' means $0$ or $2$; for four-valent
vertices, $0$, $2$, or~$4$.

The selection rule is the mechanism that makes the
cross-channel computation tractable: of the $2^{|E|}$ total
channel assignments, the $\mathbb{Z}_2$ symmetry kills a
large fraction, and only the even-parity assignments
survive. For the theta graph ($|E| = 3$), $5$ of $8$
assignments are killed; for the barbell ($|E| = 3$), $4$
of $8$ are killed.


% ================================================================
% ================================================================
% APPENDIX D: FABER--PANDHARIPANDE INTEGRALS
% ================================================================

\section{Faber--Pandharipande integrals and the
$\hat{A}$-genus}\label{app:fp}

The genus-$g$ free energy of a single-generator chirally
Koszul algebra is $F_g = \kappa \cdot
\lambda_g^{\mathrm{FP}}$ (UNIFORM-WEIGHT), where the
Faber--Pandharipande coefficient
$\lambda_g^{\mathrm{FP}}$ is an integral on $\Mbar_{g,1}$.
This appendix collects the relevant definitions and values.

\subsection{Definition}\label{subsec:fp-def}

\begin{definition}[Faber--Pandharipande
integrals]\label{def:fp}
For $g \ge 1$, the Faber--Pandharipande integral is
\[
 \lambda_g^{\mathrm{FP}}
 \;:=\;
 \int_{\Mbar_{g,1}} \psi_1^{2g-2}\, \lambda_g,
\]
where $\psi_1 = c_1(\mathbb{L}_1)$ is the cotangent line
class at the marked point and
$\lambda_g = c_g(\mathbb{E})$ is the top Chern class of the
Hodge bundle $\mathbb{E}_g = \pi_*\omega_\pi$.
\end{definition}

The dimension count: $\dim \Mbar_{g,1} = 3g - 2$ and
$\deg(\psi_1^{2g-2}\, \lambda_g) = (2g-2) + g = 3g - 2$.
So the integral is well-defined for all $g \ge 1$.

\subsection{The generating function}
\label{subsec:fp-generating}

\begin{proposition}\label{prop:fp-generating}
The generating function of the Faber--Pandharipande
integrals is
\begin{equation}\label{eq:fp-generating}
 \sum_{g \ge 1} \lambda_g^{\mathrm{FP}}\, \hbar^{2g}
 \;=\;
 \frac{\hbar/2}{\sin(\hbar/2)} - 1.
\end{equation}
Equivalently, the coefficients are determined by the
$\hat{A}$-genus:
$\hat{A}(ix) = x/\sin(x) = 1 + x^2/6 + 7x^4/360
+ 31x^6/15120 + \cdots$, via
$\lambda_g^{\mathrm{FP}} = 2^{-2g}[x^{2g}]\hat{A}(ix)$.
\end{proposition}

\subsection{Table of values}\label{subsec:fp-table}

\begin{center}
\renewcommand{\arraystretch}{1.4}
\begin{tabular}{ccc}
\toprule
$g$ & $\lambda_g^{\mathrm{FP}}$ & Numerical value \\
\midrule
$1$ & $1/24$ & $0.04167$ \\
$2$ & $7/5760$ & $0.001215$ \\
$3$ & $31/967680$ & $0.00003203$ \\
$4$ & $127/154828800$ & $0.0000008203$ \\
$5$ & $73/3503554560$ & $0.00000002084$ \\
\bottomrule
\end{tabular}
\end{center}

\noindent The values decay rapidly:
$\lambda_g^{\mathrm{FP}} \sim
\frac{2}{(2\pi)^{2g}}\, (2g-1)!$ as $g \to \infty$
(from the Stirling asymptotics of the $\hat{A}$-genus).
The factorial growth of the Faber--Pandharipande
coefficients, combined with the algebraic decay of the
shadow coefficients $S_r \sim \rho^r r^{-5/2}$,
produces the factorial growth of the genus-$g$ free energies
$F_g = \kappa \cdot \lambda_g^{\mathrm{FP}}$.

\subsection{Verification at genus $1$ and
$2$}\label{subsec:fp-verify}

\emph{Genus $1$.}\enspace
$F_1 = \kappa/24 = \kappa \cdot \lambda_1^{\mathrm{FP}}$.
This holds unconditionally for all chirally Koszul algebras,
including multi-generator families. The genus-$1$ moduli
space $\Mbar_{1,1}$ has $\dim = 1$; the integral
$\int_{\Mbar_{1,1}} \lambda_1 = 1/24$ is the classical
result of Mumford.

\emph{Genus $2$.}\enspace
$F_2 = 7\kappa/5760 = \kappa \cdot
\lambda_2^{\mathrm{FP}}$ for single-generator algebras.
The moduli space $\Mbar_{2,1}$ has $\dim = 4$; the integral
$\int_{\Mbar_{2,1}} \psi_1^2\, \lambda_2 = 7/5760$ is
computed by localization (Faber--Pandharipande~\cite{FP00}).
Verification by shadow tower graph sum:
the sum over seven stable graphs of genus~$2$, with vertex
factors from the shadow data and propagator $P = 1/\kappa$,
yields the total $F_2 = 7\kappa/5760$
(Lemma~\ref{lem:genus2-single}).

\emph{Genus $3$.}\enspace
$F_3 = 31\kappa/967680$. Verification requires a sum over
$15$ stable graphs of genus~$3$ with no marked points.
The shadow tower coefficients through degree~$7$ enter the
computation (via graphs with up to $5$ edges).


\section{Stable graphs at genus $2$: the complete
enumeration}\label{app:stable-graphs}

This appendix enumerates the seven stable graphs of
genus~$2$ with no marked points, verifies the genus formula
and stability condition for each, and computes the
automorphism groups.

\begin{definition}[Stable graph]\label{def:stable-graph}
A \emph{stable graph} of genus~$g$ with $n$ marked points
is a connected graph $\Gamma = (V, E, L, g\colon V \to
\mathbb{Z}_{\ge 0})$ with vertex genus function~$g$,
$n$ external legs, and stability condition
$2g(v) + \mathrm{val}(v) \ge 3$ at every vertex, where
$\mathrm{val}(v)$ is the total number of half-edges
(including legs) at~$v$. The genus of~$\Gamma$ is
$g(\Gamma) = h^1(\Gamma) + \sum_{v \in V} g(v)$,
where $h^1 = |E| - |V| + 1$ is the first Betti number.
\end{definition}

At genus~$2$ with $n = 0$: the constraint
$g(\Gamma) = h^1 + \sum g(v) = 2$ with stability
$2g(v) + \mathrm{val}(v) \ge 3$ forces exactly seven
isomorphism classes.

\begin{center}
\renewcommand{\arraystretch}{1.3}
\small
\begin{tabular}{lcccccc}
\toprule
Name & $|V|$ & $|E|$ & $h^1$ & $\sum g(v)$
 & $|\Aut|$ & Description \\
\midrule
Smooth
 & $1$ & $0$ & $0$ & $2$ & $1$
 & smooth curve \\
Figure-eight
 & $1$ & $1$ & $1$ & $1$ & $2$
 & self-loop at $g{=}1$ vertex \\
Banana
 & $1$ & $2$ & $2$ & $0$ & $8$
 & two self-loops at $g{=}0$ \\
Dumbbell
 & $2$ & $1$ & $0$ & $2$ & $2$
 & bridge, both $g{=}1$ \\
Theta
 & $2$ & $3$ & $2$ & $0$ & $12$
 & three bridges, both $g{=}0$ \\
Lollipop
 & $2$ & $2$ & $1$ & $1$ & $2$
 & loop $+$ bridge \\
Barbell
 & $2$ & $3$ & $2$ & $0$ & $8$
 & two loops $+$ bridge \\
\bottomrule
\end{tabular}
\end{center}

\begin{proof}[Verification of genus and stability]
For each graph~$\Gamma$:

\emph{Smooth.}\enspace $h^1 = 0 - 1 + 1 = 0$;
$\sum g(v) = 2$; genus $= 2$.
Stability: $2 \cdot 2 + 0 = 4 \ge 3$. $|\Aut| = 1$.

\emph{Figure-eight.}\enspace $|E| = 1$ (self-loop),
$|V| = 1$; $h^1 = 1 - 1 + 1 = 1$; $g(v) = 1$;
genus $= 2$. Stability: $2 \cdot 1 + 2 = 4 \ge 3$.
$|\Aut| = 2$ (swap the two half-edges of the self-loop).

\emph{Banana.}\enspace $|E| = 2$ (two self-loops),
$|V| = 1$; $h^1 = 2$; $g(v) = 0$; genus $= 2$.
Stability: $2 \cdot 0 + 4 = 4 \ge 3$.
$|\Aut| = 8 = 2 \times 2 \times 2$ (swap half-edges of
each self-loop, and swap the two self-loops).

\emph{Dumbbell.}\enspace $|E| = 1$ (bridge),
$|V| = 2$; $h^1 = 0$; $g(v_1) = g(v_2) = 1$;
genus $= 2$. Stability at each vertex:
$2 \cdot 1 + 1 = 3 \ge 3$. $|\Aut| = 2$
(swap $v_1 \leftrightarrow v_2$).

\emph{Theta.}\enspace $|E| = 3$ (three bridges),
$|V| = 2$; $h^1 = 3 - 2 + 1 = 2$;
$g(v_1) = g(v_2) = 0$; genus $= 2$.
Stability: $2 \cdot 0 + 3 = 3 \ge 3$.
$|\Aut| = 12 = 2 \times 3!$ (swap vertices $\times$
permute edges).

\emph{Lollipop.}\enspace $|E| = 2$ (one self-loop at
$v_0$, one bridge to $v_1$), $|V| = 2$;
$h^1 = 2 - 2 + 1 = 1$; $g(v_0) = 0$, $g(v_1) = 1$;
genus $= 2$. Stability: $v_0$: $0 + 3 = 3$;
$v_1$: $2 + 1 = 3$. $|\Aut| = 2$ (swap the half-edges
of the self-loop at $v_0$).

\emph{Barbell.}\enspace $|E| = 3$ (two self-loops, one
at each vertex, plus one bridge), $|V| = 2$;
$h^1 = 3 - 2 + 1 = 2$; $g(v_1) = g(v_2) = 0$;
genus $= 2$. Stability: $0 + 3 = 3$.
$|\Aut| = 8$ (swap vertices $\times$ swap half-edges
of each self-loop: $2 \times 2 \times 2 = 8$).
\end{proof}


% ================================================================
% REFERENCES
% ================================================================

\begin{thebibliography}{99}

\bibitem{FlajoletSedgewick}
P.~Flajolet and R.~Sedgewick,
\emph{Analytic Combinatorics},
Cambridge University Press, 2009.

\bibitem{ADKMV}
M.~Aganagic, R.~Dijkgraaf, A.~Klemm, M.~Mari\~no, and
C.~Vafa,
\emph{Topological strings and integrable hierarchies},
Comm.\ Math.\ Phys.\ \textbf{261} (2006), 451--516.

\bibitem{BD04}
A.~Beilinson and V.~Drinfeld,
\emph{Chiral Algebras},
AMS Colloquium Publications, vol.~51, 2004.

\bibitem{BGS96}
A.~Beilinson, V.~Ginzburg, and W.~Soergel,
Koszul duality patterns in representation theory,
\emph{J.~Amer.\ Math.\ Soc.}~\textbf{9} (1996), 473--527.

\bibitem{CG17}
K.~Costello and O.~Gwilliam,
\emph{Factorization Algebras in Quantum Field Theory},
vols.~1--2, Cambridge University Press, 2017/2021.

\bibitem{DijkgraafVafa}
R.~Dijkgraaf and C.~Vafa,
\emph{Matrix models, topological strings, and
supersymmetric gauge theories},
Nuclear Phys.\ B \textbf{644} (2002), 3--20.

\bibitem{EO}
B.~Eynard and N.~Orantin,
\emph{Invariants of algebraic curves and topological
expansion},
Comm.\ Number Theory Phys.\ \textbf{1} (2007), 347--452.

\bibitem{FG12}
J.~Francis and D.~Gaitsgory,
Chiral Koszul duality,
\emph{Selecta Math.\ (N.S.)}~\textbf{18} (2012), 27--87.

\bibitem{FP00}
C.~Faber and R.~Pandharipande,
Hodge integrals and moduli spaces of curves,
in \emph{Moduli of Curves and Abelian Varieties},
Aspects Math.~E33, Vieweg, 2000, pp.~109--130.

\bibitem{FrenkelBenZvi}
E.~Frenkel and D.~Ben-Zvi,
\emph{Vertex Algebras and Algebraic Curves},
2nd ed., Mathematical Surveys and Monographs, vol.~88,
AMS, 2004.

\bibitem{GK}
E.~Getzler and M.~M.~Kapranov,
Modular operads,
\emph{Compositio Math.}~\textbf{110} (1998), 65--126.

\bibitem{JPPZ}
F.~Janda, R.~Pandharipande, A.~Pixton, and D.~Zvonkine,
Double ramification cycles on the moduli spaces of curves,
\emph{Publ.\ Math.\ Inst.\ Hautes \'Etudes Sci.}
\textbf{125} (2017), 221--266.

\bibitem{Kac-vertex}
V.~G.~Kac,
\emph{Vertex Algebras for Beginners},
2nd ed., University Lecture Series, vol.~10, AMS, 1998.

\bibitem{KontsevichSoibelman}
M.~Kontsevich and Y.~Soibelman,
Stability structures, motivic Donaldson--Thomas
invariants and cluster transformations,
arXiv:0811.2435.

\bibitem{LV12}
J.-L.~Loday and B.~Vallette,
\emph{Algebraic Operads},
Grundlehren der Mathematischen Wissenschaften, vol.~346,
Springer, 2012.

\bibitem{Priddy70}
S.~Priddy,
Koszul resolutions,
\emph{Trans.\ Amer.\ Math.\ Soc.}~\textbf{152} (1970),
39--60.

\bibitem{Zhu}
Y.~Zhu,
Modular invariance of characters of vertex operator
algebras,
\emph{J.~Amer.\ Math.\ Soc.}~\textbf{9} (1996), 237--302.

\bibitem{LorgatVirR}
R.~Lorgat,
\emph{The Virasoro spectral $R$-matrix: closed form and
non-formality},
preprint, 2026.

\bibitem{LorgatSevenFaces}
R.~Lorgat,
\emph{The seven faces of the collision residue},
preprint, 2026.

\end{thebibliography}

\end{document}
